\documentclass[11pt]{amsart}

\begin{filecontents*}[overwrite]{vacuum-bell.bib}
@article{Araki64,
  author = {Araki, H.},
  title = {Von {N}eumann algebras of local observables for free scalar field},
  journal = {Journal of Mathematical Physics},
  volume = {5},
  pages = {1--13},
  year = {1964},
  note = {\doi{10.1063/1.1704063}}
}
@article{Bargmann61,
  author = {Bargmann, V.},
  title = {On a {H}ilbert space of analytic functions and an associated integral transform. {P}art {I}},
  journal = {Communications on Pure and Applied Mathematics},
  volume = {14},
  pages = {187--214},
  year = {1961},
  note = {\doi{10.1002/cpa.3160140303}}
}
@article{Beigi13,
  author = {Beigi, S.},
  title = {Sandwiched {R}\'enyi divergence satisfies data processing inequality},
  journal = {Journal of Mathematical Physics},
  volume = {54},
  pages = {122202},
  year = {2013},
  note = {\doi{10.1063/1.4838855}; \arxiv{1306.5920v6}}
}
@book{Bhatia97,
  author = {Bhatia, R.},
  title = {Matrix Analysis},
  series = {Graduate Texts in Mathematics},
  volume = {169},
  publisher = {Springer},
  address = {New York},
  year = {1997},
  note = {\doi{10.1007/978-1-4612-0653-8}}
}
@article{BW75,
  author = {Bisognano, J. J. and Wichmann, E. H.},
  title = {On the duality condition for a {H}ermitian scalar field},
  journal = {Journal of Mathematical Physics},
  volume = {16},
  pages = {985--1007},
  year = {1975},
  note = {\doi{10.1063/1.522605}}
}
@book{BR87,
  author = {Bratteli, O. and Robinson, D. W.},
  title = {Operator Algebras and Quantum Statistical Mechanics 1:
           {$C^*$}- and {$W^*$}-Algebras, Symmetry Groups, Decomposition of States},
  edition = {Second},
  publisher = {Springer},
  address = {Berlin},
  year = {1987},
  note = {\doi{10.1007/978-3-662-02520-8}}
}
@article{BGL02,
  author = {Brunetti, R. and Guido, D. and Longo, R.},
  title = {Modular localization and {W}igner particles},
  journal = {Reviews in Mathematical Physics},
  volume = {14},
  pages = {759--785},
  year = {2002},
  note = {\doi{10.1142/S0129055X02001387}; \arxiv{math-ph/0203021v2}}
}
@article{Buchholz74,
  author = {Buchholz, D.},
  title = {Product states for local algebras},
  journal = {Communications in Mathematical Physics},
  volume = {36},
  pages = {287--304},
  year = {1974},
  note = {\doi{10.1007/BF01646201}}
}
@article{BDL07,
  author = {Buchholz, D. and D'Antoni, C. and Longo, R.},
  title = {Nuclearity and thermal states in conformal field theory},
  journal = {Communications in Mathematical Physics},
  volume = {270},
  pages = {267--293},
  year = {2007},
  note = {\doi{10.1007/s00220-006-0127-9}; \arxiv{math-ph/0603083v2}}
}
@article{BL04,
  author = {Buchholz, D. and Lechner, G.},
  title = {Modular nuclearity and localization},
  journal = {Annales Henri Poincar\'e},
  volume = {5},
  pages = {1065--1080},
  year = {2004},
  note = {\doi{10.1007/s00023-004-0190-8}; \arxiv{math-ph/0402072v1}}
}
@article{Tsirelson80,
  author = {Cirel'son, B. S.},
  title = {Quantum generalizations of {B}ell's inequality},
  journal = {Letters in Mathematical Physics},
  volume = {4},
  pages = {93--100},
  year = {1980},
  note = {\doi{10.1007/BF00417500}; the author's name is also transliterated as B.~S. Tsirelson}
}
@article{CHSH69,
  author = {Clauser, J. F. and Horne, M. A. and Shimony, A. and Holt, R. A.},
  title = {Proposed experiment to test local hidden-variable theories},
  journal = {Physical Review Letters},
  volume = {23},
  pages = {880--884},
  year = {1969},
  note = {\doi{10.1103/PhysRevLett.23.880}}
}
@article{EO73,
  author = {Eckmann, J.-P. and Osterwalder, K.},
  title = {An application of {T}omita's theory of modular {H}ilbert algebras:
           duality for free {B}ose fields},
  journal = {Journal of Functional Analysis},
  volume = {13},
  pages = {1--12},
  year = {1973},
  note = {\doi{10.1016/0022-1236(73)90062-1}}
}
@article{FG89,
  author = {Figliolini, F. and Guido, D.},
  title = {The {T}omita operator for the free scalar field},
  journal = {Annales de l'Institut Henri Poincar\'e, Physique th\'eorique},
  volume = {51},
  pages = {419--435},
  year = {1989},
  note = {\url{https://www.numdam.org/item/AIHPA_1989__51_4_419_0/}}
}
@article{FL13,
  author = {Frank, R. L. and Lieb, E. H.},
  title = {Monotonicity of a relative {R}\'enyi entropy},
  journal = {Journal of Mathematical Physics},
  volume = {54},
  pages = {122201},
  year = {2013},
  note = {\doi{10.1063/1.4838835}; \arxiv{1306.5358v3}}
}
@book{Haag96,
  author = {Haag, R.},
  title = {Local Quantum Physics: Fields, Particles, Algebras},
  edition = {Second},
  publisher = {Springer},
  address = {Berlin},
  year = {1996},
  note = {\doi{10.1007/978-3-642-61458-3}}
}
@article{HC00,
  author = {Halvorson, H. and Clifton, R.},
  title = {Generic {B}ell correlation between arbitrary local algebras in quantum field theory},
  journal = {Journal of Mathematical Physics},
  volume = {41},
  pages = {1711--1717},
  year = {2000},
  note = {\doi{10.1063/1.533253}; \arxiv{math-ph/9909013v3}}
}
@article{Hansen80,
  author = {Hansen, F.},
  title = {An operator inequality},
  journal = {Mathematische Annalen},
  volume = {246},
  pages = {249--250},
  year = {1980},
  note = {\doi{10.1007/BF01371046}}
}
@article{Heinz51,
  author = {Heinz, E.},
  title = {Beitr{\"a}ge zur {S}t{\"o}rungstheorie der {S}pektralzerlegung},
  journal = {Mathematische Annalen},
  volume = {123},
  pages = {415--438},
  year = {1951},
  note = {\doi{10.1007/BF02054965}}
}
@book{HS18,
  author = {Hollands, S. and Sanders, K.},
  title = {Entanglement Measures and Their Properties in Quantum Field Theory},
  series = {SpringerBriefs in Mathematical Physics},
  volume = {34},
  publisher = {Springer},
  address = {Cham},
  year = {2018},
  note = {\doi{10.1007/978-3-319-94902-4}; \arxiv{1702.04924v3}}
}
@book{Hormander03,
  author = {H{\"o}rmander, L.},
  title = {The Analysis of Linear Partial Differential Operators {I}:
           Distribution Theory and {F}ourier Analysis},
  edition = {Second},
  series = {Classics in Mathematics},
  publisher = {Springer},
  address = {Berlin},
  year = {2003},
  note = {\doi{10.1007/978-3-642-61497-2}}
}
@misc{OQP12,
  author = {{IQOQI Vienna}},
  title = {{O}pen {Q}uantum {P}roblems, {P}roblem 12:
           {B}ell inequalities for long range vacuum correlations},
  howpublished = {\url{https://oqp.iqoqi.oeaw.ac.at/bell-inequalities-for-long-range-vacuum-correlations}},
  note = {Archived copy of 29 October 2023,
          \url{https://web.archive.org/web/20231029013530/https://oqp.iqoqi.oeaw.ac.at/bell-inequalities-for-long-range-vacuum-correlations};
          accessed October 21, 2026}
}
@misc{KW05,
  author = {Kr{\"u}ger, O. and Werner, R. F.},
  title = {Some open problems in quantum information theory},
  howpublished = {\arxiv{quant-ph/0504166v1}},
  year = {2005},
  note = {Open problems in quantum information theory, Institut f{\"u}r
          Mathematische Physik, TU Braunschweig; version of 21 April 2005}
}
@article{KA80,
  author = {Kubo, F. and Ando, T.},
  title = {Means of positive linear operators},
  journal = {Mathematische Annalen},
  volume = {246},
  pages = {205--224},
  year = {1980},
  note = {\doi{10.1007/BF01371042}}
}
@article{Landau87,
  author = {Landau, L. J.},
  title = {On the violation of {B}ell's inequality in quantum theory},
  journal = {Physics Letters A},
  volume = {120},
  pages = {54--56},
  year = {1987},
  note = {\doi{10.1016/0375-9601(87)90075-2}}
}
@article{LS16,
  author = {Lechner, G. and Sanders, K.},
  title = {Modular nuclearity: a generally covariant perspective},
  journal = {Axioms},
  volume = {5},
  pages = {5},
  year = {2016},
  note = {\doi{10.3390/axioms5010005}; \arxiv{1511.09027v1}}
}
@article{MN20,
  author = {Matsumura, A. and Nambu, Y.},
  title = {Violation of {B}ell-{CHSH} inequalities through optimal local filters in the vacuum},
  journal = {Quantum Reports},
  volume = {2},
  pages = {542--559},
  year = {2020},
  note = {\doi{10.3390/quantum2040038}; \arxiv{1808.03896v2}}
}
@article{ML13,
  author = {M{\"u}ller-Lennert, M. and Dupuis, F. and Szehr, O. and
            Fehr, S. and Tomamichel, M.},
  title = {On quantum {R}\'enyi entropies: a new generalization and some properties},
  journal = {Journal of Mathematical Physics},
  volume = {54},
  pages = {122203},
  year = {2013},
  note = {\doi{10.1063/1.4838856}; \arxiv{1306.3142v3}}
}
@misc{DLMF,
  author = {{National Institute of Standards and Technology}},
  title = {{NIST} {D}igital {L}ibrary of {M}athematical {F}unctions},
  year = {2026},
  howpublished = {\url{https://dlmf.nist.gov/10.32}},
  note = {Release 1.2.8 of 2026-09-15, F.~W.~J. Olver, A.~B. Olde~Daalhuis,
          D.~W. Lozier, B.~I. Schneider, R.~F. Boisvert, C.~W. Clark,
          B.~R. Miller, B.~V. Saunders, H.~S. Cohl, and M.~A. McClain, eds.;
          Chapter 10, Section 10.32, Integral representations;
          accessed October 21, 2026}
}
@article{RS61,
  author = {Reeh, H. and Schlieder, S.},
  title = {Bemerkungen zur {Unit{\"a}r{\"a}quivalenz} von lorentzinvarianten {Feldern}},
  journal = {Il Nuovo Cimento},
  volume = {22},
  pages = {1051--1068},
  year = {1961},
  note = {\doi{10.1007/BF02787889}}
}
@article{RRS05,
  author = {Reznik, B. and Retzker, A. and Silman, J.},
  title = {Violating {B}ell's inequalities in vacuum},
  journal = {Physical Review A},
  volume = {71},
  pages = {042104},
  year = {2005},
  note = {\doi{10.1103/PhysRevA.71.042104}; \arxiv{quant-ph/0310058v2}}
}
@article{RvD77,
  author = {Rieffel, M. A. and {van Daele}, A.},
  title = {A bounded operator approach to {T}omita--{T}akesaki theory},
  journal = {Pacific Journal of Mathematics},
  volume = {69},
  pages = {187--221},
  year = {1977},
  note = {\doi{10.2140/pjm.1977.69.187}}
}
@article{SLW18,
  author = {Seshadreesan, K. P. and Lami, L. and Wilde, M. M.},
  title = {{R}\'enyi relative entropies of quantum {G}aussian states},
  journal = {Journal of Mathematical Physics},
  volume = {59},
  pages = {072204},
  year = {2018},
  note = {\doi{10.1063/1.5007167}; \arxiv{1706.09885v2}}
}
@book{Simon05,
  author = {Simon, B.},
  title = {Trace Ideals and Their Applications},
  edition = {Second},
  series = {Mathematical Surveys and Monographs},
  volume = {120},
  publisher = {American Mathematical Society},
  address = {Providence, RI},
  year = {2005},
  note = {\doi{10.1090/surv/120}}
}
@article{Summers90,
  author = {Summers, S. J.},
  title = {On the independence of local algebras in quantum field theory},
  journal = {Reviews in Mathematical Physics},
  volume = {2},
  pages = {201--247},
  year = {1990},
  note = {\doi{10.1142/S0129055X90000090}}
}
@incollection{Summers11,
  author = {Summers, S. J.},
  title = {Yet more ado about nothing: the remarkable relativistic vacuum state},
  booktitle = {Deep Beauty: Understanding the Quantum World through Mathematical Innovation},
  editor = {Halvorson, H.},
  publisher = {Cambridge University Press},
  address = {Cambridge},
  year = {2011},
  note = {pp.~317--342, \doi{10.1017/CBO9780511976971.009}; \arxiv{0802.1854v1}}
}
@article{SW85,
  author = {Summers, S. J. and Werner, R. F.},
  title = {The vacuum violates {B}ell's inequalities},
  journal = {Physics Letters A},
  volume = {110},
  pages = {257--259},
  year = {1985},
  note = {\doi{10.1016/0375-9601(85)90093-3}}
}
@article{SW87a,
  author = {Summers, S. J. and Werner, R. F.},
  title = {{B}ell's inequalities and quantum field theory.
           {I}. {G}eneral setting},
  journal = {Journal of Mathematical Physics},
  volume = {28},
  pages = {2440--2447},
  year = {1987},
  note = {\doi{10.1063/1.527733}}
}
@article{SW87b,
  author = {Summers, S. J. and Werner, R. F.},
  title = {{B}ell's inequalities and quantum field theory.
           {II}. {B}ell's inequalities are maximally violated in the vacuum},
  journal = {Journal of Mathematical Physics},
  volume = {28},
  pages = {2448--2456},
  year = {1987},
  note = {\doi{10.1063/1.527734}}
}
@article{SW87c,
  author = {Summers, S. J. and Werner, R. F.},
  title = {Maximal violation of {B}ell's inequalities is generic in quantum field theory},
  journal = {Communications in Mathematical Physics},
  volume = {110},
  pages = {247--259},
  year = {1987},
  note = {\doi{10.1007/BF01207366}}
}
@article{SW88,
  author = {Summers, S. J. and Werner, R. F.},
  title = {Maximal violation of {B}ell's inequalities for algebras of observables
           in tangent spacetime regions},
  journal = {Annales de l'Institut Henri Poincar\'e, Physique th\'eorique},
  volume = {49},
  pages = {215--243},
  year = {1988},
  note = {\url{https://www.numdam.org/item/AIHPA_1988__49_2_215_0/}}
}
@article{VW05,
  author = {Verch, R. and Werner, R. F.},
  title = {Distillability and positivity of partial transposes in general quantum field systems},
  journal = {Reviews in Mathematical Physics},
  volume = {17},
  pages = {545--576},
  year = {2005},
  note = {\doi{10.1142/S0129055X05002364}; \arxiv{quant-ph/0403089v3}}
}
@article{WWY14,
  author = {Wilde, M. M. and Winter, A. and Yang, D.},
  title = {Strong converse for the classical capacity of entanglement-breaking
           and {H}adamard channels via a sandwiched {R}\'enyi relative entropy},
  journal = {Communications in Mathematical Physics},
  volume = {331},
  pages = {593--622},
  year = {2014},
  note = {\doi{10.1007/s00220-014-2122-x}; \arxiv{1306.1586v4}}
}
@article{Williamson36,
  author = {Williamson, J.},
  title = {On the algebraic problem concerning the normal forms of linear dynamical systems},
  journal = {American Journal of Mathematics},
  volume = {58},
  pages = {141--163},
  year = {1936},
  note = {\doi{10.2307/2371062}}
}
@article{Witten18,
  author = {Witten, E.},
  title = {{APS} {M}edal for {E}xceptional {A}chievement in {R}esearch:
           {I}nvited article on entanglement properties of quantum field theory},
  journal = {Reviews of Modern Physics},
  volume = {90},
  pages = {045003},
  year = {2018},
  note = {\doi{10.1103/RevModPhys.90.045003}; \arxiv{1803.04993v6}, titled
          ``Notes on some entanglement properties of quantum field theory''}
}
@article{Zych10,
  author = {Zych, M. and Costa, F. and Kofler, J. and Brukner, {\v C}.},
  title = {Entanglement between smeared field operators in the {K}lein--{G}ordon vacuum},
  journal = {Physical Review D},
  volume = {81},
  pages = {125019},
  year = {2010},
  note = {\doi{10.1103/PhysRevD.81.125019}; \arxiv{1003.3354v2}}
}
\end{filecontents*}

\usepackage[letterpaper,margin=1in]{geometry}
\usepackage[T1]{fontenc}
\usepackage{lmodern}
\usepackage{amsmath,amssymb,amsthm,mathtools,mathrsfs}
\usepackage{microtype}
\usepackage{booktabs,array,tabularx,longtable}
\usepackage{enumitem}
\usepackage[most]{tcolorbox}
\usepackage[hidelinks]{hyperref}
\usepackage{etoolbox}
\usepackage{xcolor}
\hypersetup{pdftitle={The vacuum Bell correlation of the massive free scalar field
  attains the classical bound for bounded regions beyond a finite separation}}

\allowdisplaybreaks[2]
\numberwithin{equation}{section}
\setlength{\emergencystretch}{1.5em}
\setlist[enumerate]{itemsep=5pt,topsep=6pt}

\newtheorem{theorem}{Theorem}[section]
\newtheorem{lemma}[theorem]{Lemma}
\newtheorem{corollary}[theorem]{Corollary}
\newtheorem*{mainrestatement}{Theorem 1.1}
\theoremstyle{definition}
\newtheorem{convention}[theorem]{Convention}

\newcommand{\A}{\mathfrak A}
\newcommand{\B}{\mathfrak B}
\newcommand{\R}{\mathbb R}
\newcommand{\C}{\mathbb C}
\newcommand{\N}{\mathbb N}
\newcommand{\Hh}{\mathcal H}
\newcommand{\Sc}{\mathcal S}
\newcommand{\vv}{\mathsf v}
\newcommand{\Res}{\mathsf R}
\newcommand{\JointRes}{\mathcal R}
\newcommand{\JointD}{\mathsf D}
\newcommand{\Therm}{\mathsf H}
\newcommand{\Qtil}{\widetilde Q}
\newcommand{\sx}{\mathsf x}
\newcommand{\sy}{\mathsf y}
\newcommand{\bx}{\mathbf x}
\newcommand{\by}{\mathbf y}
\newcommand{\bp}{\mathbf p}
\newcommand{\bz}{\mathbf z}
\newcommand{\eone}{\mathbf e_1}
\newcommand{\Energy}{\mathcal E}
\newcommand{\Tr}{\operatorname{Tr}}
\newcommand{\Ran}{\operatorname{Ran}}
\newcommand{\Dom}{\operatorname{Dom}}
\newcommand{\supp}{\operatorname{supp}}
\newcommand{\rank}{\operatorname{rank}}
\newcommand{\diag}{\operatorname{diag}}
\newcommand{\arctanh}{\operatorname{arctanh}}
\newcommand{\vol}{\operatorname{vol}}
\newcommand{\Rea}{\operatorname{Re}}
\newcommand{\Ima}{\operatorname{Im}}
\newcommand{\abs}[1]{\left|#1\right|}
\newcommand{\norm}[1]{\left\|#1\right\|}
\newcommand{\ip}[2]{\langle #1,#2\rangle}
\newcommand{\wh}[1]{\widehat{#1}}
\newcommand{\wt}[1]{\widetilde{#1}}
\newcommand{\doi}[1]{\href{https://doi.org/#1}{\nolinkurl{https://doi.org/#1}}}
\newcommand{\arxiv}[1]{\href{https://arxiv.org/abs/#1}{arXiv:\nolinkurl{#1}}}
% Reference list: ragged right, and no URL break after the colon of "https:".
\makeatletter
\apptocmd{\thebibliography}{\raggedright
  \def\UrlBigBreaks{\do@url@hyp}%
  \expandafter\def\expandafter\UrlOrds\expandafter{\UrlOrds\do\:}}{}{}
\makeatother

\title[Vacuum Bell correlation beyond a finite separation]
{The vacuum Bell correlation of the massive free scalar field
attains the classical bound for bounded regions beyond a finite separation}
\author{David Turturean}\thanks{D.T. built and ran the automated research and formalization workflows at PSI that produced these results, steered the writing of the manuscript, and asked for AI reviews and revisions of the drafts.}
\address{Physical Superintelligence PBC, Cambridge, Massachusetts, USA}
\date{}
\subjclass[2020]{81T05, 81T10, 81P40, 46L60}
\keywords{Vacuum correlations, Bell inequalities, free scalar field,
local algebras, standard subspaces, Gaussian states}

\begin{document}

\begin{abstract}
Summers and Werner showed that the vacuum of a massive relativistic quantum field violates the Clauser--Horne--Shimony--Holt (CHSH) inequality maximally for suitable spacelike separated regions, and that the violation decays as the regions are moved apart. Problem~12 of the IQOQI Vienna Open Quantum Problems list asks whether some violation is possible for regions arbitrarily far apart. We answer it in the negative for fixed bounded regions. For the real massive free scalar field in $3+1$ dimensions, the vacuum Bell correlation of two double cones of radius $R$ centred at equal times, the supremum of the CHSH values of their local observables, equals the classical bound once the spatial distance between their centres exceeds a fixed multiple of $R$, independent of the mass. Hence the vacuum satisfies the CHSH inequality for any two bounded regions moved sufficiently far apart in a fixed spacelike direction, although it remains entangled across them.
\end{abstract}

\maketitle

\clearpage
\vspace*{\fill}
\section*{Acknowledgements}

We thank Adam Levine for reading the manuscript, for his comments on the presentation and for condensing the manuscript into the letter version of the paper. We also thank Physical Superintelligence PBC, in which the author has a financial interest, for use of its resources, including its technology and frontier AI models. The author used these models to develop, derive, and compose this manuscript and its Lean~4 supplements.

Results were formalized in Lean~4 by AI agents and checked by Lean's kernel; the formalization is at \url{https://github.com/physical-superintelligence/oqp-batch-2026/tree/main/oqp12-vacuum-bell-correlations/lean}.

David Turturean reviewed and edited this paper and is responsible for its correctness. David Turturean is employed by PSI and has a financial interest in it.

\vspace*{\fill}
\newpage

{\small\tableofcontents}

\section{Introduction}

Vacuum entanglement is a structural feature of relativistic quantum field
theory. The Reeh--Schlieder theorem makes the vacuum cyclic for local
algebras \cite{RS61}; together with noncommutativity, this cyclicity implies
entanglement across spacelike separated regions \cite{HC00}.
The operator-algebraic and quantitative aspects of vacuum entanglement
are developed in the account of Hollands and Sanders \cite{HS18}, and
the operator-algebraic ideas are explained informally in Witten's notes
\cite{Witten18}. Bell
correlation is a stronger property: some local measurements in the vacuum
violate a Bell inequality. At large spacelike separation the vacuum stays
entangled \cite{HC00,VW05}, while its Bell correlation decays towards the
classical value, as the mass-gap bound below shows.

Summers and Werner had established maximal vacuum violation of Bell's
inequalities for suitable localizations in free Bose and Fermi theories
\cite{SW85,SW87b}. Their further results had identified complementary wedges
and tangent double cones as geometries with maximal Bell correlation;
under the hypotheses of those results, maximality holds in every normal
state \cite{SW87c,SW88}.

At positive separation, a mass gap instead gives
the upper bound on the Bell correlation \eqref{eq:beta}
\[
 \beta\bigl(\omega;\A(O_1),\A(O_2)\bigr)
 \le 1+2e^{-m d(O_1,O_2)}.
\]
Here \(\omega\) is the vacuum state of a theory with lowest mass \(m>0\),
the normalization has classical bound \(1\), and
\(d(O_1,O_2)\) is the maximal timelike distance through which one region
can be translated while preserving spacelike separation
\cite{SW87a}, \cite[Proposition~4.5]{Summers11}. That is, \(d(O_1,O_2)\) is the largest
\(\tau\) such that \(O_1+(t,0)\) and \(O_2\) are spacelike separated for
all \(|t|<\tau\); for the two double cones considered below, it equals
\(L-2R\).

The relation between Bell correlation and statistical independence had been
systematized by Summers \cite{Summers90}. His 2008 survey, subsequently
published as \cite{Summers11}, had noted that model-independent lower
bounds on the distant vacuum Bell correlation were not yet available. Whether two distant regions show Bell correlation depends on the state: Landau's
result gives states with maximal violation for commuting nonabelian
algebras satisfying the Schlieder property \cite{Landau87}.
Halvorson and Clifton had proved generic Bell correlation of states for
spacelike local algebras \cite{HC00}; they had noted that none of their results ``decide the question of
whether the vacuum state always retains some Bell correlation across
arbitrary spacelike separated regions'' \cite[Section~IV]{HC00}. Verch and Werner had further proved vacuum
distillability across arbitrarily large spacelike distances
\cite{VW05}.

Reznik, Retzker and Silman
\cite{RRS05}, and later Matsumura and Nambu \cite{MN20}, had obtained
Bell--CHSH violation after filtering locally coupled probes, in
\cite{RRS05} at every separation of the probes. Their
unfiltered detector states satisfy CHSH in the regimes analysed.
Filtering changes the experiment by selecting an output state, whereas
the regional Bell correlation considered here retains every outcome of
the original vacuum measurement. Verch and Werner had remarked that
the authors of \cite{RRS05} do not demonstrate Bell correlations of the
vacuum over arbitrarily separated and arbitrarily small regions ``in the
sense of proving a violation of the CHSH inequalities by the quantum field
observables themselves'' \cite[Section~7]{VW05}.

The question whether the vacuum can violate the CHSH inequality for
regions arbitrarily far apart, with every outcome retained, is Problem~12
of the IQOQI Vienna Open
Quantum Problems list \cite{OQP12} and of the collection of open
problems in quantum information of Krüger and Werner \cite{KW05}, where
the problem is dated 22~January 2002 and has R.~Verch as contact; its statement
is reproduced verbatim.\footnotemark[1]

\begin{tcolorbox}[enhanced,breakable,sharp corners,boxrule=0.5pt,
colframe=black!60,colback=gray!4,
title={Problem 12: Bell Inequalities for long range vacuum correlations (IQOQI Vienna Open Quantum Problems \cite{OQP12})},
fonttitle=\bfseries]
It is well known [1] that vacuum fluctuations maximally violate the CHSH-Bell inequalities for suitable spacelike separated observables, and that this violation goes to zero as the two localization regions are moved apart.

Decide whether some (necessarily small) violation of the inequalities is possible for regions arbitrarily far apart. For definiteness, consider a massive scalar free relativistic Bose field.
\end{tcolorbox}
\footnotetext[1]{The box quotes the IQOQI list, which gives the same
date and contact. The reference ``[1]'' in the quoted problem statement is
the paper of Summers and Werner \cite{SW85}.}

In this paper we prove that the answer is negative for two double cones
of equal radius moved apart (Theorem~\ref{thm:main} below), and hence for
any two bounded regions moved apart in a fixed spacelike direction
(Corollary~\ref{cor:fixed}). Problem~12 remains open for regions that
grow in proportion to their separation, at ratios that
Theorem~\ref{thm:main} does not cover, and for wedges at a positive gap, as does the
relation of these results to Bell tests after local filtering
(Section~\ref{sec:discussion}).

For \(R>0\), the double cone
\[
 D_R=\{(x^0,\bx)\in\R^{1+3}: |x^0|+|\bx|<R\}
\]
has base \(B_R(0)\) at time zero. We consider the real free scalar field
of mass \(m>0\), its vacuum state, and the two local von Neumann algebras (defined in
Section~\ref{sec:setting})
\[
 \A_m(D_R),\qquad \A_m\bigl(D_R+(0,L\eone)\bigr),\qquad L>2R.
\]
Thus \(L\) is the spatial distance between the centres of the two
double cones, and their bases are separated by the gap \(L-2R\).

The Summers--Werner Bell correlation of the two algebras in the vacuum
state is denoted by
\(\beta_m(R,L)\), with classical bound \(1\) and Tsirelson bound
\(\sqrt2\); its algebraic definition is \eqref{eq:beta}. The supremum defining
\(\beta_m(R,L)\) includes all binary local measurements, represented by the
self-adjoint contractions in \eqref{eq:beta}. The measurements may depend on \(L\), and their outcome probabilities may be
arbitrarily small.

\begin{theorem}[Separation threshold]\label{thm:main}
There is a finite constant \(C_*>2\), independent of \(m\) and \(R\),
such that
\begin{equation}
 \forall m,R>0,\qquad
 L\ge C_*R\quad\Longrightarrow\quad \beta_m(R,L)=1.
 \label{eq:main}
\end{equation}
\end{theorem}

The constant produced by the argument is astronomically large: for the
parameters in Table~\ref{tab:constants} of Section~\ref{sec:completion}, \(C_*\ge C_{\mathrm{prob}}=Q_0\cdot10^{362\,232.284166\ldots}\), where the
numerical enclosure of the correlation coefficient \(Q_0\) of
\eqref{eq:mass-uniform-cluster} remains open. Theorem~\ref{thm:main} is therefore a pure
existence statement about eventual behaviour, with no information at
physically relevant separations (Section~\ref{sec:discussion},
item~(4)).

\begin{corollary}[Fixed bounded regions]\label{cor:fixed}
For every \(m>0\), any two fixed bounded localization regions
\(O_1,O_2\), and any fixed spacelike translation direction, the vacuum
Bell correlation of the local algebras of \(O_1\) and of the translate
of \(O_2\) equals \(1\) for all sufficiently large translations of
\(O_2\) in that direction.
\end{corollary}

Corollary~\ref{cor:fixed} follows from Theorem~\ref{thm:main} by
Poincar\'e covariance and isotony (Section~\ref{subsec:main-proof}).

\subsection{The mechanism of the threshold}

The central estimates compare a joint event probability with the
product of its marginals. For projections \(E\in\A_m(D_R)\) and \(P\in\A_m(D_R+(0,L\eone))\), with
subscripts \(A\) and \(B\) for the two double cones, we write
\[
 e_A=\omega(E),\qquad e_B=\omega(P),\qquad w=\omega(EP).
\]
With the constants \(k,\lambda\) of \eqref{eq:event-orders}, and with the
exponents \(a,b\) and \(\eta=1-(1-b)/(ab^2)\) of
\eqref{eq:probability-exponents}, the two event bounds have the form
\[
 w\le k(e_Ae_B)^b,\qquad
 w\ge\lambda(e_Ae_B)^{1/a},\qquad
 0<a<\tfrac12<b<1,
 \qquad
 1-\frac{1-b}{ab^2}>0.
\]
The constants \(k,\lambda\) approach \(1\) as \(L/R\) increases,
uniformly in \(m\) (see \eqref{eq:full-K}). By \eqref{eq:envelopes}, the
bounds on \(k-1\) and \(1-\lambda\) decay like powers of \(R/L\), more slowly
than the mass-gap bound above, because the uniformity extends to small
\(mR\).

In a violating CHSH experiment, the two event
bounds force a positive lower bound on the least likely winning cell, a joint outcome probability defined in
Step~2 of the proof of Theorem~\ref{thm:probability}.
The correlation estimate \eqref{eq:additive-cluster}, supplied by
\eqref{eq:mass-uniform-cluster}, forces that same cell to have probability
tending to zero as \(L/R\) grows, and beyond a finite separation the two
requirements are incompatible (Theorem~\ref{thm:probability}).
The upper event bound and the correlation estimate are quantitative
forms of clustering; the lower event bound, which controls the joint
probability of arbitrarily rare events through the joint covariance
norm (Sections~\ref{sec:evaluation} and~\ref{sec:reverse}), is the
ingredient specific to this argument, and it excludes a
violation carried by events of very small probability.

The two event bounds arise from opposite orders of comparison between a
finite Gaussian vacuum restriction \(\rho_F\) and the product \(\sigma_F\) of
its marginals, through the sandwiched Rényi quantities
\(\Qtil_{r_+}(\rho_F\|\sigma_F)\) in \eqref{eq:Qforward} and
\(\Qtil_{r_-}(\sigma_F\|\rho_F)\) in \eqref{eq:Qreverse}.
The forward comparison uses the uncertainty reserves \eqref{eq:reserve}
of the two balls (Section~\ref{sec:forward}). The reverse comparison uses
the uncertainty reserve of their union in the joint covariance norm
(Section~\ref{sec:reverse}). Estimates on the complete source spaces precede
all finite-dimensional restrictions, and compression preserves the
bounds. Hence the constants are independent of the
selected modes. Buchholz's construction of normal product states for
separated local algebras \cite{Buchholz74} supplies the broader
split-property context for the comparison with the product of the
marginals, which the argument carries out on finite source systems
(Section~\ref{sec:completion}).

\section{Setting and notation}\label{sec:setting}

\subsection{The field, local algebras, the Bell correlation, and scaling}

We use units \(\hbar=c=1\) and metric signature \((+,-,-,-)\).
Thus \(x\cdot y=x^0y^0-\bx\cdot\by\). Momentum components are written
\(p_0,\bp=(p_1,p_2,p_3)\) without index lowering. Spatial Fourier transforms are
\[
 \wh f(\bp)=\int_{\R^3}e^{-i\bp\cdot\bx}f(\bx)\,d^3x,\qquad
 f(\bx)=\frac1{(2\pi)^3}\int_{\R^3}e^{i\bp\cdot\bx}\wh f(\bp)\,d^3p.
\]
The unitary Fourier transform is
\(\mathcal Ff=(2\pi)^{-3/2}\wh f\). Euclidean norms on complex
spacetime are written \(|\cdot|_{\mathrm E}\).

For mass \(\mu>0\), the positive mass shell and one-particle space are
\[
 H_\mu^+=\{p\in\R^{1+3}:p\cdot p=\mu^2,\ p_0>0\},\qquad
 \Hh_{1,\mu}=L^2(H_\mu^+,d\mathfrak m_\mu),\qquad
 d\mathfrak m_\mu(p)=\frac{d^3\bp}{2(2\pi)^3p_0}.
\]
The vacuum representation acts on the symmetric Fock space
\(\mathscr F_\mu=\Gamma_{\mathrm s}(\Hh_{1,\mu})\), with vacuum vector
\(\Omega_\mu\). We abbreviate its vector state by
\(\omega(X)=\ip{\Omega_\mu}{X\Omega_\mu}\) when the mass is fixed.

The time-zero field and its conjugate momentum, smeared with real
\(f,g\in C_c^\infty(\R^3)\), are denoted by \(\phi_\mu(f)\) and
\(\pi_\mu(g)\). For a source pair \(\sx=(f,g)\), we set
\[
 h_\mu=(-\nabla^2+\mu^2)^{1/2},\qquad
 \Phi_\mu(\sx)=\phi_\mu(f)+\pi_\mu(g),\qquad
 \mathsf W_\mu(\sx)=e^{i\Phi_\mu(\sx)}.
\]
The symplectic form and vacuum covariance are
\begin{align}
 \varsigma(\sx,\sy)
 &=\int_{\R^3}(fg'-gf')\,d^3x,
 \qquad \sy=(f',g'),                                      \label{eq:symplectic}\\
 \vv_{\mu,U}(\sx,\sy)
 &=\frac12\left(\ip{f}{h_\mu^{-1}f'}+\ip{g}{h_\mu g'}\right)
 \quad\text{for sources supported in }U.                  \label{eq:covariance}
\end{align}
The subscript \(U\) in \eqref{eq:covariance} labels the source Hilbert space
introduced in Section~\ref{subsec:sources}; the right-hand side is
defined for any pair of sources and is independent of \(U\).
The free-field Weyl algebra satisfies
\[
 \mathsf W_\mu(\sx)\mathsf W_\mu(\sy)
 =e^{-i\varsigma(\sx,\sy)/2}\mathsf W_\mu(\sx+\sy),\qquad
 \mathsf W_\mu(\sx)^*=\mathsf W_\mu(-\sx).
\]
The vacuum is the centred quasi-free state characterized by
\[
 \omega(\mathsf W_\mu(\sx))
 =\exp\!\left[-\frac12\vv_{\mu,U}(\sx,\sx)\right].
\]
These formulas are the usual free-field and Weyl conventions
\cite{Haag96,SLW18}; the covariances of \cite{SLW18} are twice ours.

For a real \(F\in C_c^\infty(\R^{1+3})\), the spacetime-smeared field
is \(\phi_\mu(F)=\int\phi_\mu(x)F(x)\,d^4x\), with spacetime Weyl
operator \(\mathsf W_\mu(F)=e^{i\phi_\mu(F)}\). With
\(\partial_0=\partial/\partial x^0\), let \(E_{\mathrm{ret}}F\) and
\(E_{\mathrm{adv}}F\) be the solutions \(v\) of
\((\partial_0^2-\nabla^2+\mu^2)v=F\) that vanish at all sufficiently
early and at all sufficiently late times, respectively. Their difference
\(v_F=E_{\mathrm{ret}}F-E_{\mathrm{adv}}F\) solves the homogeneous
equation, and the field equation expresses \(\phi_\mu(F)\) through the
Cauchy data of \(v_F\) at time zero:
\[
 \phi_\mu(F)=\Phi_\mu(\sx),\qquad
 \sx=\bigl(\partial_0v_F(0,\cdot),\,-v_F(0,\cdot)\bigr).
\]
By finite propagation speed, \(v_F\) vanishes outside the causal future
and past of \(\supp F\), so the time-zero pair \(\sx\) is supported in
the base \(B_R(\bx_0)\) of any double cone \(D_R+(0,\bx_0)\) containing
\(\supp F\).

For a spacetime region \(O\), the notation \(\A_\mu(O)\) denotes the von
Neumann algebra generated by the Weyl operators of real spacetime test
functions supported in \(O\). Araki's time-zero identification
\cite{Araki64}, together with the free Bose-field modular description
\cite{EO73}, gives
\begin{equation}
 \A_\mu(D_R)
 =\{\mathsf W_\mu(f,g):f,g\in C_c^\infty(B_R(0);\R)\}''.
 \label{eq:time-slice}
\end{equation}
Finite propagation gives the inclusion from spacetime smearings to
time-zero sources: if \(\supp F\subset D_R\), the formula for
\(\phi_\mu(F)\) above gives \(\mathsf W_\mu(F)=\mathsf W_\mu(\sx)\) with
\(\sx\) supported in \(B_R(0)\), so \(\A_\mu(D_R)\) lies in the
right-hand side of \eqref{eq:time-slice}.

For the converse inclusion, compact Cauchy data strictly inside \(B_R(0)\) determine a
homogeneous Klein--Gordon solution. Let \(f,g\in C_c^\infty(B_R(0);\R)\),
choose \(0<\delta<R\) with
\(\supp f\cup\supp g\subset\overline{B_{R-\delta}(0)}\), and let \(u\)
be the solution of \((\partial_0^2-\nabla^2+\mu^2)u=0\) with
\(u(0,\cdot)=-g\) and \(\partial_0u(0,\cdot)=f\). By finite propagation
speed, \(u(x^0,\cdot)\) is supported in
\(\overline{B_{R-\delta+|x^0|}(0)}\).

Multiplying \(u\) by a temporal cutoff that changes from \(0\) to \(1\) in
a sufficiently short interval around zero, and then applying the
Klein--Gordon operator, produces a test function supported in \(D_R\). Let
\(\chi\in C^\infty(\R)\) vanish on \((-\infty,-\epsilon]\) and equal
\(1\) on \([\epsilon,\infty)\), where \(0<2\epsilon<\delta\). With
\(\chi=\chi(x^0)\),
\[
 F=(\partial_0^2-\nabla^2+\mu^2)(\chi u)=\chi''u+2\chi'\,\partial_0u
\]
is real and smooth, and vanishes unless \(|x^0|\le\epsilon\) and
\(|\bx|\le R-\delta+\epsilon\). On that set
\(|x^0|+|\bx|\le R-\delta+2\epsilon<R\), so \(\supp F\subset D_R\).

The causal solution \(v_F\) of \(F\) recovers the original data. The functions
\(\chi u\) and \((\chi-1)u\) solve \((\partial_0^2-\nabla^2+\mu^2)v=F\)
and vanish for \(x^0\le-\epsilon\) and for \(x^0\ge\epsilon\),
respectively, so \(E_{\mathrm{ret}}F=\chi u\),
\(E_{\mathrm{adv}}F=(\chi-1)u\) and \(v_F=u\). The formula for
\(\phi_\mu(F)\) above gives
\[
 \phi_\mu(F)
 =\Phi_\mu\bigl(\partial_0u(0,\cdot),\,-u(0,\cdot)\bigr)
 =\Phi_\mu(f,g),
\]
so \(\mathsf W_\mu(f,g)=\mathsf W_\mu(F)\in\A_\mu(D_R)\). Since \(f\) and \(g\) were
arbitrary, the converse inclusion in \eqref{eq:time-slice} follows.

An \emph{effect} is an operator \(E\) with \(0\le E\le I\).
A self-adjoint contraction \(A\) represents the binary measurement with
effects \((I\pm A)/2\). For commuting von Neumann algebras \(\A,\B\), we
use the Summers--Werner normalization
\begin{equation}
 \beta(\omega;\A,\B)
 =
 \sup_{\substack{A_i=A_i^*\in\A,\ \norm{A_i}\le1\\
                 B_j=B_j^*\in\B,\ \norm{B_j}\le1}}
 \frac12\abs{\omega\bigl(A_0(B_0+B_1)+A_1(B_0-B_1)\bigr)}.
 \label{eq:beta}
\end{equation}

The classical CHSH bound is \(1\) \cite{CHSH69,SW87a}, and Tsirelson's
bound is \(\sqrt2\) \cite{Tsirelson80}. The identities \(A_0=A_1=B_0=B_1=I\) give the value \(1\), so
\(\beta(\omega;\A,\B)\ge1\).
Some authors, including Verch and Werner \cite{VW05}, omit the factor
\(1/2\), so their classical and Tsirelson bounds are \(2\) and
\(2\sqrt2\), respectively. In the normalization \eqref{eq:beta},
\[
 \beta_m(R,L)=
 \beta\bigl(\omega;\A_m(D_R),\A_m(D_R+(0,L\eone))\bigr).
\]

\begin{lemma}[Dilation]\label{lem:scaling}
The source transformation
\begin{equation}
 \mu=mR,\qquad s=L/R,\qquad
 \wt f(\by)=R^2f(R\by),\qquad
 \wt g(\by)=Rg(R\by)
 \label{eq:scaling}
\end{equation}
preserves the symplectic form and vacuum covariance, with mass \(m\)
replaced by \(\mu\). The transformation induces a vacuum-preserving isomorphism of the
corresponding local Weyl algebras and their von Neumann closures.
Consequently,
\[
 \beta_m(R,L)=\beta_{mR}(1,L/R).
\]
\end{lemma}

\begin{proof}
\emph{Step 1: the symplectic form.} Let \(\sx=(f,g)\) and
\(\sy=(f',g')\) be real smooth sources supported in a bounded open set
\(U\), and let \(\wt\sx=(\wt f,\wt g)\) and \(\wt\sy=(\wt{f'},\wt{g'})\)
be their transforms by \eqref{eq:scaling}, supported in
\(\wt U=\{\by:R\by\in U\}\). A change of variables in
\eqref{eq:symplectic} preserves the symplectic form: with \(\bx=R\by\),
\[
 \int_{\R^3}\wt f\,\wt{g'}\,d^3y
 =R^3\int_{\R^3}f(R\by)\,g'(R\by)\,d^3y
 =\int_{\R^3}fg'\,d^3x,
\]
and likewise for \(\wt g\,\wt{f'}\), so that
\(\varsigma(\wt\sx,\wt\sy)=\varsigma(\sx,\sy)\).

\smallskip
\emph{Step 2: the covariance.} The same change of variables in the
Fourier integral gives
\[
 \wh{\wt f}(\mathbf q)=R^{-1}\wh f(\mathbf q/R),\qquad
 \wh{\wt g}(\mathbf q)=R^{-2}\wh g(\mathbf q/R).
\]
The Fourier multiplier \(h_m\) in \eqref{eq:covariance} satisfies
\(h_{mR}(\mathbf q)=R\,h_m(\mathbf q/R)\). With
Plancherel's formula and the substitution \(\mathbf q=R\bp\), the factors
\(R^2\) and \(R\) in \eqref{eq:scaling} therefore preserve the two
covariance terms:
\begin{align*}
 \ip{\wt f}{h_{mR}^{-1}\wt{f'}}
 &=\frac1{(2\pi)^3}\int_{\R^3}
   \frac{\overline{\wh{\wt f}(\mathbf q)}\,\wh{\wt{f'}}(\mathbf q)}
        {h_{mR}(\mathbf q)}\,d^3q\\
 &=R^{-2}R^{-1}R^3\,\frac1{(2\pi)^3}\int_{\R^3}
   \frac{\overline{\wh f(\bp)}\,\wh{f'}(\bp)}{h_m(\bp)}\,d^3p
  =\ip{f}{h_m^{-1}f'},\\
 \ip{\wt g}{h_{mR}\wt{g'}}
 &=\frac1{(2\pi)^3}\int_{\R^3}
   \overline{\wh{\wt g}(\mathbf q)}\,\wh{\wt{g'}}(\mathbf q)\,
   h_{mR}(\mathbf q)\,d^3q\\
 &=R^{-4}R\,R^3\,\frac1{(2\pi)^3}\int_{\R^3}
   \overline{\wh g(\bp)}\,\wh{g'}(\bp)\,h_m(\bp)\,d^3p
  =\ip{g}{h_mg'}.
\end{align*}
Hence \(\vv_{mR,\wt U}(\wt\sx,\wt\sy)=\vv_{m,U}(\sx,\sy)\).

\smallskip
\emph{Step 3: the unitary.} The Weyl relations and the quasi-free
characteristic function are therefore preserved: \(\sx\mapsto\wt\sx\) is
a real linear bijection of the smooth sources, with inverse
\(f(\bx)=R^{-2}\wt f(\bx/R)\), \(g(\bx)=R^{-1}\wt g(\bx/R)\), and
\[
 \ip{\Omega_{mR}}{\mathsf W_{mR}(\wt\sx)\Omega_{mR}}
 =\exp\!\Bigl[-\frac12\vv_{mR,\wt U}(\wt\sx,\wt\sx)\Bigr]
 =\exp\!\Bigl[-\frac12\vv_{m,U}(\sx,\sx)\Bigr]
 =\ip{\Omega_m}{\mathsf W_m(\sx)\Omega_m}.
\]
In both representations the vacuum vector is cyclic for the Weyl
operators of the smooth sources. Uniqueness of the GNS representation
therefore gives a unitary \(\mathsf V\colon\mathscr F_m\to\mathscr F_{mR}\)
with
\[
 \mathsf V\Omega_m=\Omega_{mR},\qquad
 \mathsf V\mathsf W_m(\sx)\mathsf V^*=\mathsf W_{mR}(\wt\sx)
\]
for every real smooth source \(\sx\).

\smallskip
\emph{Step 4: the two regions.} The transformation maps the sources
supported in \(B_R(0)\) onto those supported in \(B_1(0)\), and those
supported in \(B_R(L\eone)\) onto those supported in \(B_1(s\eone)\).
By \eqref{eq:time-slice}, and by the same identity for the translated
double cones with bases \(B_R(L\eone)\) and \(B_1(s\eone)\), conjugation
by \(\mathsf V\) therefore maps \(\A_m(D_R)\) onto \(\A_{mR}(D_1)\) and
\(\A_m(D_R+(0,L\eone))\) onto \(\A_{mR}(D_1+(0,s\eone))\). This conjugation preserves
vacuum expectations,
\(\ip{\Omega_{mR}}{\mathsf VX\mathsf V^*\Omega_{mR}}=\ip{\Omega_m}{X\Omega_m}\).
Formula \eqref{eq:beta} then gives the stated identity:
\(A\mapsto\mathsf VA\mathsf V^*\) is a
bijection between the self-adjoint contractions of corresponding
algebras and leaves the expression under the supremum in
\eqref{eq:beta} unchanged, so \(\beta_m(R,L)=\beta_{mR}(1,L/R)\).
\end{proof}

Lemma~\ref{lem:scaling} reduces Theorem~\ref{thm:main} to double cones
of unit radius, whose algebras are generated by the Weyl operators of
sources in unit balls by \eqref{eq:time-slice}. Sections~\ref{sec:setting}--\ref{sec:evaluation} use an
arbitrary mass \(\mu>0\) and arbitrary length scales; from
Section~\ref{sec:forward} on, \(\mu\) denotes the dimensionless mass
\(mR\) and \(s=L/R\) the dimensionless separation of unit balls.

\subsection{Source Hilbert spaces}\label{subsec:sources}

For a nonempty bounded open \(U\subset\R^3\), let
\(\mathcal X_{\mu,U}\) be the completion of the real smooth source pairs
supported in \(U\), in the norm of \eqref{eq:covariance}. Its elements
are distributions in \(H^{-1/2}\oplus H^{1/2}\) supported in \(\overline U\). Here
\(H^\alpha(\R^3)\) is the Sobolev space defined by the Fourier weight
\((1+|\bp|^2)^\alpha\).
The complexification is \(\Hh_{\mu,U}\), with conjugation
\(\mathsf c_U\). Inner products are conjugate-linear in the first
variable.

The complexified source map is
\[
 (\kappa_U(f,g))(p)=\wh f(\bp)+ip_0\wh g(\bp).
\]
For real sources \(\sx=(f,g)\) and \(\sy=(f',g')\), the measure
\(d\mathfrak m_\mu\) and Plancherel's formula give
\begin{align*}
 \ip{\kappa_U\sx}{\kappa_U\sy}
 &=\frac1{(2\pi)^3}\int_{\R^3}\biggl(
   \frac{\overline{\wh f}\,\wh{f'}}{2p_0}
   +\frac{p_0\,\overline{\wh g}\,\wh{g'}}2
   +\frac i2\bigl(\overline{\wh f}\,\wh{g'}-\overline{\wh g}\,\wh{f'}\bigr)
   \biggr)\,d^3p\\
 &=\frac12\bigl(\ip{f}{h_\mu^{-1}f'}+\ip{g}{h_\mu g'}\bigr)
   +\frac i2\int_{\R^3}(fg'-gf')\,d^3x,
\end{align*}
that is,
\begin{equation}
 \ip{\kappa_U\sx}{\kappa_U\sy}
 =\vv_{\mu,U}(\sx,\sy)+\frac i2\varsigma(\sx,\sy).
 \label{eq:gram}
\end{equation}
Thus the real source map is isometric:
\(\norm{\kappa_U\sx}^2=\vv_{\mu,U}(\sx,\sx)\), since
\(\varsigma(\sx,\sx)=0\). Its complexification is bounded: with the norm
\(\norm{\sx+i\sy}^2=\vv_{\mu,U}(\sx,\sx)+\vv_{\mu,U}(\sy,\sy)\) of
\(\Hh_{\mu,U}\),
\[
 \norm{\kappa_U(\sx+i\sy)}^2
 =\norm{\kappa_U\sx}^2+\norm{\kappa_U\sy}^2
  -2\Ima\ip{\kappa_U\sx}{\kappa_U\sy}
 \overset{\eqref{eq:gram}}{=}\norm{\sx+i\sy}^2-\varsigma(\sx,\sy)
 \le2\norm{\sx+i\sy}^2,
\]
because
\(|\varsigma(\sx,\sy)|=2\lvert\Ima\ip{\kappa_U\sx}{\kappa_U\sy}\rvert
\le2\norm{\kappa_U\sx}\norm{\kappa_U\sy}\le\norm{\sx+i\sy}^2\).

\begin{lemma}[Standardness and injectivity]\label{lem:standard}
For every nonempty bounded open \(U\), the map \(\kappa_U\) from
\(\Hh_{\mu,U}\) to \(\Hh_{1,\mu}\) is injective and has dense complex range. Its real range
\(\mathcal K_U=\kappa_U\mathcal X_{\mu,U}\) is closed and satisfies
\begin{equation}
 \mathcal K_U\cap i\mathcal K_U=\{0\},\qquad
 \overline{\mathcal K_U+i\mathcal K_U}=\Hh_{1,\mu}.
 \label{eq:standard}
\end{equation}
\end{lemma}

\begin{proof}
\emph{Step 1: entire transforms.} The components
of an element of \(\Hh_{\mu,U}\) are distributions supported in the
compact set \(\overline U\), and by the Paley--Wiener--Schwartz theorem
their Fourier transforms extend from \(\R^3\) to entire functions on
\(\C^3\).

\smallskip
\emph{Step 2: injectivity.} If \(\kappa_U(f,g)=0\), then on real momenta
\[
 \wh f(\bp)+i\sqrt{\mu^2+|\bp|^2}\,\wh g(\bp)=0.
\]
For fixed real transverse momentum \(\bp_\perp\), squaring and
continuing in \(z=p_1\) gives
\[
 \wh f(z,\bp_\perp)^2
 =-(z^2+M^2)\wh g(z,\bp_\perp)^2,\qquad
 M^2=\mu^2+|\bp_\perp|^2>0,
\]
first for real \(z\), and then for all \(z\in\C\), because both sides
are entire functions of \(z\) by Step~1.

At \(z=iM\), the factor
\(z^2+M^2=(z-iM)(z+iM)\) has a simple zero, since \(M>0\). So a nonzero
right-hand side has odd vanishing order, whereas the left-hand side, a
square, has even vanishing order. Both sides therefore
vanish identically. Since \(\bp_\perp\) was arbitrary, \(f=g=0\).

\smallskip
\emph{Step 3: a vector orthogonal to the range.} For density, the
equivalent spatial \(L^2\) realization, obtained with the unitary map
\(\psi\mapsto\mathcal F^{-1}\bigl((2\pi)^{-3/2}(2p_0)^{-1/2}\psi\bigr)\)
of \(\Hh_{1,\mu}\) onto \(L^2(\R^3)\), is
\[
 \kappa_U(f,g)=2^{-1/2}
 \bigl(h_\mu^{-1/2}f+ih_\mu^{1/2}g\bigr).
\]
A vector \(v\) perpendicular to the complex local range gives a
distribution \(u=h_\mu^{-1/2}v\) with \(u=h_\mu u=0\) on \(U\): for all
\(f,g\in C_c^\infty(U)\), orthogonality to \(\kappa_U(f,0)\) and to
\(\kappa_U(0,g)\) reads
\[
 \ip{u}{f}=\ip{h_\mu^{-1/2}v}{f}=0,\qquad
 \ip{h_\mu u}{g}=\ip{h_\mu^{1/2}v}{g}=0,
\]
where \(u\in H^{1/2}(\R^3)\) and \(h_\mu u\in H^{-1/2}(\R^3)\).

\smallskip
\emph{Step 4: the elliptic continuation.} The function
\(w(\tau)=e^{-\tau h_\mu}u\), \(\tau>0\), is smooth and solves
\[
 (-\partial_\tau^2-\nabla^2+\mu^2)w=0,
\]
and, as \(\tau\downarrow0\), \(w(\tau)\to u\) in \(H^{1/2}(\R^3)\) and
\(\partial_\tau w(\tau)=-h_\mu w(\tau)\to-h_\mu u\) in
\(H^{-1/2}(\R^3)\). Its Dirichlet and Neumann traces, \(u\) and
\(-h_\mu u\), vanish on any ball \(B\) contained in \(U\).

Extension by
zero across the boundary patch \(\{0\}\times B\) is a weak solution of the same elliptic
equation on the connected open set
\(\mathcal D=((0,\infty)\times\R^3)\cup(\R\times B)\). Indeed, for
\(\varphi\in C_c^\infty(\mathcal D)\) and \(\delta>0\), integration by
parts on \([\delta,\infty)\times\R^3\) moves the operator onto \(w\),
where it gives zero, and leaves the boundary terms at \(\tau=\delta\):
\[
 \int_\delta^\infty
 \ip{w(\tau)}{(-\partial_\tau^2-\nabla^2+\mu^2)\varphi(\tau)}\,d\tau
 =\ip{w(\delta)}{\partial_\tau\varphi(\delta)}
 -\ip{\partial_\tau w(\delta)}{\varphi(\delta)}.
\]
As \(\delta\downarrow0\), the left-hand side tends to the pairing of the
extension with \((-\partial_\tau^2-\nabla^2+\mu^2)\varphi\), and the
right-hand side tends to
\(\ip{u}{\partial_\tau\varphi(0)}+\ip{h_\mu u}{\varphi(0)}\), a sum of
pairings of \(H^{1/2}\) with \(H^{-1/2}\), which vanishes by Step~3
because \(\varphi(0)\) and \(\partial_\tau\varphi(0)\) are supported in
\(B\subset U\).

The extension vanishes on the open subset
\((-\infty,0)\times B\) of \(\mathcal D\). Analytic elliptic regularity
and uniqueness \cite{Hormander03} imply that it vanishes on the
connected domain. Hence \(w(\tau)=0\) for \(\tau>0\),
\(u=\lim_{\tau\downarrow0}w(\tau)=0\) in \(H^{1/2}(\R^3)\), and
\(v=h_\mu^{1/2}u=0\).

\smallskip
\emph{Step 5: standardness.} The real isometry has closed range,
because \(\mathcal X_{\mu,U}\) is complete. Injectivity of the
complexification gives the first relation in \eqref{eq:standard}: if
\(\kappa_U\sx=i\kappa_U\sy\) with \(\sx,\sy\in\mathcal X_{\mu,U}\), then
\(\kappa_U(\sx-i\sy)=0\), so \(\sx=\sy=0\) by Step~2. Density gives the
second, since \(\mathcal K_U+i\mathcal K_U=\kappa_U\Hh_{\mu,U}\), and
this complex range is dense by Steps~3 and~4.

Standardness of local subspaces of the free field is known; for
regions with \(C^1\) boundary, Figliolini and Guido prove it in a similar
way, from the anti-locality of \(h_\mu\)
\cite[Propositions~2.4 and~2.7]{FG89}.
\end{proof}

\subsection{Finite canonical systems and local approximation}

A finite-dimensional real subspace of smooth local sources is a
\emph{canonical source system} when the restriction of
\eqref{eq:symplectic} is nondegenerate. A symplectic basis gives
quadratures with commutation matrix
\[
 J_n=\bigoplus_{j=1}^n
 \begin{pmatrix}0&1\\-1&0\end{pmatrix}.
\]

By the Stone--von Neumann theorem, the regular Weyl representation of a
canonical source system is
the oscillator representation with a possible multiplicity
\cite{Haag96}. The vacuum restriction is a Gaussian density operator.
Its covariance \(V\) satisfies
\begin{equation}
 V+\frac i2J_n>0.
 \label{eq:strict-uncertainty}
\end{equation}

For the symplectic basis
\(\sx_1,\dots,\sx_{2n}\), supported in a bounded open set \(U\), the
covariance is \(V_{kl}=\vv_{\mu,U}(\sx_k,\sx_l)\), since
\(\omega(\mathsf W_\mu(\sum_k\xi_k\sx_k))
=\exp(-\frac12\sum_{k,l}\xi_k\xi_l\vv_{\mu,U}(\sx_k,\sx_l))\) for
\(\xi\in\R^{2n}\), and \(\varsigma(\sx_k,\sx_l)=(J_n)_{kl}\); so
\[
 \Bigl(V+\frac i2J_n\Bigr)_{kl}
 =\vv_{\mu,U}(\sx_k,\sx_l)+\frac i2\varsigma(\sx_k,\sx_l)
 \overset{\eqref{eq:gram}}{=}\ip{\kappa_U\sx_k}{\kappa_U\sx_l}.
\]
For \(c\in\C^{2n}\setminus\{0\}\), the vector \(\sum_lc_l\sx_l\) of
\(\Hh_{\mu,U}\) is nonzero; since
\(\kappa_U\) is injective by Lemma~\ref{lem:standard},
\(c^*(V+\frac i2J_n)c=\norm{\kappa_U\sum_lc_l\sx_l}^2>0\).

Williamson's normal form \cite{Williamson36} writes
\[
 V=S\left(\bigoplus_{j=1}^n\upsilon_j I_2\right)S^T,\qquad
 SJ_nS^T=J_n.
\]
The numbers \(\upsilon_j\) are the symplectic eigenvalues. Relation \eqref{eq:strict-uncertainty} therefore
gives \(\upsilon_j>1/2\), so the Gaussian density operator is faithful.

Every finite collection of smooth sources supported in an open set
\(U\) lies in a finite canonical system of sources supported in \(U\).
The form \(\varsigma\) is nondegenerate on the
smooth sources supported in \(U\): a nonzero \(\sx=(f,g)\) has the
partner \((0,f)\) if \(f\ne0\), and \((-g,0)\) if \(f=0\), with
\[
 \varsigma\bigl((f,g),(0,f)\bigr)=\int_{\R^3}f^2\,d^3x,\qquad
 \varsigma\bigl((0,g),(-g,0)\bigr)=\int_{\R^3}g^2\,d^3x.
\]

Let \(\mathcal C\) be the span of the collection,
\(\mathsf r_1,\dots,\mathsf r_k\) a basis of its radical
\(\{\sx\in\mathcal C:\varsigma(\sx,\sy)=0\text{ for all }\sy\in\mathcal C\}\),
and \(\mathcal C_0\) a complement of the radical in \(\mathcal C\), on
which \(\varsigma\) is nondegenerate. Since \(\mathcal C_0\) is
finite-dimensional, \(\varsigma\) is also nondegenerate on the
symplectic complement \(\mathcal C_0^{\varsigma}\) of \(\mathcal C_0\)
among the smooth sources supported in \(U\), which contains the
radical; so there are local partners
\(\sy_1,\dots,\sy_k\in\mathcal C_0^{\varsigma}\) with
\(\varsigma(\mathsf r_i,\sy_j)=\delta_{ij}\). The matrix of
\(\varsigma\) on \(\mathsf r_1,\dots,\mathsf r_k,\sy_1,\dots,\sy_k\) is
\[
 \begin{pmatrix}0&I_k\\-I_k&\ast\end{pmatrix},
\]
which is invertible. So these vectors form a basis of their span,
\(\varsigma\) is nondegenerate on the span, and the span is
\(\varsigma\)-orthogonal to \(\mathcal C_0\). Hence \(\varsigma\) is
nondegenerate on the sum of \(\mathcal C_0\) and the span, a finite
canonical system that contains \(\mathcal C\). The finite
canonical systems in \(U\) form a directed family: two of them lie in the system obtained
in this way from their sum.

\begin{lemma}[Passage to all local effects]\label{lem:exhaustion}
Suppose an inequality involving
\[
 \omega(E),\qquad \omega(P),\qquad \omega(EP)
\]
holds, with constants independent of the canonical systems, for every
effect \(E\) of the von Neumann algebra generated by the Weyl operators
of a finite canonical source system in a first region and every effect
\(P\) of such an algebra in a second region. If the inequality has the
form \(\Phi(\omega(E),\omega(P),\omega(EP))\le0\) with a continuous
function \(\Phi\), it holds for every pair of effects in the full
regional algebras, generated by the Weyl operators of all smooth
sources in the two regions (for two balls, the double-cone algebras of
\eqref{eq:time-slice}). The same conclusion holds for a bound on
\(|\omega(AB)-\omega(A)\omega(B)|\), with pairs of self-adjoint
contractions \(A,B\) in place of the effects \(E,P\).
\end{lemma}

\begin{proof}
\emph{Step 1: a dense union.} Let \(\mathfrak M\) be the regional
algebra of the first region and \(\mathfrak M_0\) the union of the von
Neumann algebras of the finite canonical source systems in that region.
The finite canonical Weyl algebras generate the local algebras: every smooth source of the region lies in a finite
canonical system of the region, so \(\mathfrak M_0\subset\mathfrak M\)
contains the Weyl operators that generate \(\mathfrak M\). The directed
union \(\mathfrak M_0\) is a strongly dense unital \(*\)-algebra: two finite canonical
systems lie in a third, so sums and products of elements of
\(\mathfrak M_0\) lie in \(\mathfrak M_0\), and by the bicommutant
theorem the strong closure of \(\mathfrak M_0\) is
\(\mathfrak M_0''=\mathfrak M\).

\smallskip
\emph{Step 2: effects of the union.} For
an effect \(E\in\mathfrak M\), the operator \(2E-I\) is a self-adjoint
contraction in \(\mathfrak M\). Kaplansky density
\cite[Theorem~2.4.16]{BR87} gives a net of
self-adjoint contractions \(X_\alpha\in\mathfrak M_0\) with \(X_\alpha\to2E-I\)
strongly, and the affine map \(t\mapsto(1+t)/2\) turns them
into effects \(E_\alpha=(I+X_\alpha)/2\) of \(\mathfrak M_0\) with
\(E_\alpha\to E\) strongly. Each \(E_\alpha\) is an effect of the von
Neumann algebra of one finite canonical system.

\smallskip
\emph{Step 3: joint convergence.} The approximations in the two regions
can be chosen independently: let \(P_\beta\to P\) be obtained in the
same way in the second region. Products of bounded strongly convergent
nets converge strongly. So the
three vacuum expectations converge along the pairs \((\alpha,\beta)\):
\(\omega(E_\alpha)\to\omega(E)\), \(\omega(P_\beta)\to\omega(P)\) and
\(\omega(E_\alpha P_\beta)\to\omega(EP)\).

\smallskip
\emph{Step 4: the limit.} By hypothesis,
\(\Phi(\omega(E_\alpha),\omega(P_\beta),\omega(E_\alpha P_\beta))\le0\)
for all \(\alpha,\beta\), and the continuity of \(\Phi\) gives
\(\Phi(\omega(E),\omega(P),\omega(EP))\le0\). The argument for
self-adjoint contractions is the same, with Kaplansky density applied to
\(A\) and \(B\) themselves.
\end{proof}

\subsection{Modular operators and uncertainty reserves}

A closed real subspace satisfying \eqref{eq:standard} is
\emph{standard}. The Tomita operator and modular operator of a standard
subspace, written here for \(\mathcal K_U\), are
\[
 \mathsf s_U(u+iv)=u-iv\quad(u,v\in\mathcal K_U),\qquad
 \Delta_U=\mathsf s_U^*\mathsf s_U.
\]
We use the standard-subspace formulation of modular theory
\cite{EO73,RvD77}.

For an operator \(T\) on the source space,
\(\overline T=\mathsf c_UT\mathsf c_U\). The source Gram operators and
their \emph{uncertainty reserve} are
\begin{equation}
 G_U=\kappa_U^*\kappa_U,\qquad
 \overline G_U=2I-G_U,\qquad
 \Res_U=(G_U\overline G_U)^{1/2}.
 \label{eq:reserve}
\end{equation}
The middle formula, an identity for
\(\overline G_U=\mathsf c_UG_U\mathsf c_U\), follows from \eqref{eq:gram}: for real
\(\sx,\sy\), the conjugation is antiunitary and fixes real vectors, so
\[
 \ip{\sx}{\overline G_U\sy}
 =\overline{\ip{\sx}{G_U\sy}}
 =\vv_{\mu,U}(\sx,\sy)-\frac i2\varsigma(\sx,\sy)
 =2\vv_{\mu,U}(\sx,\sy)-\ip{\sx}{G_U\sy}
 =\ip{\sx}{(2I-G_U)\sy},
\]
and complex linearity extends the identity to \(\Hh_{\mu,U}\).

Both Gram operators are positive and injective, and they commute:
\(\ip{\sx}{G_U\sx}=\norm{\kappa_U\sx}^2\) and
\(\ip{\sx}{\overline G_U\sx}=\norm{\kappa_U\mathsf c_U\sx}^2\), where
\(\kappa_U\) is injective by Lemma~\ref{lem:standard}, and
\(\overline G_U=2I-G_U\) commutes with \(G_U\). So the spectrum of
\(G_U\) lies in \([0,2]\), and
\[
 G_U\overline G_U=G_U(2I-G_U)=I-(I-G_U)^2\le I.
\]
Since \(G_U\) and
\(2I-G_U\) are injective, \(G_U\overline G_U\) is injective. Hence
\(0<\Res_U\le I\); the spectrum of \(\Res_U\) can accumulate at
zero. The reserve
measures the room left by the uncertainty relation: \(\Res_U\) equals
\(I\) on source directions along which the symplectic form vanishes
relative to the covariance, where \(G_U=I\), and \(\Res_U\) is small
on directions along which the covariance nearly saturates the
uncertainty bound, where \(G_U\) is close to \(0\) or to \(2\).

In the polar decomposition
\(\kappa_U=\mathcal U_UG_U^{1/2}\), the operator \(\mathcal U_U\)
is unitary by Lemma~\ref{lem:standard}, since \(\kappa_U\) is injective
and has dense range. The relation
\(\mathsf s_U\kappa_U\sx=\kappa_U\mathsf c_U\sx\), which is the
definition of \(\mathsf s_U\) applied to \(u=\kappa_U\sx_1\) and
\(v=\kappa_U\sx_2\) for \(\sx=\sx_1+i\sx_2\) with real \(\sx_1,\sx_2\),
gives, for every \(\sx\in\Hh_{\mu,U}\),
\[
 \norm{\Delta_U^{1/2}\kappa_U\sx}^2
 =\norm{\mathsf s_U\kappa_U\sx}^2
 =\norm{\kappa_U\mathsf c_U\sx}^2
 \overset{\eqref{eq:reserve}}{=}\ip{\mathsf c_U\sx}{G_U\mathsf c_U\sx}
 =\ip{\sx}{\overline G_U\sx}.
\]
The first equality is the polar decomposition of \(\mathsf s_U\), whose
antiunitary factor preserves norms, and the last holds because
\(\mathsf c_U\) is antiunitary and
\(\overline G_U=\mathsf c_UG_U\mathsf c_U\).

Let
\(\overline G_UG_U^{-1}\) and its powers be defined by the joint
spectral calculus of the commuting operators \(G_U\) and
\(\overline G_U\). With \(\xi=G_U^{1/2}\sx\), so that
\(\kappa_U\sx=\mathcal U_U\xi\) and \(\sx=G_U^{-1/2}\xi\), the identity for
\(\lVert\Delta_U^{1/2}\kappa_U\sx\rVert^2\) reads
\[
 \norm{\bigl(\mathcal U_U^*\Delta_U\mathcal U_U\bigr)^{1/2}\xi}^2
 =\norm{\Delta_U^{1/2}\mathcal U_U\xi}^2
 =\ip{G_U^{-1/2}\xi}{\overline G_UG_U^{-1/2}\xi}
 =\norm{\bigl(\overline G_UG_U^{-1}\bigr)^{1/2}\xi}^2
\]
for \(\xi\in\Ran G_U^{1/2}=\Dom G_U^{-1/2}\).
This space is the form domain of both operators: that of
\(\mathcal U_U^*\Delta_U\mathcal U_U\) is
\(\mathcal U_U^*\Dom\Delta_U^{1/2}=\mathcal U_U^*\Dom\mathsf s_U
=\mathcal U_U^*\kappa_U\Hh_{\mu,U}=\Ran G_U^{1/2}\), and that of
\(\overline G_UG_U^{-1}\) is \(\Dom G_U^{-1/2}\), because
\((2-\lambda)/\lambda\le2/\lambda=(2-\lambda)/\lambda+1\) for
\(0<\lambda\le2\). Positive self-adjoint operators with the same form
domain and the same quadratic form coincide, so
\(\mathcal U_U^*\Delta_U\mathcal U_U=\overline G_UG_U^{-1}\), and
\begin{equation}
 \Delta_U=\mathcal U_U\overline G_UG_U^{-1}\mathcal U_U^*,
 \qquad
 \Delta_U^c\kappa_U=\mathcal U_Uh_{c,U},
 \qquad
 h_{c,U}=G_U^{1/2-c}\overline G_U^{\,c}.
 \label{eq:modular-source}
\end{equation}

For \(0\le c\le1/2\), \(h_{c,U}\) is bounded, because the function
\(\lambda^{1/2-c}(2-\lambda)^c\) is bounded on \([0,2]\). The second formula holds on the whole source
space: \(\kappa_U\Hh_{\mu,U}=\Dom\Delta_U^{1/2}\subset\Dom\Delta_U^c\),
and the first formula gives
\[
 \Delta_U^c\kappa_U
 =\mathcal U_U\bigl(\overline G_UG_U^{-1}\bigr)^c\mathcal U_U^*\,
  \mathcal U_UG_U^{1/2}
 =\mathcal U_UG_U^{1/2-c}\overline G_U^{\,c}.
\]

In the notation of Rieffel and van Daele
\cite[Sections~2 and~3]{RvD77}, with \(P\) and \(Q\) the real orthogonal
projections onto \(\mathcal K_U\) and \(i\mathcal K_U\), their operators
\(R=P+Q\), \(2-R\) and \(T=R^{1/2}(2-R)^{1/2}\) are
\(\mathcal U_UG_U\mathcal U_U^*\), \(\mathcal U_U\overline G_U\mathcal U_U^*\)
and \(\mathcal U_U\Res_U\mathcal U_U^*\), because
\(P+Q=\kappa_U\kappa_U^*\), and the first formula in
\eqref{eq:modular-source} is their relation \(\Delta=(2-R)R^{-1}\).
For quasi-free states, Lechner and Sanders obtain the same two formulas
from the same relation between the Tomita operator and the conjugation,
in their notation with \(G_U=1+\Sigma\) and \(\overline G_U=1-\Sigma\)
\cite[Proposition~5.2 and the proof of Proposition~5.3]{LS16}.

\begin{lemma}[Containing-region inequality]\label{lem:inclusion}
If \(\mathcal K_U\subset\mathcal K_W\) are standard real subspaces,
\(\Delta_W\) is the modular operator of \(\mathcal K_W\), and
\(0<c<1/2\), then every \(\sx\in\Hh_{\mu,U}\) satisfies
\begin{equation}
 \norm{\Delta_W^c\kappa_U\sx}\le\norm{h_{c,U}\sx}.
 \label{eq:inclusion}
\end{equation}
\end{lemma}

\begin{proof}
\emph{Step 1: the form order.} The inclusion gives
\(\mathsf s_U\subset\mathsf s_W\): the domain
\(\mathcal K_U+i\mathcal K_U\) of \(\mathsf s_U\) lies in the domain
\(\mathcal K_W+i\mathcal K_W\) of \(\mathsf s_W\), and both operators map
\(u+iv\) to \(u-iv\). The quadratic form of \(\Delta_U\) is the
restriction of the form of \(\Delta_W\): by the polar decompositions of
\(\mathsf s_U\) and \(\mathsf s_W\), every
\(\xi\in\Dom\Delta_U^{1/2}=\Dom\mathsf s_U\) satisfies
\[
 \norm{\Delta_U^{1/2}\xi}=\norm{\mathsf s_U\xi}
 =\norm{\mathsf s_W\xi}=\norm{\Delta_W^{1/2}\xi}.
\]
Therefore \(\Delta_U\ge\Delta_W\) in form order:
\[
 \Dom\Delta_U^{1/2}\subset\Dom\Delta_W^{1/2},\quad\text{and}\quad
 \norm{\Delta_W^{1/2}\xi}\le\norm{\Delta_U^{1/2}\xi}\quad\text{for }
 \xi\in\Dom\Delta_U^{1/2}.
\]

\smallskip
\emph{Step 2: the Löwner--Heinz inequality.} The
Löwner--Heinz inequality for \(2c\in(0,1)\), in Heinz's form for
unbounded positive operators ordered as quadratic forms
\cite[Satz~2]{Heinz51}, gives
\(\Delta_U^{2c}\ge\Delta_W^{2c}\). In form order, the inequality
\(\Delta_U^{2c}\ge\Delta_W^{2c}\) says that \(\Dom\Delta_U^c\subset\Dom\Delta_W^c\) and
\[
 \norm{\Delta_W^c\xi}\le\norm{\Delta_U^c\xi},\qquad
 \xi\in\Dom\Delta_U^c.
\]

\smallskip
\emph{Step 3: the source bound.} For \(\sx\in\Hh_{\mu,U}\), the vector
\(\kappa_U\sx\) lies in \(\Dom\mathsf s_U=\Dom\Delta_U^{1/2}\), which is
contained in \(\Dom\Delta_U^c\) since \(c<1/2\). Step~2 at
\(\xi=\kappa_U\sx\) gives
\eqref{eq:inclusion}:
\[
 \norm{\Delta_W^c\kappa_U\sx}
 \overset{\text{Step 2}}{\le}\norm{\Delta_U^c\kappa_U\sx}
 \overset{\eqref{eq:modular-source}}{=}\norm{\mathcal U_Uh_{c,U}\sx}
 =\norm{h_{c,U}\sx},
\]
the last step because \(\mathcal U_U\) is unitary.
Equivalently, the closure of \(\Delta_W^c\Delta_U^{-c}\) is a
contraction, which is the bound of \cite[Section~2]{BDL07} for the
inclusion \(\mathcal K_U\subset\mathcal K_W\) and the exponent \(c\).
\end{proof}

We write \(\Sc_j\) for the Schatten class, with
\(\norm{T}_j=(\Tr|T|^j)^{1/j}\) for real \(1\le j<\infty\).
The unadorned operator norm is \(\norm{T}\).
Our uses of Schatten Hölder inequalities and complex interpolation
follow~\cite{Simon05}.

\begin{convention}[Negative reserve powers]\label{conv:factorization}
An expression such as
\[
 \Res_U^bX\Res_{U'}^b,\qquad X\in\Sc_1,
\]
with \(b>0\) and the uncertainty reserves \eqref{eq:reserve} of nonempty
bounded open sets \(U,U'\subset\R^3\), denotes the
bounded operator obtained by composition, and an equation between a
bounded operator and such an expression is an equality of bounded
operators. Such an equation takes the place of multiplying the bounded
operator by \(\Res_U^{-b}\) on the left and by \(\Res_{U'}^{-b}\) on the
right. For a linear functional
\(\ell\), the notation \(\lVert\ell\Res_U^{-b}\rVert\le M\) means
\[
 |\ell(\sx)|\le M\norm{\Res_U^b\sx},
\]
for every \(\sx\in\Hh_{\mu,U}\), followed by continuous extension of \(\ell\Res_U^{-b}\)
from \(\Ran\Res_U^b\). All negative
reserve powers below have this factorization interpretation.
\end{convention}

Throughout, \(C\) denotes a positive universal constant whose value may
change from line to line. Dependence on a fixed parameter is indicated
by a subscript. Named constants such as \(C_*\), \(C_c\), \(C_1\), \(C_2\),
\(C_{\mathrm{sub}}\) and \(C_{\mathrm{prob}}\) keep their values throughout, and
\(C_{\mu,s}\) in Section~\ref{sec:forward} is an operator.

Table~\ref{tab:constants} in Section~\ref{sec:completion} lists the
constants that enter \(C_*\). Some letters occur in two fonts with
different meanings: \(E\) and \(\Energy\) (an effect and the energy
scale), \(P\) and \(\mathfrak P_\mu\) (an effect and the momentum
covariance kernel), \(\mathcal X_{\mu,U}\) and \(\mathfrak X_\mu\) (the
real source space and the field covariance kernel), \(R\) and \(\Res_U\)
(the radius and the uncertainty reserve), and \(J_n\) and \(\mathsf J\)
(the canonical commutation matrix and the normalized commutation form).

\section{A probability criterion for the classical CHSH bound}
\label{sec:probability}

Let \(\omega\) be any state on the von Neumann algebra generated by two
commuting von Neumann algebras \(\A\) and \(\B\). For \(r>1\), the event comparison constants are
\begin{equation}
 K_{r,+}
 =\sup_{E,P}
 \frac{\omega(EP)^r}{(\omega(E)\omega(P))^{r-1}},
 \qquad
 K_{r,-}
 =\sup_{E,P}
 \frac{(\omega(E)\omega(P))^r}{\omega(EP)^{r-1}}.
 \label{eq:K}
\end{equation}
The suprema range over projections \(E\in\A\) and \(P\in\B\) with
\(\omega(E)\omega(P)>0\). A positive numerator with zero denominator gives
\(+\infty\). At \(E=P=I\) each quotient equals \(1\), so both
constants are at least \(1\).

\begin{theorem}[Two event bounds and clustering]\label{thm:probability}
Let \(r_+>2\) and \(1<r_-<2\), and set
\begin{equation}
 b=1-\frac1{r_+},\qquad
 a=1-\frac1{r_-},\qquad
 \eta=1-\frac{1-b}{ab^2}.
 \label{eq:probability-exponents}
\end{equation}
Suppose \(\eta>0\), \(K_{r_+,+},K_{r_-,-}<\infty\), and, for some
\(q_{\mathrm{cl}}\ge0\),
\begin{equation}
 |\omega(AB)-\omega(A)\omega(B)|\le q_{\mathrm{cl}}
 \label{eq:additive-cluster}
\end{equation}
for all self-adjoint contractions \(A\in\A\) and \(B\in\B\). Define
\begin{equation}
 \nu_+=\frac{1+b}{ar_+b^2},\qquad
 \nu_-=\frac{1-a}{a}.
 \label{eq:nus}
\end{equation}
If
\begin{equation}
 q_{\mathrm{cl}}^\eta
 K_{r_+,+}^{\nu_+}K_{r_-,-}^{\nu_-}
 \le 2^{-2/(ab)},
 \label{eq:probability-cutoff}
\end{equation}
then the Summers--Werner Bell correlation \(\beta(\omega;\A,\B)\) of
\eqref{eq:beta} equals \(1\).
\end{theorem}

\begin{proof}
\emph{Step 1: reduction to quartets of reflections.} Identities attain
the value \(1\) in \eqref{eq:beta}, so \(\beta(\omega;\A,\B)=1\) unless
some self-adjoint contractions \(A_0,A_1\in\A\) and \(B_0,B_1\in\B\) give
the CHSH expression an absolute value larger than \(1\). Since \(A_i\)
and \(B_j\) commute, the expression is real, and negating both
observables at one site changes its sign; a value smaller than \(-1\)
thus gives a value larger than \(1\).

Every self-adjoint contraction is a spectral average of reflections:
\[
 A=\frac12\int_{-1}^1
 \bigl(2\,\mathbf1_{\{A\ge u\}}-I\bigr)\,du.
\]
The reflections \(2\,\mathbf1_{\{A\ge u\}}-I\) lie in every von Neumann
algebra that contains \(A\), and the integral converges in norm, because
its Riemann sums are step functions of \(A\) whose scalar counterparts
converge to \(t\mapsto t\) uniformly on \([-1,1]\).

The CHSH expression is separately affine in its four observables and
continuous in norm. With \(A_1,B_0,B_1\) fixed, its value at \(A_0\) is
therefore a limit of averages, with weights of total mass \(1\), of its
values at the reflections \(2\,\mathbf1_{\{A_0\ge u\}}-I\), and a value
larger than \(1\) at \(A_0\) gives a value larger than \(1\) at one of
these reflections. Replacing \(A_0\), \(A_1\), \(B_0\) and \(B_1\) in
turn in this way turns a strict violation into a strict violation by
reflections.

For a quartet of reflections \(A_0,A_1\in\A\) and \(B_0,B_1\in\B\),
the \emph{score} is the value
\[
 S=\tfrac12\omega(A_0(B_0+B_1)+A_1(B_0-B_1))
\]
of the CHSH expression in \eqref{eq:beta}.

\smallskip
\emph{Step 2: the winning cells and their symmetries.} The
\emph{cells} of the experiment are the sixteen joint probabilities
\(\omega(E_i^\pm P_j^\pm)\) of the outcome effects
\(E_i^\pm=(I\pm A_i)/2\) and \(P_j^\pm=(I\pm B_j)/2\), which are
projections because \(A_i\) and \(B_j\) are reflections.
A cell is \emph{winning} when the product of its outcomes agrees with
the sign of the corresponding correlator in the expanded score
\[
 S=\tfrac12\bigl(\omega(A_0B_0)+\omega(A_0B_1)+\omega(A_1B_0)
 -\omega(A_1B_1)\bigr).
\]
The winning cells of the setting pair \((i,j)\) are thus
\(\omega(E_i^+P_j^+)\) and \(\omega(E_i^-P_j^-)\) if \((i,j)\ne(1,1)\),
and \(\omega(E_1^+P_1^-)\) and \(\omega(E_1^-P_1^+)\) if \((i,j)=(1,1)\).

Let \(p_{ij}\) be the sum of the two winning cells of the setting pair
\((i,j)\). Since \(A_iB_j=(E_i^+-E_i^-)(P_j^+-P_j^-)\) and the four cells
of a setting pair add up to \(1\), the correlator \(\omega(A_iB_j)\), with
its sign in \(S\), equals \(p_{ij}-(1-p_{ij})\), and the sum over the four
setting pairs is
\[
 S=\frac12\sum_{i,j=0}^1\bigl(2p_{ij}-1\bigr)
 =\sum_{i,j=0}^1p_{ij}-2,
\]
so \(S+2\) is the sum of the eight winning cells.

The maps that replace \((A_0,A_1,B_0,B_1)\) by
\[
 (A_0,-A_1,B_1,B_0),\qquad
 (A_1,A_0,B_0,-B_1),\qquad
 (-A_0,-A_1,-B_0,-B_1)
\]
carry reflection quartets to reflection quartets and permute the eight
winning cells. The first exchanges the two settings in \(\B\), which
moves the cells of the setting pair \((i,j)\) to \((i,1-j)\), and flips
the outcome of \(A_1\), which compensates the opposite signs of
\(\omega(A_1B_0)\) and \(\omega(A_1B_1)\) in \(S\). The second exchanges
the two settings in \(\A\), which moves the cells of \((i,j)\) to
\((1-i,j)\), and flips the outcome of \(B_1\), which compensates the
opposite signs of \(\omega(A_0B_1)\) and \(\omega(A_1B_1)\). The third
flips the outcomes at both sites, which exchanges the two winning cells
of each setting pair.

The three maps therefore leave \(S\) unchanged, and the eight winning
cells form one orbit under them: compositions of the first two carry
any setting pair to \((0,0)\), and the third exchanges the two winning
cells of \((0,0)\). Replacing the quartet by its image under a
composition of the three maps, we may assume that a least likely
winning cell is \(\omega(E_0^-P_0^-)\), the cell of the outcomes \(-1\)
of \(A_0\) and \(B_0\).

\smallskip
\emph{Step 3: the probability form of the score.} By Step~2, the least
probability of a winning cell is \(w=\omega(E_0^-P_0^-)\), and we show
that \(S>1\) takes the probability form
\begin{equation}
 w>
 \omega(E_0^-P_1^+)
 +\omega(E_1^+P_0^-)
 +\omega(E_1^-P_1^-).
 \label{eq:probability-CHSH}
\end{equation}

We write
\[
 e_A=\omega(E_0^-),\qquad e_B=\omega(P_0^-),\qquad
 x_A=\omega(E_1^-),\qquad x_B=\omega(P_1^-).
\]
Since \(A_i=I-2E_i^-\) and \(B_j=I-2P_j^-\), and
\(\omega((I-2E)(I-2P))=1-2\omega(E)-2\omega(P)+4\omega(EP)\) for all
effects \(E,P\), the four correlators are
\begin{align*}
 \omega(A_0B_0)&=1-2e_A-2e_B+4w, &
 \omega(A_0B_1)&=1-2e_A-2x_B+4\omega(E_0^-P_1^-),\\
 \omega(A_1B_0)&=1-2x_A-2e_B+4\omega(E_1^-P_0^-), &
 \omega(A_1B_1)&=1-2x_A-2x_B+4\omega(E_1^-P_1^-).
\end{align*}
We insert them into the expanded score of Step~2, in which the terms in
\(x_A\) and \(x_B\) cancel, and use
\(\omega(E_0^-P_1^+)=e_A-\omega(E_0^-P_1^-)\) and
\(\omega(E_1^+P_0^-)=e_B-\omega(E_1^-P_0^-)\), which hold because
\(P_1^+=I-P_1^-\) and \(E_1^+=I-E_1^-\):
\begin{align*}
 S&=1-2e_A-2e_B+2w+2\omega(E_0^-P_1^-)+2\omega(E_1^-P_0^-)
  -2\omega(E_1^-P_1^-)\\
 &=1+2\left\{
 w-\omega(E_0^-P_1^+)
  -\omega(E_1^+P_0^-)-\omega(E_1^-P_1^-)\right\}.
\end{align*}
Since \(S>1\), the expression in braces is positive, and
\eqref{eq:probability-CHSH} follows. The three terms
on the right of \eqref{eq:probability-CHSH} are values of \(\omega\) at
products of commuting projections, hence nonnegative; so \(w>0\), and
each of the three terms is smaller than \(w\).

\smallskip
\emph{Step 4: the two neighbouring winning cells.} The cells
\(E_0^-P_1^-\) and \(E_1^-P_0^-\) are winning, because
\(\omega(A_0B_1)\) and \(\omega(A_1B_0)\) have the sign
\(+\) in \(S\); so their probabilities are at least \(w\). By Step~3, the
first two terms on the right of \eqref{eq:probability-CHSH} are smaller
than \(w\), so
\[
 \omega(E_0^-P_1^-)=e_A-\omega(E_0^-P_1^+)>e_A-w,\qquad
 \omega(E_1^-P_0^-)=e_B-\omega(E_1^+P_0^-)>e_B-w.
\]
Hence
\[
 \omega(E_0^-P_1^-)\ge\max\{w,e_A-w\}\ge e_A/2,\qquad
 \omega(E_1^-P_0^-)\ge\max\{w,e_B-w\}\ge e_B/2.
\]

\smallskip
\emph{Step 5: the two event bounds and the marginals \(x_A,x_B\).} From
\eqref{eq:K}, for all projections \(E\in\A\) and \(P\in\B\),
\begin{equation}
 \omega(EP)\le k\bigl(\omega(E)\omega(P)\bigr)^b,\qquad
 \omega(EP)\ge\lambda
       \bigl(\omega(E)\omega(P)\bigr)^{1/a},
 \label{eq:event-orders}
\end{equation}
where
\[
 k=K_{r_+,+}^{1/r_+},\qquad
 \lambda=K_{r_-,-}^{-1/(r_--1)}.
\]
If \(\omega(E)\omega(P)>0\), the two quotients in \eqref{eq:K} at \(E,P\)
are at most \(K_{r_+,+}\) and \(K_{r_-,-}\). The second quotient has a
positive numerator; since a zero denominator would make it \(+\infty\)
and \(K_{r_-,-}<\infty\), its denominator is positive, that is,
\(\omega(EP)>0\). Multiplying out the two bounds on the quotients gives
\[
 \omega(EP)^{r_+}\le K_{r_+,+}\bigl(\omega(E)\omega(P)\bigr)^{r_+-1},
 \qquad
 \omega(EP)^{r_--1}\ge K_{r_-,-}^{-1}\bigl(\omega(E)\omega(P)\bigr)^{r_-},
\]
and the roots of orders \(r_+\) and \(r_--1\) give \eqref{eq:event-orders},
since \((r_+-1)/r_+=b\) and \(r_-/(r_--1)=1/a\) by
\eqref{eq:probability-exponents}. If \(\omega(E)\omega(P)=0\), then
\(\omega(EP)=0\), because the projection \(EP\) satisfies \(EP\le E\) and
\(EP\le P\), and both bounds in \eqref{eq:event-orders} hold trivially.

The upper bound in \eqref{eq:event-orders}, applied to the two
neighbouring winning cells, gives
\[
 \frac{e_A}2\overset{\text{Step 4}}{\le}\omega(E_0^-P_1^-)
 \overset{\eqref{eq:event-orders}}{\le}k(e_Ax_B)^b,\qquad
 \frac{e_B}2\overset{\text{Step 4}}{\le}\omega(E_1^-P_0^-)
 \overset{\eqref{eq:event-orders}}{\le}k(x_Ae_B)^b.
\]
Since \(e_A\ge\omega(E_0^-P_0^-)=w>0\) and likewise \(e_B>0\), division by
\(ke_A^b\) and by \(ke_B^b\), followed by the power \(1/b\), gives
\[
 x_B\ge(2k)^{-1/b}e_A^{(1-b)/b},\qquad
 x_A\ge(2k)^{-1/b}e_B^{(1-b)/b}.
\]

\smallskip
\emph{Step 6: a lower bound on \(w\).} At the cell \(E_0^-P_0^-\), the upper
bound in \eqref{eq:event-orders} reads \(w\le k(e_Ae_B)^b\), that is,
\(e_Ae_B\ge(w/k)^{1/b}\). By the lower
bound in \eqref{eq:event-orders} at \(E_1^-P_1^-\), by the bounds of Step~5
on \(x_A\) and \(x_B\) raised to the power \(1/a\), and by the bound
\(e_Ae_B\ge(w/k)^{1/b}\) at \(E_0^-P_0^-\) raised to the power \((1-b)/(ab)\), both powers being
positive,
\begin{align*}
 w&\overset{\eqref{eq:probability-CHSH}}{>}\omega(E_1^-P_1^-)
 \overset{\eqref{eq:event-orders}}{\ge}\lambda(x_Ax_B)^{1/a}
 \overset{\text{Step 5}}{\ge}\lambda(2k)^{-2/(ab)}(e_Ae_B)^{(1-b)/(ab)}\\
 &\overset{\eqref{eq:event-orders}}{\ge}
 \lambda(2k)^{-2/(ab)}\Bigl(\frac wk\Bigr)^{(1-b)/(ab^2)}
 =\lambda\,2^{-2/(ab)}\,k^{-(1+b)/(ab^2)}\,w^{(1-b)/(ab^2)},
\end{align*}
where the exponents of \(k\) add up to
\(-2/(ab)-(1-b)/(ab^2)=-(1+b)/(ab^2)\). Since \(w>0\), division by
\(w^{(1-b)/(ab^2)}\), with \(\eta=1-(1-b)/(ab^2)\) from
\eqref{eq:probability-exponents}, gives
\[
 w^\eta>\lambda\,2^{-2/(ab)}\,k^{-(1+b)/(ab^2)}
 =2^{-2/(ab)}K_{r_+,+}^{-\nu_+}K_{r_-,-}^{-\nu_-}.
\]
Here
\[
 k^{-(1+b)/(ab^2)}=K_{r_+,+}^{-(1+b)/(ar_+b^2)}=K_{r_+,+}^{-\nu_+}
 \quad\text{by \eqref{eq:nus}, and}\quad
 \lambda=K_{r_-,-}^{-1/(r_--1)}=K_{r_-,-}^{-\nu_-},
\]
because
\(r_--1=a/(1-a)\) by \eqref{eq:probability-exponents}. Multiplication by
\(K_{r_+,+}^{\nu_+}K_{r_-,-}^{\nu_-}\) gives
\begin{equation}
 w^\eta K_{r_+,+}^{\nu_+}K_{r_-,-}^{\nu_-}
 >2^{-2/(ab)}.
 \label{eq:necessary-cell}
\end{equation}

\smallskip
\emph{Step 7: an upper bound on \(w\) from clustering.} Let
\(\mathsf m_i=\omega(A_i)\) and \(\mathsf n_j=\omega(B_j)\) be the
marginal means, and let \(S_{\mathrm{prod}}\) be the score obtained by
replacing each correlator \(\omega(A_iB_j)\) by
\(\mathsf m_i\mathsf n_j\). Since \(|\mathsf m_i|\le1\) and
\(\frac12(|x+y|+|x-y|)=\max\{|x|,|y|\}\) for real \(x,y\),
\[
 S_{\mathrm{prod}}
 =\tfrac12\{\mathsf m_0(\mathsf n_0+\mathsf n_1)
             +\mathsf m_1(\mathsf n_0-\mathsf n_1)\}
 \le\tfrac12\{|\mathsf n_0+\mathsf n_1|+|\mathsf n_0-\mathsf n_1|\}
 =\max\{|\mathsf n_0|,|\mathsf n_1|\}.
\]
The correlation bound \eqref{eq:additive-cluster}, applied to the four
pairs \(A_i,B_j\), bounds each difference
\(\omega(A_iB_j)-\mathsf m_i\mathsf n_j\) by \(q_{\mathrm{cl}}\) in
absolute value, so
\(|S-S_{\mathrm{prod}}|\le\frac12\cdot4q_{\mathrm{cl}}=2q_{\mathrm{cl}}\),
and
\[
 1<S\le S_{\mathrm{prod}}+2q_{\mathrm{cl}}
 \le\max\{|\mathsf n_0|,|\mathsf n_1|\}+2q_{\mathrm{cl}}.
\]

The \emph{minority outcome} of \(B_j\) is the less probable of its two
effects \(P_j^\pm\); since \(\omega(P_j^\pm)=(1\pm\mathsf n_j)/2\), its
probability is \((1-|\mathsf n_j|)/2\). Thus \(S>1\) forces one of the
minority outcome probabilities of \(B_0,B_1\) to be smaller than \(q_{\mathrm{cl}}\):
for \(j\) with \(|\mathsf n_j|=\max\{|\mathsf n_0|,|\mathsf n_1|\}\), the
last display gives \((1-|\mathsf n_j|)/2<q_{\mathrm{cl}}\). A winning
cell lies in the minority outcome of \(B_j\): if the outcome
is \(P_j^\pm\), the cell \(\omega(E_0^\pm P_j^\pm)\) with the same sign is
winning, because \(\omega(A_0B_j)\) has the sign \(+\) in \(S\). Since
\(E_0^\pm P_j^\pm\le P_j^\pm\), the minimality of \(w\) gives
\[
 w\le\omega(E_0^\pm P_j^\pm)\le\omega(P_j^\pm)<q_{\mathrm{cl}}.
\]

\smallskip
\emph{Conclusion.} Since \(\eta>0\), \(0<w<q_{\mathrm{cl}}\) and
\(K_{r_+,+},K_{r_-,-}\ge1\),
\[
 w^\eta K_{r_+,+}^{\nu_+}K_{r_-,-}^{\nu_-}
 <q_{\mathrm{cl}}^\eta K_{r_+,+}^{\nu_+}K_{r_-,-}^{\nu_-}
 \overset{\eqref{eq:probability-cutoff}}{\le}2^{-2/(ab)},
\]
which contradicts \eqref{eq:necessary-cell}. Hence no reflection quartet
has \(S>1\), and \(\beta(\omega;\A,\B)=1\) by Step~1.
\end{proof}

\section{Gaussian comparison and source compression}
\label{sec:gaussian}

\subsection{Thermal powers and a covariance perturbation bound}

For quadratures with commutation matrix \(J_n\), a centred Gaussian
state with covariance \(V\) has characteristic function
\(\exp(-z^TVz/2)\). We use the Weyl and Gaussian formalism of
\cite{SLW18}, normalized so that the oscillator ground state has covariance \(I/2\).

For a faithful Gaussian state \(\rho_V\), let \(V^{[\tau]}\) be the
covariance of \(\rho_V^\tau/{\Tr\rho_V^\tau}\). In Williamson coordinates,
Proposition~1 of \cite{SLW18} gives, in our normalization,
\begin{equation}
 \upsilon_j^{[\tau]}
 =\frac12\coth\!\left(
      \tau\arctanh\frac1{2\upsilon_j}\right).
 \label{eq:thermal-power}
\end{equation}
For one mode with number operator \(N\), the thermal state
\(\rho_V\propto e^{-\beta N}\) has \(V=\upsilon I_2\) with
\(\upsilon=\Tr(\rho_VN)+\frac12=\frac12\coth(\beta/2)\), and
\(\rho_V^\tau\propto e^{-\tau\beta N}\) has
\(V^{[\tau]}=\frac12\coth(\tau\beta/2)I_2\), which is
\eqref{eq:thermal-power} because \(\beta/2=\arctanh(1/(2\upsilon))\).

For \(0<\tau<1\), the thermal increment is
\[
 \Therm_\tau(V)=V^{[\tau]}-V>0.
\]
In Williamson coordinates, \(\Therm_\tau(V)\) is diagonal with the
entries \(\upsilon_j^{[\tau]}-\upsilon_j\), which are positive because
\(\upsilon_j=\frac12\coth(\arctanh\frac1{2\upsilon_j})\) and \(\coth\)
decreases on \((0,\infty)\).

For faithful density operators and \(r>1\), the sandwiched Rényi
quantity is
\begin{equation}
 \Qtil_r(\rho\|\sigma)
 =\Tr\left(
 \sigma^{-(r-1)/(2r)}\rho
 \sigma^{-(r-1)/(2r)}\right)^r.
 \label{eq:renyi}
\end{equation}
Its logarithm divided by \(r-1\) is the sandwiched Rényi divergence,
introduced independently by Wilde, Winter and Yang \cite{WWY14} and by
Müller-Lennert, Dupuis, Szehr, Fehr and Tomamichel \cite{ML13}. The value in \eqref{eq:renyi} is \(+\infty\) when
the positive sandwich \(\sigma^{-(r-1)/(2r)}\rho\sigma^{-(r-1)/(2r)}\)
does not lie in \(\Sc_r\).
Theorem~\ref{thm:gaussian} bounds \(\Qtil_r\) with constants independent of
the number of oscillators and of their thermal parameters.

\begin{theorem}[Gaussian Rényi comparison]\label{thm:gaussian}
Let \(\rho,\sigma\) be centred faithful finite-mode Gaussian states
on the same canonical system, with covariance matrices
\(V_\rho,V_\sigma\). For \(r>1\), set
\[
 \tau=\frac{r-1}{r},\qquad k_r=\max\{2,\tau^{-1}\},\qquad
 \mathsf Z=
 \Therm_\tau(V_\sigma)^{-1/2}(V_\rho-V_\sigma)
 \Therm_\tau(V_\sigma)^{-1/2}.
\]
If \(\norm{\mathsf Z}\le(4k_r)^{-1}\), then
\begin{equation}
 \log\Qtil_r(\rho\|\sigma)\le2rk_r\norm{\mathsf Z}_1.
 \label{eq:gaussian-bound}
\end{equation}
\end{theorem}

\begin{proof}
\emph{Step 1: independence of the quadrature coordinates.} For a real
invertible matrix \(S\), a positive matrix \(A\) and a real symmetric
matrix \(X\) of the same size, the matrix
\(W=(SAS^T)^{-1/2}SA^{1/2}\) satisfies \(WW^T=I\), so it is
orthogonal, and
\[
 (SAS^T)^{-1/2}\,SXS^T\,(SAS^T)^{-1/2}
 =W\bigl(A^{-1/2}XA^{-1/2}\bigr)W^T.
\]
A real change of quadrature coordinates by \(S\), with the commutation
form transformed at the same time, replaces \(V_\rho\), \(V_\sigma\) and
\(\Therm_\tau(V_\sigma)\) by \(SV_\rho S^T\), \(SV_\sigma S^T\) and
\(S\Therm_\tau(V_\sigma)S^T\).
With \(A=\Therm_\tau(V_\sigma)\) and \(X=V_\rho-V_\sigma\), the identity
above shows that \(\mathsf Z\) changes to \(W\mathsf ZW^T\). Hence
\(\norm{\mathsf Z}\), \(\norm{\mathsf Z}_1\) and the norm hypothesis do
not depend on the quadrature coordinates; neither does
\(\Qtil_r(\rho\|\sigma)\), which is defined by the two states.

\smallskip
\emph{Step 2: the sandwich formulas.} In Williamson coordinates for
\(\sigma\), which Step~1 allows, we write
\[
 \begin{gathered}
 V_0=V_\sigma,\qquad
 \mathsf T_\tau=V_0^{[\tau]},\qquad \mathsf H=\mathsf T_\tau-V_0,\\
 \mathsf T_{1-\tau}=V_0^{[1-\tau]},\qquad
 \mathsf L=(\mathsf T_\tau^2-I/4)^{1/2}.
 \end{gathered}
\]
For the sandwich \(\Theta=\sigma^{-\tau/2}\rho\sigma^{-\tau/2}\) and the
covariance \(V_\tau\) of the normalized sandwich \(\Theta/{\Tr\Theta}\),
Appendix~\ref{app:gaussian} gives the Gaussian sandwich formulas
\begin{equation}
 \Tr\Theta
 =\frac{\Tr\sigma^{1-\tau}}{\sqrt{\det(I-\mathsf Z)}},
 \qquad
 V_\tau=-\mathsf T_\tau+
       \mathsf L(\mathsf T_\tau-V_\rho)^{-1}\mathsf L,
 \label{eq:sandwich}
\end{equation}
whenever \(V_\rho<\mathsf T_\tau\); under
this condition Seshadreesan, Lami and Wilde give \(\Qtil_r(\rho\|\sigma)\)
in closed form \cite[Theorem~22]{SLW18}.

The norm hypothesis implies the condition \(V_\rho<\mathsf T_\tau\): since \(\norm{\mathsf Z}\le(4k_r)^{-1}<1\),
\[
 \mathsf T_\tau-V_\rho
 =\mathsf H-(V_\rho-V_0)
 =\mathsf H^{1/2}(I-\mathsf Z)\mathsf H^{1/2}>0.
\]
Then, by Appendix~\ref{app:gaussian}, \(\Theta\) is a positive
trace-class operator with the trace \eqref{eq:traceTheta}.

\smallskip
\emph{Step 3: the perturbation of the covariance.} At
\(V_\rho=V_0\), the sandwich is \(\sigma^{1-\tau}\), and the normalized
state \(\sigma^{1-\tau}/{\Tr\sigma^{1-\tau}}\) has covariance
\(\mathsf T_{1-\tau}\); so \eqref{eq:sandwich} at \(V_\rho=V_0\) reads
\(\mathsf T_{1-\tau}=-\mathsf T_\tau+\mathsf L\mathsf H^{-1}\mathsf L\).
For general \(V_\rho\), Step~2 gives
\((\mathsf T_\tau-V_\rho)^{-1}
=\mathsf H^{-1/2}(I-\mathsf Z)^{-1}\mathsf H^{-1/2}\). Subtracting the
formula at \(V_\rho=V_0\) from \eqref{eq:sandwich} gives
\[
 V_\tau-\mathsf T_{1-\tau}
 \overset{\eqref{eq:sandwich}}{=}
 \mathsf L\bigl((\mathsf T_\tau-V_\rho)^{-1}-\mathsf H^{-1}\bigr)\mathsf L
 \overset{\text{Step 2}}{=}
 \mathsf L\mathsf H^{-1/2}\bigl((I-\mathsf Z)^{-1}-I\bigr)
 \mathsf H^{-1/2}\mathsf L.
\]
In Williamson coordinates \(V_0\) is diagonal, and so are
\(\mathsf T_\tau\) (by \eqref{eq:thermal-power}), \(\mathsf H\) and
\(\mathsf L\). Hence \(\mathsf P=\mathsf L\mathsf H^{-1/2}\) is a diagonal
matrix equal to \(\mathsf H^{-1/2}\mathsf L\), so
\(\mathsf P^2=\mathsf L\mathsf H^{-1}\mathsf L\), and the formula at
\(V_\rho=V_0\) and the last display give
\[
 \mathsf P^2=\mathsf T_{1-\tau}+\mathsf T_\tau,\qquad
 V_\tau-\mathsf T_{1-\tau}
 =\mathsf P\bigl((I-\mathsf Z)^{-1}-I\bigr)\mathsf P.
\]

For each thermal mode, with \(u=\arctanh(1/(2\upsilon_j))\), formula
\eqref{eq:thermal-power} gives the entries \(\frac12\coth(\tau u)\) of
\(\mathsf T_\tau\) and \(\frac12\coth((1-\tau)u)\) of
\(\mathsf T_{1-\tau}\), the two terms of \(\mathsf P^2\). If \(\tau\le1/2\), then \(\tau u\le(1-\tau)u\),
and since \(x\coth x\) increases on \((0,\infty)\),
\(\tau u\coth(\tau u)\le(1-\tau)u\coth((1-\tau)u)\). If
\(\tau\ge1/2\), then \(\tau u\ge(1-\tau)u\), and since \(\coth\)
decreases, \(\coth(\tau u)\le\coth((1-\tau)u)\). The two cases give
\[
 \frac{\mathsf T_\tau}{\mathsf T_{1-\tau}}
 =\frac{\coth(\tau u)}{\coth((1-\tau)u)}
 \le\max\{1,(1-\tau)/\tau\},
\]
and hence, since \(1+\max\{1,(1-\tau)/\tau\}=\max\{2,\tau^{-1}\}=k_r\),
\(\mathsf T_{1-\tau}+\mathsf T_\tau\le k_r\mathsf T_{1-\tau}\).

For
\[
 \delta=\norm{\mathsf Z},\qquad
 \mathsf Y=\mathsf T_{1-\tau}^{-1/2}(V_\tau-\mathsf T_{1-\tau})\mathsf T_{1-\tau}^{-1/2},
\]
the identity \((I-\mathsf Z)^{-1}-I=\mathsf Z(I-\mathsf Z)^{-1}\) gives
the first formula below, and the bound
\(\mathsf T_{1-\tau}+\mathsf T_\tau\le k_r\mathsf T_{1-\tau}\) gives the second, since
\(\mathsf P\) and \(\mathsf T_{1-\tau}\) are diagonal:
\begin{align*}
 \mathsf Y
 &=\mathsf T_{1-\tau}^{-1/2}\mathsf P\,\mathsf Z(I-\mathsf Z)^{-1}\,
   \mathsf P\mathsf T_{1-\tau}^{-1/2},\\
 \norm{\mathsf T_{1-\tau}^{-1/2}\mathsf P}^2
 &=\norm{\mathsf T_{1-\tau}^{-1}(\mathsf T_{1-\tau}+\mathsf T_\tau)}
 \le k_r.
\end{align*}
The eigenvalues of \(\mathsf Z\) lie in \([-\delta,\delta]\), so
\(\norm{\mathsf Z(I-\mathsf Z)^{-1}}\le\delta/(1-\delta)\) and
\(\norm{\mathsf Z(I-\mathsf Z)^{-1}}_1\le\norm{\mathsf Z}_1/(1-\delta)\).
We obtain
\begin{equation}
 \norm{\mathsf Y}\le\frac{k_r\delta}{1-\delta},\qquad
 \norm{\mathsf Y}_1\le\frac{k_r}{1-\delta}\norm{\mathsf Z}_1.
 \label{eq:purity-perturbation}
\end{equation}
The norm hypothesis gives \(k_r\delta\le1/4\), hence \(\delta\le1/8\)
since \(k_r\ge2\), and \eqref{eq:purity-perturbation} gives
\(\norm{\mathsf Y}\le\frac14\cdot\frac87=\frac27\).

\smallskip
\emph{Step 4: the purity function and its derivative.} For a physical
covariance, define
\[
 \mathscr P_r(V)=\log\Tr\rho_V^r
 =-\sum_j\log f_r(\upsilon_j),\qquad
 f_r(x)=(x+\tfrac12)^r-(x-\tfrac12)^r.
\]
The second expression holds mode by mode: a thermal mode with
symplectic eigenvalue \(\upsilon\) is the state \((1-\zeta)\zeta^N\)
with \(\zeta=(\upsilon-1/2)/(\upsilon+1/2)\), and
\(\Tr((1-\zeta)\zeta^N)^r=(1-\zeta)^r/(1-\zeta^r)=1/f_r(\upsilon)\).
For \(x\ge1/2\), the substitution
\(\zeta=(x-1/2)/(x+1/2)\) gives \(f_r(x)=(x+\frac12)^r(1-\zeta^r)\),
\(f_r'(x)=r(x+\frac12)^{r-1}(1-\zeta^{r-1})\) and
\(x=(x+\frac12)(1+\zeta)/2\), hence
\[
 0\le\frac{xf_r'(x)}{f_r(x)}
 =\frac{r(1+\zeta)}2
   \frac{1-\zeta^{r-1}}{1-\zeta^r}
 \le r,
\]
where the upper bound is equivalent to \((1-\zeta)(1+\zeta^{r-1})\ge0\),
since \(2(1-\zeta^r)-(1+\zeta)(1-\zeta^{r-1})=(1-\zeta)(1+\zeta^{r-1})\).

The derivative of \(\mathscr P_r\) involves the derivatives of the
symplectic eigenvalues of \(V\), which are the positive eigenvalues of
\(iJ_nV\). At a Williamson-diagonal covariance
\(\Upsilon=\bigoplus_j\upsilon_jI_2\) with distinct \(\upsilon_j\),
the vector \(\psi_j=(e_{2j-1}-ie_{2j})/\sqrt2\), where the \(e_k\) are
the standard basis vectors, satisfies \(iJ_n\psi_j=\psi_j\). Since
\(iJ_n\) is self-adjoint, \(\psi_j\) is both a right and a left
eigenvector of \(iJ_n\Upsilon\) for its simple eigenvalue
\(\upsilon_j\), and first-order perturbation of this eigenvalue shows
that the differential of a symplectic eigenvalue is half the trace of
the corresponding two-dimensional diagonal block \(\dot V_{jj}\):
\[
 d\upsilon_j[\dot V]
 =\psi_j^*\,iJ_n\dot V\,\psi_j
 =\psi_j^*\dot V\psi_j
 =\tfrac12\Tr\dot V_{jj}.
\]
With \(0\le f_r'(\upsilon_j)/f_r(\upsilon_j)\le r/\upsilon_j\), the bound on
\(xf_r'(x)/f_r(x)\) at \(x=\upsilon_j\), and
since the diagonal blocks of a matrix have total trace norm at most
its trace norm, it follows that
\begin{align*}
 |d\mathscr P_r(\Upsilon)[\dot V]|
 &\le\sum_j\frac{f_r'(\upsilon_j)}{f_r(\upsilon_j)}\,
   \frac{\lvert\Tr\dot V_{jj}\rvert}2
 \le\frac r2\sum_j\bigl|\Tr(\upsilon_j^{-1}\dot V_{jj})\bigr|\\
 &\le\frac r2\norm{\Upsilon^{-1/2}\dot V\Upsilon^{-1/2}}_1.
\end{align*}

By Williamson's normal form, a physical covariance is
\(V=S\Upsilon S^T\) with \(S\) symplectic. Under this symplectic
congruence \(\mathscr P_r\) is unchanged, so
\(d\mathscr P_r(V)[\dot V]=d\mathscr P_r(\Upsilon)[S^{-1}\dot VS^{-T}]\),
and by Step~1, with \(A=\Upsilon\) and \(X=S^{-1}\dot VS^{-T}\), the
relative perturbation \(V^{-1/2}\dot VV^{-1/2}\) is orthogonally
similar to \(\Upsilon^{-1/2}S^{-1}\dot VS^{-T}\Upsilon^{-1/2}\), so the two
matrices have the same trace norm. Therefore, when the symplectic eigenvalues of
\(V\) are distinct,
\begin{equation}
 |d\mathscr P_r(V)[\dot V]|
 \le\frac r2\norm{V^{-1/2}\dot V V^{-1/2}}_1.
 \label{eq:purity-derivative}
\end{equation}
A covariance \(V\) is
faithful when \(V+\frac i2J_n>0\), which by Williamson's normal form, as
after \eqref{eq:strict-uncertainty}, holds exactly when
\(\upsilon_j>1/2\) for all \(j\). The Hermitian matrix
\(V^{1/2}iJ_nV^{1/2}=V^{1/2}(iJ_nV)V^{-1/2}\) has the eigenvalues of
\(iJ_nV\), which by Williamson's normal form are
\(\pm\upsilon_1,\dots,\pm\upsilon_n\); so
\[
 \mathscr P_r(V)
 =-\frac12\Tr F\bigl(V^{1/2}iJ_nV^{1/2}\bigr),\qquad
 F(x)=\log f_r(|x|).
\]
For faithful \(V\), the eigenvalues \(\pm\upsilon_j\) lie in the open set \(\{|x|>1/2\}\),
on which \(F\) is real-analytic, since \(f_r\) is real-analytic and
positive on \((1/2,\infty)\). By the holomorphic functional calculus,
\(M\mapsto\Tr F(M)\) is smooth on the Hermitian matrices with spectrum
in \(\{|x|>1/2\}\), also where eigenvalues coincide, and \(V\mapsto V^{1/2}\)
is smooth on positive definite matrices. Hence \(\mathscr P_r\) is
smooth on the open set of faithful covariances, and
\(d\mathscr P_r(V)\) is continuous in \(V\) there. Covariances with
distinct symplectic eigenvalues are dense in the faithful covariances,
and the right
side of \eqref{eq:purity-derivative} is continuous in \(V\); so
\eqref{eq:purity-derivative} holds at every faithful covariance.

\smallskip
\emph{Step 5: integration along a segment.} The segment
\[
 V(v)=\mathsf T_{1-\tau}+v(V_\tau-\mathsf T_{1-\tau})
 =\mathsf T_{1-\tau}^{1/2}(I+v\mathsf Y)\mathsf T_{1-\tau}^{1/2},
 \qquad 0\le v\le1,
\]
consists of physical covariances, as convex combinations of the
covariances \(\mathsf T_{1-\tau}\) and \(V_\tau\) of two states. For
\(v<1\) the covariances \(V(v)\) are faithful: \(\mathsf T_{1-\tau}\) is
faithful, since its symplectic eigenvalues \(\frac12\coth((1-\tau)u)\)
from Step~3 exceed \(\frac12\), and \(V_\tau+\frac i2J_n\ge0\), so
\[
 V(v)+\tfrac i2J_n
 =(1-v)\bigl(\mathsf T_{1-\tau}+\tfrac i2J_n\bigr)
 +v\bigl(V_\tau+\tfrac i2J_n\bigr)>0,
 \qquad 0\le v<1.
\]

The segment has the derivative
\(\dot V=\mathsf T_{1-\tau}^{1/2}\mathsf Y\mathsf T_{1-\tau}^{1/2}\).
By Step~1, with \(S=\mathsf T_{1-\tau}^{1/2}\), \(A=I+v\mathsf Y\) and
\(X=\mathsf Y\), the matrix \(V(v)^{-1/2}\dot VV(v)^{-1/2}\) is
orthogonally similar to
\((I+v\mathsf Y)^{-1/2}\mathsf Y(I+v\mathsf Y)^{-1/2}
=\mathsf Y(I+v\mathsf Y)^{-1}\), whose trace norm is at most
\(\norm{\mathsf Y}_1/(1-\norm{\mathsf Y})\). On physical covariances,
\(\mathscr P_r=-\sum_j\log f_r(\upsilon_j)\) is
continuous, because \(f_r\) is continuous and positive on
\([1/2,\infty)\) and the symplectic eigenvalues depend continuously on \(V\); so
integrating
\eqref{eq:purity-derivative} from \(0\) to \(v<1\), letting \(v\to1\), and
using \eqref{eq:purity-perturbation} gives
\[
 |\mathscr P_r(V_\tau)-\mathscr P_r(\mathsf T_{1-\tau})|
 \le\frac r2\,\frac{\norm{\mathsf Y}_1}{1-\norm{\mathsf Y}}
 \le
 \frac r2
 \frac{k_r}{(1-\delta)(1-\norm{\mathsf Y})}
 \norm{\mathsf Z}_1.
\]

\smallskip
\emph{Step 6: assembly and the constant.} Since \((r-1)/(2r)=\tau/2\),
and since the normalized sandwich \(\Theta/{\Tr\Theta}\), a Gaussian
state, has covariance \(V_\tau\),
\begin{align*}
 \Qtil_r(\rho\|\sigma)
 &\overset{\eqref{eq:renyi}}{=}\Tr\Theta^r
 =(\Tr\Theta)^r\,\Tr\biggl(\frac{\Theta}{\Tr\Theta}\biggr)^r\\
 &\overset{\eqref{eq:sandwich}}{=}
 \frac{(\Tr\sigma^{1-\tau})^r}{\det(I-\mathsf Z)^{r/2}}\,
 e^{\mathscr P_r(V_\tau)}.
\end{align*}
Because \(r(1-\tau)=1\), the factor \((\Tr\sigma^{1-\tau})^r\) equals
\(e^{-\mathscr P_r(\mathsf T_{1-\tau})}\): the state
\(\sigma^{1/r}/{\Tr\sigma^{1/r}}\) has covariance \(\mathsf T_{1-\tau}\),
so
\[
 (\Tr\sigma^{1/r})^re^{\mathscr P_r(\mathsf T_{1-\tau})}
 =(\Tr\sigma^{1/r})^r
 \Tr\biggl(\frac{\sigma^{1/r}}{\Tr\sigma^{1/r}}\biggr)^r=\Tr\sigma=1.
\]
Taking logarithms, with \(\log\det=\Tr\log\), gives
\[
 \log\Qtil_r(\rho\|\sigma)
 =-\frac r2\Tr\log(I-\mathsf Z)
  +\mathscr P_r(V_\tau)-\mathscr P_r(\mathsf T_{1-\tau}).
\]

The bound \(\lvert\log(1-t)\rvert\le|t|/(1-|t|)\), \(|t|<1\), applied to the
eigenvalues of \(\mathsf Z\), gives
\[
 \lvert\Tr\log(I-\mathsf Z)\rvert
 \le\frac{\norm{\mathsf Z}_1}{1-\delta}.
\]
With Step~5, and with \(1/(1-\delta)\le8/7\) and
\(1/(1-\norm{\mathsf Y})\le7/5\) from \(\delta\le1/8\) and
\(\norm{\mathsf Y}\le2/7\),
\begin{align*}
 \log\Qtil_r(\rho\|\sigma)
 &\le\frac r2\,\frac{\norm{\mathsf Z}_1}{1-\delta}
 +\frac r2\,\frac{k_r\norm{\mathsf Z}_1}{(1-\delta)(1-\norm{\mathsf Y})}\\
 &\le r\Bigl(\frac47+\frac45k_r\Bigr)\norm{\mathsf Z}_1
 \overset{k_r\ge2}{\le}2rk_r\norm{\mathsf Z}_1,
\end{align*}
proving \eqref{eq:gaussian-bound}.
\end{proof}

The effect estimate below is the two-outcome consequence of data
processing for the sandwiched Rényi divergence, proved for all \(r>1\)
independently by Frank and Lieb \cite{FL13} and by Beigi
\cite{Beigi13}.

\begin{lemma}[Effects and trace distance]\label{lem:renyi-consequences}
Let \(\rho,\sigma\) be faithful density operators and \(r>1\), and
write \(\rho(M)=\Tr(\rho M)\) and \(\sigma(M)=\Tr(\sigma M)\). If
\(\Qtil_r(\rho\|\sigma)<\infty\), every effect \(M\)
satisfies
\begin{equation}
 \rho(M)^r\le
 \Qtil_r(\rho\|\sigma)\,\sigma(M)^{r-1}.
 \label{eq:effect-renyi}
\end{equation}
For \(r\ge2\), one also has
\begin{equation}
 \norm{\rho-\sigma}_1
 \le\sqrt{\Qtil_r(\rho\|\sigma)^{1/(r-1)}-1}.
 \label{eq:trace-distance}
\end{equation}
\end{lemma}

\begin{proof}
\emph{Step 1: the effect bound.} With \(\tau=(r-1)/r\), finiteness of
\(\Qtil_r(\rho\|\sigma)\) means that the sandwich
\(\Theta=\sigma^{-\tau/2}\rho\sigma^{-\tau/2}\) of \eqref{eq:sandwich}
lies in \(\Sc_r\), with
\(\rho=\sigma^{\tau/2}\Theta\sigma^{\tau/2}\), as \(\sigma^{\tau/2}\) maps every
vector into the domain of \(\sigma^{-\tau/2}\). Since
\(\sigma^{\tau/2}\in\Sc_{2/\tau}\) and \(\tau+1/r=1\), the trace is
cyclic on the products below, and the Schatten Hölder inequality with
the exponents \(r\) and \(1/\tau\) bounds \(\rho(M)\) by the product of
\(\norm{\Theta}_r\), which is \(\Qtil_r(\rho\|\sigma)^{1/r}\) by
\eqref{eq:renyi}, and the factor
\(\norm{\sigma^{\tau/2}M\sigma^{\tau/2}}_{1/\tau}\):
\[
 \rho(M)
 =\Tr(\sigma^{\tau/2}\Theta\sigma^{\tau/2}M)
 =\Tr(\Theta\,\sigma^{\tau/2}M\sigma^{\tau/2})
 \overset{\text{Hölder}}{\le}\norm{\Theta}_r\,
   \norm{\sigma^{\tau/2}M\sigma^{\tau/2}}_{1/\tau}.
\]

The operator \(\sigma^{\tau/2}M\sigma^{\tau/2}\) has the same nonzero
spectrum as \(M^{1/2}\sigma^\tau M^{1/2}\). Hansen's inequality for operator monotone
functions and contractions \cite[Theorem]{Hansen80}, applied to
\(t\mapsto t^\tau\) and the contraction \(M^{1/2}\), gives
\[
 M^{1/2}\sigma^\tau M^{1/2}
 \le(M^{1/2}\sigma M^{1/2})^\tau.
\]
Monotonicity of the trace of a positive power bounds the remaining
factor by \(\sigma(M)^\tau\):
\begin{align*}
 \norm{\sigma^{\tau/2}M\sigma^{\tau/2}}_{1/\tau}
 &=\norm{M^{1/2}\sigma^\tau M^{1/2}}_{1/\tau}
 \overset{\text{Hansen}}{\le}\norm{(M^{1/2}\sigma M^{1/2})^\tau}_{1/\tau}\\
 &=\bigl(\Tr M^{1/2}\sigma M^{1/2}\bigr)^\tau
 =\sigma(M)^\tau.
\end{align*}
Hence \(\rho(M)\le\Qtil_r(\rho\|\sigma)^{1/r}\sigma(M)^\tau\), and
raising to the power \(r\), with \(r\tau=r-1\), proves
\eqref{eq:effect-renyi}.

\smallskip
\emph{Step 2: comparison of the orders \(2\) and \(r\).}
For \(r\ge2\), Schatten interpolation \cite{Simon05} of
\[
 z\longmapsto\sigma^{-\tau z/2}\rho\sigma^{-\tau z/2}
\]
between the trace norm and the \(r\)-norm, at \(z=1/(2\tau)\), gives
\[
 \Qtil_2(\rho\|\sigma)
 \le\Qtil_r(\rho\|\sigma)^{1/(r-1)}.
\]
Finite spectral compressions of \(\sigma\) justify the interpolation for
unbounded inverse powers. Let \(\Pi_n\) be the spectral projection of
\(\sigma\) for \([1/n,1]\), which has finite rank, and let
\(\Theta_{1/2}=\sigma^{-1/4}\rho\sigma^{-1/4}\), the sandwich of
\eqref{eq:sandwich} for \(r=2\). The compressed family
\(z\mapsto\Pi_n\sigma^{-\tau z/2}\rho\sigma^{-\tau z/2}\Pi_n\) is entire
and bounded on the strip \(0\le\Rea z\le1\). Since
\(\sigma^{-i\tau t/2}\Pi_n\), \(t\in\R\), is a partial isometry
commuting with \(\Pi_n\), the trace-norm endpoint is at most \(1\), and
the other endpoint is a compression of the full \(r\)-sandwich:
\begin{align*}
 \norm{\Pi_n\sigma^{-i\tau t/2}\rho\sigma^{-i\tau t/2}\Pi_n}_1
 &\le\norm{\rho}_1=1,\\
 \norm{\Pi_n\sigma^{-\tau(1+it)/2}\rho\sigma^{-\tau(1+it)/2}\Pi_n}_r
 &=\norm{\Pi_n\Theta\Pi_n}_r\le\norm{\Theta}_r.
\end{align*}

The point \(\theta=1/(2\tau)\) lies in \((0,1]\) because \(r\ge2\), and
the interpolated exponent \(p_\theta\) satisfies
\(1/p_\theta=(1-\theta)+\theta/r=1-\theta\tau=1/2\). The three-lines
bound at \(z=\theta\) therefore gives
\[
 \norm{\Pi_n\Theta_{1/2}\Pi_n}_2
 \le\norm{\Theta}_r^{\theta}
 =\Qtil_r(\rho\|\sigma)^{1/(2(r-1))}.
\]
As \(n\to\infty\), the left-hand side increases to
\(\lVert\Theta_{1/2}\rVert_2=\Qtil_2(\rho\|\sigma)^{1/2}\), since
\(\lVert\Pi_n\Theta_{1/2}\Pi_n\rVert_2^2\) is the sum of the squared moduli
of the matrix entries of \(\Theta_{1/2}\) between eigenvectors of
\(\sigma\) with eigenvalues at least \(1/n\).
The Hilbert--Schmidt limit gives
\(\Qtil_2(\rho\|\sigma)\le\Qtil_r(\rho\|\sigma)^{1/(r-1)}\). For density
matrices, this inequality is the monotonicity in the order of the
sandwiched Rényi divergence, proved independently by Beigi
\cite[Theorem~7]{Beigi13} and by Müller-Lennert et al.\
\cite[Theorem~7]{ML13}, and the argument above is Beigi's proof, applied
to the orders \(2\) and \(r\).

\smallskip
\emph{Step 3: the trace distance.} Since
\(\Tr(\Theta_{1/2}\sigma^{1/2})
=\Tr(\sigma^{1/4}\Theta_{1/2}\sigma^{1/4})=\Tr\rho=1\),
\[
 \norm{\Theta_{1/2}-\sigma^{1/2}}_2^2
 =\Tr\Theta_{1/2}^2-2\Tr(\Theta_{1/2}\sigma^{1/2})+\Tr\sigma
 =\Qtil_2(\rho\|\sigma)-1.
\]
The compressions \(\Pi_n\) of Step~2 justify this identity. Finally,
\[
 \rho-\sigma
 =\sigma^{1/4}(\Theta_{1/2}-\sigma^{1/2})\sigma^{1/4},
\]
and Hölder with \(\norm{\sigma^{1/4}}_4=1\) gives
\begin{align*}
 \norm{\rho-\sigma}_1
 &\overset{\phantom{\text{Step 2}}}{\le}\norm{\sigma^{1/4}}_4\,\norm{\Theta_{1/2}-\sigma^{1/2}}_2\,
   \norm{\sigma^{1/4}}_4
 =\sqrt{\Qtil_2(\rho\|\sigma)-1}\\
 &\overset{\text{Step 2}}{\le}\sqrt{\Qtil_r(\rho\|\sigma)^{1/(r-1)}-1},
\end{align*}
which is \eqref{eq:trace-distance}.
\end{proof}

\subsection{Thermal weights and compression}

A full source space is a space \(\mathcal X_{\mu,U}\) of
Section~\ref{subsec:sources}, and a finite source space is a finite
canonical source system; both carry the covariance inner product. On a
source space of either kind, the symplectic
form is represented by the bounded skew-adjoint operator \(\mathsf J\)
with
\[
 \vv_{\mu,U}(\sx,\mathsf J\sy)=\varsigma(\sx,\sy),
 \qquad \norm{\mathsf J}\le2,
\]
the norm bound being the imaginary part of \eqref{eq:gram}. In coordinates
of a finite source space that are orthonormal for the covariance,
\(\mathsf J\) is the commutation matrix. On the complexification, with
\(G,\overline G,\Res\) the operators \(G_U,\overline G_U,\Res_U\) of
\eqref{eq:reserve} written without the subscript, \eqref{eq:gram} gives
\(\ip{\sx}{G\sy}=\vv_{\mu,U}(\sx,\sy)+\frac i2\vv_{\mu,U}(\sx,\mathsf J\sy)\)
for real \(\sx,\sy\), hence
\[
 G=I+i\mathsf J/2,\qquad
 \overline G=I-i\mathsf J/2,\qquad
 d=|i\mathsf J/2|,\qquad
 \Res=(G\overline G)^{1/2}=(I-d^2)^{1/2},
\]
where the last identity holds because \(i\mathsf J/2\) is self-adjoint,
and where \(d\) and \(\Res\) act on the complexification and preserve
the real source space.

In finite canonical dimension, \(0<d<I\): \(d\) is invertible because \(\varsigma\) is
nondegenerate, and \(d<I\) because \(G\) and \(\overline G\) are positive
definite by \eqref{eq:strict-uncertainty}. An
orthogonal change of basis then brings \(\mathsf J\) to blocks
\(2d_jJ_1\), on which \(d=d_jI_2\) with \(0<d_j<1\), and keeps the
covariance \(I\). Rescaling the two basis vectors of a block by
\((2d_j)^{-1/2}\) gives the commutation matrix \(J_1\) and the
covariance \(I_2/(2d_j)\), so the block has symplectic eigenvalue
\(\upsilon_j=1/(2d_j)\), and \eqref{eq:thermal-power} gives
\(\upsilon_j^{[\tau]}=\frac12\coth(\tau\arctanh d_j)\). Undoing the
rescaling multiplies covariances by \(2d_j\), so formula
\eqref{eq:thermal-power} becomes
\[
 \Therm_\tau(I)=d\coth(\tau\arctanh d)-I.
\]

For \(0<d<1\) the scalar function \(d\coth(\tau\arctanh d)-1\) is positive.
The function extends continuously at \(d=0\), with the value \((1-\tau)/\tau\),
and scalar functional
calculus defines the corresponding bounded positive operator on the
full source space, with limiting value zero at \(d=1\).

\begin{lemma}[Thermal comparison]\label{lem:thermal-comparison}
For \(0<\tau<1\),
\begin{equation}
 \Therm_\tau(I)\ge c_\tau\Res^{2\tau},\qquad
 c_\tau=
 \begin{cases}
  2^{1-2\tau},&0<\tau\le1/2,\\
  (1-\tau)/\tau,&1/2\le\tau<1.
 \end{cases}
 \label{eq:scalar-reserve}
\end{equation}
These constants are the infima over \(0<d<1\) of the scalar ratios
\((d\coth(\tau\arctanh d)-1)(1-d^2)^{-\tau}\).
\end{lemma}

\begin{proof}
\emph{Step 1: the ratio as a function of \(u\).} For \(d=\tanh u\),
\(u>0\), one has \((1-d^2)^{-\tau}=\cosh(u)^{2\tau}\) and
\[
 \tanh u\coth(\tau u)-1
 =\frac{\sinh u\cosh(\tau u)-\cosh u\sinh(\tau u)}
       {\cosh u\sinh(\tau u)}
 =\frac{\sinh((1-\tau)u)}{\cosh u\sinh(\tau u)},
\]
so the ratio in \eqref{eq:scalar-reserve} is
\[
 \mathfrak q_\tau(u)
 =\frac{\sinh((1-\tau)u)}{\sinh(\tau u)}
   \cosh(u)^{2\tau-1}.
\]
Its limits at zero and infinity are \((1-\tau)/\tau\) and
\(2^{1-2\tau}\); at infinity, \(\sinh((1-\tau)u)/{\sinh(\tau u)}\)
behaves as \(e^{(1-2\tau)u}\) and \(\cosh(u)^{2\tau-1}\) as
\(2^{1-2\tau}e^{(2\tau-1)u}\).

\smallskip
\emph{Step 2: an auxiliary convex function.} For fixed \(u>0\), define
\(\mathfrak f(v)=v\coth(vu)\). Its derivative
\[
 \mathfrak f'(v)
 =\coth(vu)-vu\operatorname{csch}^2(vu)
\]
is increasing, because
\[
 \mathfrak f''(v)
 =2u\operatorname{csch}^2(vu)(vu\coth(vu)-1)>0.
\]
For \(0<v\le1\), one also has
\(\mathfrak f'(v)\le\mathfrak f'(1)\le\tanh u\);
the last inequality is equivalent to \(\tanh u\le u\), because
\[
 \tanh u-\mathfrak f'(1)
 =\frac u{\sinh^2u}-(\coth u-\tanh u)
 =\frac u{\sinh^2u}-\frac1{\sinh u\cosh u}
 =\frac{u-\tanh u}{\sinh^2u}.
\]

\smallskip
\emph{Step 3: monotonicity of the ratio.} By the mean value theorem,
\(\mathfrak f(1-\tau)-\mathfrak f(\tau)=(1-2\tau)\mathfrak f'(\vartheta)\)
for some \(\vartheta\) between \(\tau\) and \(1-\tau\), so
\(0<\vartheta<1\). Consequently,
\begin{align*}
 \frac d{du}\log\mathfrak q_\tau(u)
 &=(1-\tau)\coth((1-\tau)u)-\tau\coth(\tau u)+(2\tau-1)\tanh u\\
 &=\mathfrak f(1-\tau)-\mathfrak f(\tau)
  +(2\tau-1)\tanh u
 =(1-2\tau)\bigl(\mathfrak f'(\vartheta)-\tanh u\bigr)
\end{align*}
is nonpositive for \(\tau\le1/2\) and nonnegative for
\(\tau\ge1/2\), since \(\mathfrak f'(\vartheta)\le\tanh u\) by Step~2.
So \(\mathfrak q_\tau\) is nonincreasing for \(\tau\le1/2\), with
infimum its limit \(2^{1-2\tau}\) at infinity, and nondecreasing for
\(\tau\ge1/2\), with infimum its limit \((1-\tau)/\tau\) at zero. In
both cases the infimum is \(c_\tau\), so
\(d\coth(\tau\arctanh d)-1\ge c_\tau(1-d^2)^\tau\) for \(0<d<1\), by
continuity also for \(0\le d\le1\), and the scalar functional calculus
of \(d\) gives \eqref{eq:scalar-reserve}.
\end{proof}

For positive operators, the weighted geometric mean is
\[
 A\#_\tau B
 =A^{1/2}(A^{-1/2}BA^{-1/2})^\tau A^{1/2},
 \qquad 0<\tau<1.
\]
We write \(A\#B=A\#_{1/2}B\). Noninvertible cases are defined by
decreasing positive regularization. Parallel sums
\((\sum_jA_j^{-1})^{-1}\) have the analogous quadratic-form
interpretation. Weighted geometric means and parallel sums are
Kubo--Ando operations, and we use their compression properties
\cite[Theorems~3.4 and~3.5]{KA80}.

\begin{lemma}[Compression of thermal weights]\label{lem:compression}
Let \(\iota\) be a real isometry from a finite source space into a
full source space, with \(\varsigma(\iota\sx,\iota\sy)=\varsigma(\sx,\sy)\) for all
\(\sx,\sy\), as for an inclusion, and write
\[
 G_F=\iota^*G\iota,\qquad
 \overline G_F=\iota^*\overline G\iota,\qquad
 \Res_F=(G_F\overline G_F)^{1/2}.
\]
For \(0<a\le1/2\),
\begin{equation}
 \Res_F^{2a}\ge\iota^*\Res^{2a}\iota.
 \label{eq:low-compression}
\end{equation}
For an integer \(r\ge2\) and \(b=(r-1)/r\),
\begin{equation}
 \Therm_b(I_F)\ge\iota^*\Therm_b(I)\iota
 \ge c_b\,\iota^*\Res^{2b}\iota.
 \label{eq:high-compression}
\end{equation}
\end{lemma}

\begin{proof}
\emph{Step 1: compression of parallel sums and geometric means.} For
positive invertible \(A,B\), the variational formula
\[
 \ip{z}{(A^{-1}+B^{-1})^{-1}z}
 =\inf_{x+y=z}\{\ip{x}{Ax}+\ip{y}{By}\}
\]
gives monotonicity and the compression inequality for parallel sums:
restricting the vectors in the infimum to \(\Ran\iota\) increases
the value, so that
\[
 \iota^*(A^{-1}+B^{-1})^{-1}\iota
 \le\bigl((\iota^*A\iota)^{-1}+(\iota^*B\iota)^{-1}\bigr)^{-1},
\]
and the analogous formula for \((\sum_jA_j^{-1})^{-1}\) gives monotonicity
and the compression inequality for more than two terms. The integral representation
\[
 A\#_\tau B
 =\frac{\sin\pi\tau}{\pi}
   \int_0^\infty v^{\tau-1}(A^{-1}+vB^{-1})^{-1}\,dv,
\]
which for commuting \(A\) and \(B\) reduces to \(A^{1-\tau}B^\tau\),
writes \(A\#_\tau B\) through the parallel sums of \(A\) and
\(v^{-1}B\), and therefore gives
\[
 \iota^*(A\#_\tau B)\iota
 \le(\iota^*A\iota)\#_\tau(\iota^*B\iota).
\]

\smallskip
\emph{Step 2: the reserve inequality \eqref{eq:low-compression}.} For
\(A=G\), \(B=\overline G\), the two operators commute both before
and after compression, because their sum is \(2I\) and
\(\iota^*(2I)\iota=2I_F\). Commuting positive operators satisfy
\(A\#B=(AB)^{1/2}\), so \(G\#\overline G=\Res\) and
\(G_F\#\overline G_F=\Res_F\), and Step~1 at \(\tau=1/2\) (for invertible
operators; Step~4 removes this restriction) gives
\(\Res_F\ge\iota^*\Res\iota\). Since \(0<2a\le1\), the function
\(x^{2a}\) is operator monotone and operator concave on \([0,\infty)\),
and these two properties give
\[
 \Res_F^{2a}
 \overset{\text{monotone}}{\ge}(\iota^*\Res\iota)^{2a}
 \overset{\text{Jensen}}{\ge}\iota^*\Res^{2a}\iota,
\]
which is \eqref{eq:low-compression}. The map \(X\mapsto\iota^*X\iota\) is
completely positive and unital, as required for this use of Jensen's
inequality in the form of Hansen \cite[Corollary]{Hansen80}.

Since \(i\mathsf J/2\) is self-adjoint, the squares of the reserves
also compare directly:
\begin{align*}
 \Res_F^2-\iota^*\Res^2\iota
 &=\bigl(I_F-(\iota^*(i\mathsf J/2)\iota)^2\bigr)
  -\bigl(I_F-\iota^*(i\mathsf J/2)^2\iota\bigr)\\
 &=\iota^*(i\mathsf J/2)(I-\iota\iota^*)(i\mathsf J/2)\iota\ge0.
\end{align*}
Applying operator monotonicity and concavity to \(x^a\) gives a
second route to \eqref{eq:low-compression}:
\[
 \Res_F^{2a}\ge(\iota^*\Res^2\iota)^a\ge\iota^*\Res^{2a}\iota.
\]

\smallskip
\emph{Step 3: the thermal weights and \eqref{eq:high-compression}.}
The thermal increment \(\Therm_{(r-1)/r}(I)\) is a parallel sum of
geometric means of \(G\) and \(\overline G\):
\begin{equation}
 \Therm_{(r-1)/r}(I)
 =
 \left[\sum_{k=1}^{r-1}
       (G\#_{k/r}\overline G)^{-1}\right]^{-1}.
 \label{eq:parallel-thermal}
\end{equation}
At a scalar spectral value \(d\in(-1,1)\) of \(i\mathsf J/2\), here
signed, let
\(g_\pm=1\pm d\) and \(\mathfrak z=(g_-/g_+)^{1/r}\), so that
\(G\#_{k/r}\overline G\) takes the value
\(g_+^{1-k/r}g_-^{k/r}=g_+\mathfrak z^k\). The right-hand side of
\eqref{eq:parallel-thermal} takes the value
\[
 \frac{g_+}{\sum_{k=1}^{r-1}\mathfrak z^{-k}}
 =\frac{g_+\mathfrak z^{r-1}(1-\mathfrak z)}{1-\mathfrak z^{r-1}}
 =d\coth\!\left(\frac{r-1}{r}\arctanh d\right)-1.
\]
For the last equality, write \(d=\tanh u\); then
\(g_+=e^u/{\cosh u}\), \(\mathfrak z=e^{-2u/r}\),
\(1-\mathfrak z=2e^{-u/r}\sinh(u/r)\) and
\(1-\mathfrak z^{r-1}=2e^{-(r-1)u/r}\sinh((r-1)u/r)\), so the middle
expression equals
\[
 \sinh(u/r)/(\cosh u\sinh((r-1)u/r)),
\]
which is
the right-hand side by Step~1 of the proof of
Lemma~\ref{lem:thermal-comparison} at \(\tau=(r-1)/r\). The expression
\(d\coth(\frac{r-1}r\arctanh d)-1\) is even in \(d\) and extends
continuously at zero and the endpoints. Compression of each geometric
mean in \eqref{eq:parallel-thermal}, followed by monotonicity and
compression of the parallel sum, gives the first inequality in
\eqref{eq:high-compression}, and Lemma~\ref{lem:thermal-comparison}
gives the second; with \eqref{eq:parallel-thermal} on the finite
source space, where \(G_F\) and \(\overline G_F\) take the places of
\(G\) and \(\overline G\) (because \(\iota\) preserves \(\varsigma\),
\(\iota^*\mathsf J\iota\) is the operator \(\mathsf J\) of the finite source
space, so \(G_F\), \(\overline G_F\) and \(\Res_F\) are its own \(G\),
\(\overline G\) and \(\Res\)),
\begin{align*}
 \Therm_b(I_F)
 &\overset{\eqref{eq:parallel-thermal}}{=}
 \Bigl[\sum_{k=1}^{r-1}(G_F\#_{k/r}\overline G_F)^{-1}\Bigr]^{-1}
 \overset{\text{Step 1}}{\ge}
 \Bigl[\sum_{k=1}^{r-1}
  \bigl(\iota^*(G\#_{k/r}\overline G)\iota\bigr)^{-1}\Bigr]^{-1}\\
 &\overset{\text{Step 1}}{\ge}
 \iota^*\Bigl[\sum_{k=1}^{r-1}(G\#_{k/r}\overline G)^{-1}\Bigr]^{-1}\iota
 \overset{\eqref{eq:parallel-thermal}}{=}\iota^*\Therm_b(I)\iota
 \overset{\text{Lem.~\ref{lem:thermal-comparison}}}{\ge}
 c_b\,\iota^*\Res^{2b}\iota.
\end{align*}

\smallskip
\emph{Step 4: the full source space.} Positive regularization and form
limits yield \eqref{eq:low-compression} and \eqref{eq:high-compression}
in infinite dimension. On the full source
space, \(G\) and \(\overline G\) are injective but need not have
bounded inverses. The inequalities of Steps 1 to 3 hold for the
commuting invertible operators \(G+\epsilon I\) and
\(\overline G+\epsilon I\), \(\epsilon>0\), in place of \(G\) and
\(\overline G\). As \(\epsilon\downarrow0\), the geometric means and
parallel sums formed from \(G+\epsilon I\) and \(\overline G+\epsilon I\)
decrease strongly to those of \(G\) and
\(\overline G\), which is how the noninvertible cases are defined,
while their compressions \(G_F+\epsilon I_F\) and
\(\overline G_F+\epsilon I_F\) converge to the invertible \(G_F\) and
\(\overline G_F\). The inequalities therefore pass to the limit, and
\eqref{eq:parallel-thermal} holds on the full source space as the
scalar identity above in the spectral calculus of \(i\mathsf J/2\).
\end{proof}

The comparisons of Lemmas~\ref{lem:thermal-comparison}
and~\ref{lem:compression} also control trace norms after normalization
by a thermal increment, as in the matrix \(\mathsf Z\) of
Theorem~\ref{thm:gaussian}.
Let \(H_F\) be a positive operator on the finite source space with
\(H_F\ge c_a\iota^*\Res^{2a}\iota\); in
Section~\ref{sec:completion}, \(H_F\) is the thermal increment
\(\Therm_a(I_F)\), which satisfies this hypothesis by
\eqref{eq:scalar-reserve} on the finite source space and
\eqref{eq:low-compression}, as displayed in
\eqref{eq:joint-compression}. For an operator \(\Res^aY\Res^a\) with
\(Y\in\Sc_1\),
\[
 H_F^{-1/2}\iota^*\Res^aY\Res^a\iota H_F^{-1/2}
 =\mathsf N^*Y\mathsf N,\qquad
 \mathsf N=\Res^a\iota H_F^{-1/2},\qquad
 \norm{\mathsf N}\le c_a^{-1/2},
\]
where the norm bound follows from the hypothesis, conjugated by
\(H_F^{-1/2}\):
\[
 \mathsf N^*\mathsf N
 =H_F^{-1/2}\iota^*\Res^{2a}\iota H_F^{-1/2}
 \le c_a^{-1}I_F.
\]
The trace norm of \(\mathsf N^*Y\mathsf N\) is therefore at most
\(\norm{\mathsf N}^2\norm{Y}_1\le c_a^{-1}\norm{Y}_1\) (the
\emph{one-space normalization bound}).

For two full source spaces with reserves \(\Res_A,\Res_B\) and real
isometries \(\iota_A,\iota_B\) into them that preserve \(\varsigma\), as in
Lemma~\ref{lem:compression},
let \(H_{A,F}\) and \(H_{B,F}\) be positive operators with
\(H_{A,F}\ge c_b\iota_A^*\Res_A^{2b}\iota_A\) and
\(H_{B,F}\ge c_b\iota_B^*\Res_B^{2b}\iota_B\); in
Section~\ref{sec:completion}, \(H_{A,F}\) and \(H_{B,F}\) are the
thermal increments \(\Therm_b(I_F)\) of the two finite source spaces,
which satisfy these hypotheses by \eqref{eq:high-compression}.
Conjugation by \(H_{A,F}^{-1/2}\) and \(H_{B,F}^{-1/2}\), as for
\(\mathsf N\), bounds the norms of
\(\mathsf N_A=\Res_A^b\iota_AH_{A,F}^{-1/2}\) and
\(\mathsf N_B=\Res_B^b\iota_BH_{B,F}^{-1/2}\) by \(c_b^{-1/2}\). The
operator in \eqref{eq:trace-compression} below is
\(\mathsf N_A^*X\mathsf N_B\), so its trace norm is at most
\(\norm{\mathsf N_A}\norm{X}_1\norm{\mathsf N_B}\), and the bounds on
\(\norm{\mathsf N_A}\) and \(\norm{\mathsf N_B}\) give
\begin{equation}
 \norm{
 H_{A,F}^{-1/2}\iota_A^*\Res_A^bX\Res_B^b\iota_BH_{B,F}^{-1/2}
 }_1
 \le c_b^{-1}\norm{X}_1,
 \qquad X\in\Sc_1.
 \label{eq:trace-compression}
\end{equation}

\section{Bounded localization and Fourier evaluation}
\label{sec:evaluation}

The complex mass shell is
\[
 \mathcal Q_\mu
 =\{z\in\C^4:z_0^2-z_1^2-z_2^2-z_3^2=\mu^2\}.
\]
For a compactly supported source \(\sx=(f,g)\), we use the entire
function
\[
 \Psi_\sx(z)=\wh f(\bz)+iz_0\wh g(\bz).
\]
The spacetime product \(z\cdot z\) is Lorentzian, while
\(\bz\cdot\bz\) is the complex bilinear Euclidean spatial product.

For a nonempty bounded open \(U\subset\R^3\), \(\sx=(f,g)\in\Hh_{\mu,U}\)
and \(z\in\C^4\), the number \(\Psi_\sx(z)\) depends continuously on
\(\sx\). Indeed, \(\wh f(\bz)\) and \(\wh g(\bz)\) are the values of the
distributions \(f\in H^{-1/2}\) and \(g\in H^{1/2}\) on a function in
\(C_c^\infty(\R^3)\) equal to \(e^{-i\bz\cdot\bx}\) near \(\overline U\), so
the two values are bounded by \(\norm{f}_{H^{-1/2}}\) and \(\norm{g}_{H^{1/2}}\) times
the \(H^{1/2}\) and \(H^{-1/2}\) norms of that function. Since
\((\mu^2+|\bp|^2)/(1+|\bp|^2)\) lies between \(\min\{1,\mu^2\}\) and
\(\max\{1,\mu^2\}\), the norms \(\norm{f}_{H^{-1/2}}\) and
\(\norm{g}_{H^{1/2}}\) are at most a factor depending only on \(\mu\) times
the norm of \(\sx\) in \(\Hh_{\mu,U}\).

\subsection{Containing wedges and weighted complex-shell evaluation}

For \(\ell\ge0\), let
\[
 \mathcal W_\ell
 =\{(x^0,\bx):x^1<\ell-|x^0|\}.
\]
The real subspace \(\mathcal K_{\mathcal W_\ell}\) of the wedge is the
closed real linear span of the subspaces \(\mathcal K_U\) of
Lemma~\ref{lem:standard} over the bounded open sets \(U\) of its time-zero
base, the half-space \(\{x^1<\ell\}\); by finite propagation speed and the
time-slice construction that proves \eqref{eq:time-slice}, the same subspace is also the
closed real span of the one-particle vectors of the real test functions
supported in \(\mathcal W_\ell\).

The subspace \(\mathcal K_{\mathcal W_\ell}\)
is standard. It contains \(\mathcal K_U\) for a ball \(U\), so the complex
span of \(\mathcal K_{\mathcal W_\ell}\) is dense by Lemma~\ref{lem:standard}. Let
\(\xi\in\mathcal K_{\mathcal W_\ell}\cap i\mathcal K_{\mathcal W_\ell}\), so
that \(-i\xi\) lies in \(\mathcal K_{\mathcal W_\ell}\) as well, and let
\(\sy\) be a smooth real source supported in a bounded open
\(U'\subset\{x^1>\ell\}\). Sources supported in the two half-spaces have
vanishing symplectic form, so by \eqref{eq:gram} the numbers
\(\ip{\xi}{\kappa_{U'}\sy}\) and
\(\ip{-i\xi}{\kappa_{U'}\sy}=i\ip{\xi}{\kappa_{U'}\sy}\) are real; hence
\(\ip{\xi}{\kappa_{U'}\sy}=0\). The vectors \(\kappa_{U'}\sy\) span a dense
complex subspace by Lemma~\ref{lem:standard}, so \(\xi=0\). We write
\(\Delta_{\mathcal W_\ell}\) for the modular operator of
\(\mathcal K_{\mathcal W_\ell}\).

The Bisognano--Wichmann theorem identifies the wedge modular group with
Lorentz boosts \cite{BW75}, and for the free field the same
identification holds in the one-particle space
\cite[Remark~3.6(2)]{BGL02}: for the left wedge \(\mathcal W_0\),
\[
 \Delta_{\mathcal W_0}^{iy}=\mathsf U(\Lambda_1(2\pi y)),\qquad
 (\Lambda_1(\chi)p)_0=p_0\cosh\chi+p_1\sinh\chi,\quad
 (\Lambda_1(\chi)p)_1=p_1\cosh\chi+p_0\sinh\chi,
\]
where \(\mathsf U\) is the one-particle representation,
\((\mathsf U(\Lambda)\psi)(p)=\psi(\Lambda^{-1}p)\), and \(\Lambda_1(\chi)\)
is the boost of rapidity \(\chi\) along \(x^1\), which preserves
\(\mathcal W_0\). For the right wedge the sign of the rapidity is
reversed; the choice of the left wedge fixes the sign of the imaginary
shift below. Analytic continuation to \(y=-ic\), \(0<c<1/2\),
replaces the real boost by the boost of imaginary rapidity \(i\phi\),
\(\phi=2\pi c\in(0,\pi)\): on one-particle wave functions \(\psi\) in the
domain of \(\Delta_{\mathcal W_0}^{c}\),
\(\Delta_{\mathcal W_0}^{c}\psi=\psi\circ b_\phi\) with
\[
 b_\phi(p)=\Lambda_1(i\phi)p=
 (p_0\cos\phi+ip_1\sin\phi,\,
  p_1\cos\phi+ip_0\sin\phi,\,
  p_2,p_3).
\]

With the Fourier convention of Section~\ref{sec:setting},
\begin{equation}
 \norm{\Delta_{\mathcal W_\ell}^{\,c}\kappa_U\sx}^2
 =
 \int_{H_\mu^+}e^{-2\ell p_0\sin\phi}
       |\Psi_\sx(b_\phi(p))|^2\,d\mathfrak m_\mu(p)
 \label{eq:wedge-contour}
\end{equation}
for every \(\sx\in\Hh_{\mu,U}\), \(U\) a bounded open subset of the base
\(\{x^1<\ell\}\). A proper orthochronous Lorentz
transformation \(\mathcal L\) replaces the argument \(b_\phi(p)\) of \(\Psi_\sx\) by
\(\mathcal Lb_\phi(p)\) for the wedge \(\mathcal L\mathcal W_\ell\).

Steps~1 to~5 below prove \eqref{eq:wedge-contour} and its form for
\(\mathcal L\mathcal W_\ell\), reading the contour and its damping
directly from the one-particle representation, as in the similar
computation of Buchholz and Lechner for the free massive field in two
dimensions at \(c=1/4\) \cite[Section~4]{BL04}. In rapidity coordinates
\[
 p_0=m_\perp\cosh\chi,\qquad p_1=m_\perp\sinh\chi,\qquad
 m_\perp^2=\mu^2+p_2^2+p_3^2,
\]
the invariant measure is
\(d\mathfrak m_\mu=(2(2\pi)^3)^{-1}d\chi\,d^2p_\perp\), since
\(dp_1=p_0\,d\chi\).

\smallskip
\emph{Step 1: the modular operator as a shift of rapidity.} Boosts
translate \(\chi\):
\((\Delta_{\mathcal W_0}^{iy}\psi)(\chi,\bp_\perp)=\psi(\chi-2\pi y,\bp_\perp)\).
The Fourier transform \(\int e^{-ik\chi}\psi(\chi,\bp_\perp)\,d\chi\) turns
\(\Delta_{\mathcal W_0}^{iy}\) into multiplication by \(e^{-2\pi iyk}\), so
\(\Delta_{\mathcal W_0}^{c}\) is multiplication by \(e^{-\phi k}\). By the
Paley--Wiener theorem for a strip, applied at almost every \(\bp_\perp\), a wave
function \(\psi\) lies in the domain of \(\Delta_{\mathcal W_0}^{c}\), with
\[
 (\Delta_{\mathcal W_0}^{c}\psi)(\chi,\bp_\perp)=\psi(\chi+i\phi,\bp_\perp),
\]
whenever each function \(\chi\mapsto\psi(\chi,\bp_\perp)\) extends
holomorphically to the strip \(0<\Ima\chi<\phi\) and continuously to its
closure, with \(|\psi(\chi+i\phi',\bp_\perp)|\le G(\chi,\bp_\perp)\) for
\(0\le\phi'\le\phi\) and a function \(G\) square integrable in
\(d\chi\,d^2p_\perp\). Since
\(\Lambda_1(\chi+i\phi)=\Lambda_1(i\phi)\Lambda_1(\chi)\), the shift
\(\chi\mapsto\chi+i\phi\) replaces \(p\) by \(b_\phi(p)\).

\smallskip
\emph{Step 2: smooth sources in the base of \(\mathcal W_0\).} Let
\(\sx=(f,g)\) be smooth with compact support in \(\{x^1<0\}\), and let
\(z=\Lambda_1(\chi+i\phi')(m_\perp,0,p_2,p_3)\) with \(0\le\phi'\le\phi\).
Then \(\Ima\bz=(m_\perp\cosh\chi\sin\phi',0,0)\) has a nonnegative first
entry, so \(|e^{-i\bz\cdot\bx}|=e^{x^1\Ima z_1}\le1\) on the support of
\(\sx\). Integration by parts in
\(\wh f(\bz)=\int e^{-i\bz\cdot\bx}f(\bx)\,d^3x\) and in \(\wh g(\bz)\),
together with \(|z_0|^2=|\mu^2+\bz\cdot\bz|\le(\mu+|\bz|)^2\), gives for
every \(N\) a constant \(C_{N,\sx}\) with
\[
 |\Psi_\sx(z)|\le C_{N,\sx}(1+|\bz|)^{-N}.
\]
Since \(|\bz|\ge\max\{|z_1|,|\bp_\perp|\}\) with
\(|z_1|=m_\perp\lvert\sinh(\chi+i\phi')\rvert\ge m_\perp\lvert\sinh\chi\rvert\), and \(m_\perp\ge\mu>0\), the right-hand side is at most a function of
\((\chi,\bp_\perp)\) alone that is square integrable in
\(d\chi\,d^2p_\perp\) for \(N>2\). The function
\(\chi\mapsto\Psi_\sx(\Lambda_1(\chi)(m_\perp,0,p_2,p_3))\) is entire and
equals \(\kappa_U\sx\) at real \(\chi\), so Step 1 gives
\(\Delta_{\mathcal W_0}^{c}\kappa_U\sx=\Psi_\sx\circ b_\phi\).

\smallskip
\emph{Step 3: translation of the wedge edge.}
Translation of the wedge edge by \(\ell\) contributes the modulus
\(e^{-\ell p_0\sin\phi}\). Since
\(\mathcal W_\ell=\mathcal W_0+\ell\eone\), one has
\(\Delta_{\mathcal W_\ell}=\mathsf U(\ell\eone)\Delta_{\mathcal W_0}
\mathsf U(\ell\eone)^*\), where \((\mathsf U(\ell\eone)\psi)(p)=e^{-i\ell p_1}\psi(p)\)
represents the translation by \(\ell\eone\). The inner factor
\(\mathsf U(\ell\eone)^*\) maps \(\kappa_U\sx\) to the one-particle vector of the
source translated by \(-\ell\eone\), whose Fourier transform is
\(e^{i\ell p_1}\Psi_\sx(p)\); the continuation
\(p_1\mapsto p_1\cos\phi+ip_0\sin\phi\) turns the factor \(e^{i\ell p_1}\) into
\(e^{i\ell p_1\cos\phi}e^{-\ell p_0\sin\phi}\), and the outer
translation is unitary. For smooth \(\sx\) with compact support in
\(\{x^1<\ell\}\), Step 2 applies to the translated source, and
\[
 |(\Delta_{\mathcal W_\ell}^{c}\kappa_U\sx)(p)|
 =e^{-\ell p_0\sin\phi}|\Psi_\sx(b_\phi(p))|,
\]
which gives \eqref{eq:wedge-contour} for such \(\sx\).

\smallskip
\emph{Step 4: Lorentz-moved wedges.} For a proper orthochronous
\(\mathcal L\), the subspace
\(\mathcal K_{\mathcal L\mathcal W_\ell}=\mathsf U(\mathcal L)\mathcal K_{\mathcal W_\ell}\)
is standard, with modular operator
\(\Delta_{\mathcal L\mathcal W_\ell}=\mathsf U(\mathcal L)\Delta_{\mathcal W_\ell}\mathsf U(\mathcal L)^*\).
Let \(U\subset B_T(0)\) with \(D_T\subset\mathcal L\mathcal W_\ell\), let
\(\sx\) be smooth with compact support in \(U\), let \(u\) be the
Klein--Gordon solution determined by \(\sx\), and let \(\sx_{\mathcal L}\)
be the time-zero pair of the solution \(x\mapsto u(\mathcal Lx)\). If
\(\bx\) lies in the support of \(\sx_{\mathcal L}\), then \(\mathcal L(0,\bx)\)
lies in the support of \(u\), whose points are, by finite propagation
speed, causally related to points of \(\{0\}\times\supp\sx\subset D_T\); so
\((0,\bx)\) is causally related to a point
\(q\in\mathcal L^{-1}D_T\subset\mathcal W_\ell\), and
\(x^1\le q^1+|q^0|<\ell\). Hence \(\sx_{\mathcal L}\) is smooth with compact
support in a bounded open \(U'\subset\{x^1<\ell\}\), and real if \(\sx\) is
real. By the Lorentz covariance of the free field,
\(\kappa_{U'}\sx_{\mathcal L}=\mathsf U(\mathcal L)^*\kappa_U\sx\), that is,
\(\Psi_{\sx_{\mathcal L}}(p)=\Psi_\sx(\mathcal Lp)\) on \(H_\mu^+\); so, for
real \(\sx\),
\(\kappa_U\sx\in\mathsf U(\mathcal L)\mathcal K_{U'}\subset
\mathcal K_{\mathcal L\mathcal W_\ell}\), and by closure
\(\mathcal K_U\subset\mathcal K_{\mathcal L\mathcal W_\ell}\).

The functions
\(\Psi_{\sx_{\mathcal L}}\) and \(\Psi_\sx\circ\mathcal L\) are holomorphic on
the connected complex manifold \(\mathcal Q_\mu\) and agree on \(H_\mu^+\), a
totally real submanifold of maximal dimension, so they agree on
\(\mathcal Q_\mu\). Step 3 for \(\sx_{\mathcal L}\) therefore gives
\[
 |(\Delta_{\mathcal W_\ell}^{c}\mathsf U(\mathcal L)^*\kappa_U\sx)(p)|
 =e^{-\ell p_0\sin\phi}|\Psi_\sx(\mathcal Lb_\phi(p))|,
\]
and \(\norm{\Delta_{\mathcal L\mathcal W_\ell}^{c}\kappa_U\sx}
=\norm{\Delta_{\mathcal W_\ell}^{c}\mathsf U(\mathcal L)^*\kappa_U\sx}\).

\smallskip
\emph{Step 5: all sources, by closure.} Let \(\sx\in\Hh_{\mu,U}\), with
\(U\) and \(\mathcal L\) as in Step 4, or with \(U\) a bounded open subset
of \(\{x^1<\ell\}\) and \(\mathcal L\) the identity. The Tomita operator
\(\mathsf s\) of \(\mathcal K_{\mathcal W_\ell}\) fixes every
\(\xi\in\mathcal K_{\mathcal W_\ell}\), so
\(\lVert\Delta_{\mathcal W_\ell}^{1/2}\xi\rVert=\lVert\mathsf s\xi\rVert=\lVert\xi\rVert\),
and \(\lambda^c\le1+\lambda^{1/2}\) for \(\lambda\ge0\) gives
\(\lVert\Delta_{\mathcal W_\ell}^{c}\xi\rVert\le2\lVert\xi\rVert\). Choose smooth
sources \(\sx_n\) with compact support in \(U\) and \(\sx_n\to\sx\), and
write \(\sx_n-\sx=\sy_n+i\sy_n'\) with real \(\sy_n,\sy_n'\to0\). The vectors
\(\mathsf U(\mathcal L)^*\kappa_U\sy_n\) and
\(\mathsf U(\mathcal L)^*\kappa_U\sy_n'\) lie in \(\mathcal K_{\mathcal W_\ell}\),
since \(\mathsf U(\mathcal L)^*\mathcal K_U\subset\mathcal K_{\mathcal W_\ell}\),
and have the norms of \(\sy_n\) and \(\sy_n'\) by \eqref{eq:gram}; so, by the bound
\(\lVert\Delta_{\mathcal W_\ell}^{c}\xi\rVert\le2\lVert\xi\rVert\),
\(\Delta_{\mathcal W_\ell}^{c}\mathsf U(\mathcal L)^*\kappa_U\sx_n\) tends to
\(\Delta_{\mathcal W_\ell}^{c}\mathsf U(\mathcal L)^*\kappa_U\sx\), and a
subsequence converges almost everywhere. For every
\(p\), \(\Psi_{\sx_n}(\mathcal Lb_\phi(p))\to\Psi_\sx(\mathcal Lb_\phi(p))\) by
the continuity shown at the beginning of Section~\ref{sec:evaluation}, so
the moduli of Steps 3 and 4
hold for \(\sx\), and with them \eqref{eq:wedge-contour} for every
\(\sx\in\Hh_{\mu,U}\), together with its form for \(\mathcal L\mathcal W_\ell\),
in which \(\Psi_\sx(\mathcal Lb_\phi(p))\) replaces \(\Psi_\sx(b_\phi(p))\).

For \(U\subset B_T(0)\) and a compact Lorentz family \(\mathscr L\), a
choice
\[
 \ell_{\mathscr L}>
 \sqrt2\sup_{\mathcal L\in\mathscr L}
       \norm{\mathcal L^{-1}}_{\mathrm E}
\]
ensures \(D_T\subset\mathcal L\mathcal W_{\ell_{\mathscr L}T}\) for every
\(\mathcal L\in\mathscr L\): the double cone \(D_T\) lies in the Euclidean
ball of radius \(T\), since \((x^0)^2+|\bx|^2\le(|x^0|+|\bx|)^2<T^2\) on
\(D_T\), so \(\mathcal L^{-1}D_T\) lies in the ball of radius
\(\norm{\mathcal L^{-1}}_{\mathrm E}T\), and every \(x\) with
\(|x|_{\mathrm E}\le r\) satisfies
\(x^1+|x^0|\le\sqrt2\,((x^0)^2+(x^1)^2)^{1/2}\le\sqrt2\,r\).
For each \(\mathcal L\in\mathscr L\), Step 4 therefore gives the inclusion
\(\mathcal K_U\subset\mathcal K_{\mathcal L\mathcal W_{\ell_{\mathscr L}T}}\)
of standard subspaces that Lemma~\ref{lem:inclusion} requires, and Step 5
gives, for \(\sx\in\Hh_{\mu,U}\),
\begin{align*}
 \int_{H_\mu^+}e^{-2\ell_{\mathscr L}Tp_0\sin\phi}
 |\Psi_\sx(\mathcal Lb_\phi(p))|^2\,d\mathfrak m_\mu(p)
 &\overset{\eqref{eq:wedge-contour}}{=}
 \norm{\Delta_{\mathcal L\mathcal W_{\ell_{\mathscr L}T}}^{\,c}\kappa_U\sx}^2\\
 &\overset{\text{Lem.~\ref{lem:inclusion}}}{\le}
 \norm{h_{c,U}\sx}^2.
\end{align*}
The integrand on the left is nonnegative and continuous in
\((\mathcal L,p)\), so by Tonelli's theorem the left-hand side is a
measurable function of \(\mathcal L\).
Integrating the bound of Lemma~\ref{lem:inclusion} against a probability measure
\(\vartheta\) on \(\mathscr L\) gives
\begin{equation}
 \int_{\mathscr L}\int_{H_\mu^+}
 e^{-2\ell_{\mathscr L}Tp_0\sin\phi}
 |\Psi_\sx(\mathcal Lb_\phi(p))|^2
 \,d\mathfrak m_\mu(p)\,d\vartheta(\mathcal L)
 \le\norm{h_{c,U}\sx}^2.
 \label{eq:averaged-wedges}
\end{equation}

\begin{theorem}[Weighted complex-shell evaluation]\label{thm:evaluation}
For every \(0<c<1/2\), there are finite constants \(C_c,\lambda_c\ge1\),
depending only on \(c\), such that for every \(\mu,T>0\), every nonempty
bounded open \(U\subset B_T(0)\), every \(\sx\in\Hh_{\mu,U}\), and every \(\Energy\ge\mu\),
\begin{equation}
 |\Psi_\sx(z)|^2
 \le C_c\Energy^{-2}e^{2\lambda_cT\Energy}
       \norm{h_{c,U}\sx}^2,
 \qquad
 z\in\mathcal Q_\mu,\quad |z|_{\mathrm E}\le2\Energy.
 \label{eq:evaluation}
\end{equation}
Here \(\Psi_\sx(z)=\wh f(\bz)+iz_0\wh g(\bz)\) for \(\sx=(f,g)\), and
\(h_{c,U}\) is the bounded operator of \eqref{eq:modular-source}. The
constants are independent of the shape and connectedness of \(U\) and
of the dimension of any finite canonical source system in \(U\). The choice
\[
 \lambda_c=\frac{2048}{\sin(2\pi c)}
\]
is admissible.
\end{theorem}

The averaged contour norm of \eqref{eq:averaged-wedges} controls \(\Psi_\sx\) on an open shell
neighbourhood (Lemma~\ref{lem:regular-patches}), and analytic discs
reach the remaining evaluation points (Lemma~\ref{lem:null-discs}).

\begin{lemma}[Regular shell patches]\label{lem:regular-patches}
For \(0<c<1/2\) and \(\phi=2\pi c\), consider the compact normalized set
\begin{equation}
 \mathscr Z=
 \left\{(t,z):
 0\le t\le1,\ z\in\mathcal Q_t,\ |z|_{\mathrm E}\le16,\ 
 (\Rea z)^2\le-16\right\}.
 \label{eq:compact-target}
\end{equation}
There are a compact set \(\mathscr L_c\) of proper orthochronous
Lorentz transformations, a smooth probability density \(\vartheta_c\)
on \(\mathscr L_c\) with respect to Haar measure, a neighbourhood
\(\mathcal V\) of \(\mathscr Z\) in the shell family
\(\{(t,z):0\le t\le1,\ z\in\mathcal Q_t\}\), and \(\delta_c>0\), such
that the map \(\Xi_t(\mathcal L,p)=\mathcal Lb_\phi(p)\) on
\(\mathscr L_c\times H_t^+\) satisfies
\begin{equation}
 (\Xi_t)_*
 \left(d\vartheta_c\otimes
       \mathbf1_{\{p_0\le M_c\}}\,d\mathfrak m_t\right)
 \ge\delta_c\,d\vol_{\mathcal Q_t}
 \quad\text{on }\mathcal V_t,
 \qquad M_c=\frac{32}{\sin\phi}.
 \label{eq:pushforward}
\end{equation}
Here \(d\vartheta_c\) is the probability measure with density
\(\vartheta_c\), \(d\vol_{\mathcal Q_t}\) is induced Euclidean
six-dimensional volume, \(\mathcal V_t=\{z:(t,z)\in\mathcal V\}\) is the
slice of \(\mathcal V\) at normalized mass \(t\), and at \(t=0\) the
shell \(H_0^+\) and its measure \(d\mathfrak m_0\) are given by the
formulas of Section~\ref{sec:setting}, so that \(H_0^+\) is the forward
light cone without its vertex. The
family \(\mathscr L_c\) can be chosen with
\[
 \sup_{\mathcal L\in\mathscr L_c}
 \norm{\mathcal L^{-1}}_{\mathrm E}<\frac{32}{\sin\phi},
\]
so that \(\ell_c=64/{\sin\phi}\) exceeds
\(\sqrt2\sup_{\mathcal L\in\mathscr L_c}\norm{\mathcal L^{-1}}_{\mathrm E}\),
as required of the containing-wedge parameter \(\ell_{\mathscr L}\) for
\eqref{eq:averaged-wedges} with \(\mathscr L=\mathscr L_c\).
Moreover, there are \(r>0\) and \(C_{\mathrm{sub}}\ge1\) such that every
\((t,z)\in\mathscr Z\) has a holomorphic shell chart of radius \(r\): an
injective holomorphic map \(\psi_{t,z}\) from a ball of radius \(r\) in
\(\C^3\) into \(\mathcal Q_t\) that maps the centre of the ball to \(z\),
whose image \(\mathcal B_z\) lies in \(\mathcal V_t\), and whose volume
Jacobian with respect to Lebesgue measure lies between \(1\) and
\(C_{\mathrm{sub}}\). In particular, \(\vol(\mathcal B_z)\ge\vol(B_r)\),
where \(B_r\) is a ball of radius \(r\) in \(\C^3=\R^6\), and every
nonnegative plurisubharmonic function \(F\) on \(\C^4\) satisfies
\[
 F(z)\le\frac{C_{\mathrm{sub}}}{\vol(\mathcal B_z)}
 \int_{\mathcal B_z}F\,{d\vol_{\mathcal Q_t}}.
\]
\end{lemma}

\begin{proof}
\emph{Step 1: a preimage of every point of \(\mathscr Z\).}
A point \(z=u+iv\in\mathcal Q_t\) satisfies
\[
 u\cdot v=0,\qquad u^2-v^2=t^2.
\]
On \(\mathscr Z\), writing \(u^2=-\rho_0^2\) with \(\rho_0>0\), the
Lorentzian length of \(u\), gives \(\rho_0\in[4,16]\) and \(v^2=-(t^2+\rho_0^2)\):
\(\rho_0\ge4\) because \(u^2\le-16\), and
\(\rho_0^2=|\mathbf u|^2-u_0^2\le|z|_{\mathrm E}^2\le256\), with \(\mathbf u\) the
spatial part of \(u\).

The ordered pair \(u,v\) spans a negative-definite real
two-dimensional subspace. This pair is a Lorentz image of the real and
imaginary parts of \(b_\phi(p)\) at
\[
 p_1=p_3=0,\qquad
 p_0=\frac{\sqrt{t^2+\rho_0^2}}{\sin\phi},\qquad
 p_2=\sqrt{p_0^2-t^2}.
\]
Indeed, at this momentum the \emph{reference pair}
\[
 \Rea b_\phi(p)=(p_0\cos\phi,0,p_2,0),\qquad
 \Ima b_\phi(p)=(0,p_0\sin\phi,0,0)
\]
is orthogonal. Since \(p_0^2=t^2+p_2^2\) and
\(p_0^2\sin^2\phi=t^2+\rho_0^2\), its Lorentzian squares are
\[
 (\Rea b_\phi(p))^2=p_0^2\cos^2\phi-p_2^2
 =t^2\cos^2\phi-p_2^2\sin^2\phi
 =t^2-p_0^2\sin^2\phi=-\rho_0^2
\]
and \((\Ima b_\phi(p))^2=-p_0^2\sin^2\phi=-(t^2+\rho_0^2)\), the same as those
of the \emph{target pair} \((u,v)\). Moreover,
\(p_2^2=p_0^2-t^2\ge(t^2+\rho_0^2)-t^2=\rho_0^2\), so \(|\bp|=p_2\ge4\) for
every \(t\in[0,1]\): the preimages stay at distance at least \(4\) from the
vertex of \(H_t^+\).

The \emph{source frame}
\(\mathcal L_{\mathrm{src}}\) is the proper orthochronous
transformation whose columns are a future-directed unit timelike
vector, the two unit vectors \(\Rea b_\phi(p)/\rho_0\) and
\(\Ima b_\phi(p)/(t^2+\rho_0^2)^{1/2}\), and a unit spacelike vector
completing the orthonormal frame; the \emph{target frame}
\(\mathcal L_{\mathrm{tar}}\) is built in the same way from
\(u/\rho_0\) and \(v/(t^2+\rho_0^2)^{1/2}\). The \emph{matching
transformation}
\(\mathcal L=\mathcal L_{\mathrm{tar}}\mathcal L_{\mathrm{src}}^{-1}\)
is proper orthochronous and carries the reference pair to the target
pair, so that \(\Xi_t(\mathcal L,p)=z\).

\smallskip
\emph{Step 2: \(\Xi_t\) is a submersion at these preimages.} The
differential of \(\Xi_t\) at \((\mathcal L,p)\) sends \((X,\delta p)\), with
\(X\) in the Lie algebra of the Lorentz group and \(\delta p\) tangent to
\(H_t^+\), to \(Xz+\mathcal Lb_\phi(\delta p)\).
The Lorentz orbit of the ordered pair \((u,v)\) has dimension five: the
stabilizer of the pair acts on the orthogonal Lorentzian two-dimensional subspace
and has dimension one, since an \(X\) with \(Xu=Xv=0\) maps that subspace
into itself and is antisymmetric there. The vectors \(Xz\) therefore form
a five-dimensional subspace of the tangent space of \(\mathcal Q_t\) at
\(z\), on which the differential of the Lorentz invariant \((\Rea z)^2\)
vanishes.

The additional variation in \(p_2\), with \(p_1=p_3=0\) and
\(p_0=(t^2+p_2^2)^{1/2}\), changes
\((\Rea b_\phi(p))^2=t^2\cos^2\phi-p_2^2\sin^2\phi\) with nonzero derivative
\(-2p_2\sin^2\phi\). The image of the differential thus contains that
five-dimensional subspace and a vector on which the differential of
\((\Rea z)^2\) does not vanish, so the image has dimension six, the real dimension
of \(\mathcal Q_t\) at \(z\); here \(z\) is a smooth point of
\(\mathcal Q_t\), because \(d(z\cdot z)=2z\cdot dz\ne0\) for \(z\ne0\) and
\(|z|_{\mathrm E}\ge\lvert\Rea\bz\rvert\ge4\) on \(\mathscr Z\). The map \(\Xi_t\) is
therefore a submersion at these preimages. Every preimage of a point of
\(\mathscr Z\) has \(\bp_\perp\ne0\): for every \(p\in H_t^+\),
\((\Rea b_\phi(p))^2=t^2\cos^2\phi-|\bp_\perp|^2\sin^2\phi\); the fold
\(u^2=t^2\cos^2\phi\), where \(\bp_\perp=0\), is avoided on \(\mathscr Z\) by
\(u^2\le-16\).

\smallskip
\emph{Step 3: bounds for the matching transformations and the preimage
energies.} For a Lorentz matrix \(\mathcal L\), the Lorentz
factor \(\gamma_{\mathcal L}\) is the time component of the image of the
time axis, and \(\gamma_{\mathcal L}^2-1\) is the sum of the squared time
components of the three spacelike columns, because \(\mathcal L^T\) is a
Lorentz matrix too, so that the first row of \(\mathcal L\) has Lorentzian
square \(1\). The Euclidean operator norm
of a proper orthochronous \(\mathcal L\), and of its inverse, is
\(\gamma_{\mathcal L}+(\gamma_{\mathcal L}^2-1)^{1/2}\), because
\(\mathcal L=R_1BR_2\) with rotations \(R_1,R_2\) of space and a boost \(B\)
along \(x^1\) of Lorentz factor \(\gamma_{\mathcal L}\), whose Euclidean
singular values are \(\gamma_{\mathcal L}\pm(\gamma_{\mathcal L}^2-1)^{1/2}\),
\(1\) and \(1\).

The target frame can be chosen with a
spacelike complementary vector of zero time component: the Lorentzian
complement of the target pair is two-dimensional and meets the hyperplane
\(\{x^0=0\}\) in a nonzero vector, which is spacelike. The Lorentz
factor of the target frame obeys
\[
 \gamma_{\mathrm{tar}}^2
 =1+\frac{u_0^2}{\rho_0^2}
    +\frac{v_0^2}{t^2+\rho_0^2}\le16,
 \qquad
 \norm{\mathcal L_{\mathrm{tar}}}_{\mathrm E}
 =\gamma_{\mathrm{tar}}
   +\sqrt{\gamma_{\mathrm{tar}}^2-1}
 \le4+\sqrt{15}<8,
\]
since \(|\mathbf u|^2=u_0^2+\rho_0^2\) and
\(|\mathbf v|^2=v_0^2+t^2+\rho_0^2\) give
\(|z|_{\mathrm E}^2=2(u_0^2+v_0^2)+2\rho_0^2+t^2\le256\), whence
\(u_0^2+v_0^2\le128\),
\(\gamma_{\mathrm{tar}}^2-1\le(u_0^2+v_0^2)/\rho_0^2\le128/16=8\) and
\(\gamma_{\mathrm{tar}}^2\le9\le16\).

For the source frame, the Lorentzian complement of the reference pair is
spanned by \(\mathbf e_3=(0,0,0,1)\) and \((p_2,0,p_0\cos\phi,0)\), whose
Lorentzian square is \(p_2^2-p_0^2\cos^2\phi=\rho_0^2\); the frame is
completed by \((p_2,0,p_0\cos\phi,0)/\rho_0\), future-directed since
\(p_2>0\), and by \(\pm\mathbf e_3\). With this completion the time
components of the three spacelike
columns are \(p_0\cos\phi/\rho_0\), \(0\), and \(0\), so that
\[
 \gamma_{\mathrm{src}}^2
 =1+\frac{p_0^2\cos^2\phi}{\rho_0^2}
 =1+\cot^2\phi\left(1+\frac{t^2}{\rho_0^2}\right)
 \le1+\frac{17}{16}\cot^2\phi,\qquad
 \norm{\mathcal L_{\mathrm{src}}}_{\mathrm E}
 \le\frac{\sqrt{17}}{2\sin\phi},
\]
using \(t\le1\), \(\rho_0\ge4\), \(\sqrt{1+\cot^2\phi}=1/{\sin\phi}\), and
\(\sqrt{1+\tfrac{17}{16}\cot^2\phi}+\tfrac{\sqrt{17}}4\lvert\cot\phi\rvert
\le\tfrac{\sqrt{17}}2\sqrt{1+\cot^2\phi}\). With \(a=\lvert\cot\phi\rvert\), squaring
reduces the last inequality to
\(\tfrac{\sqrt{17}}2a(1+\tfrac{17}{16}a^2)^{1/2}\le\tfrac{13}4+\tfrac{17}8a^2\),
which holds because
\((\tfrac{13}4+\tfrac{17}8a^2)^2-\tfrac{17}4a^2(1+\tfrac{17}{16}a^2)
=\tfrac{169+153a^2}{16}>0\).

Consequently the matching transformation satisfies
\[
 \norm{\mathcal L}_{\mathrm E}
 =\lVert\mathcal L^{-1}\rVert_{\mathrm E}
 \le\norm{\mathcal L_{\mathrm{tar}}}_{\mathrm E}
   \norm{\mathcal L_{\mathrm{src}}}_{\mathrm E}
 <\frac{4\sqrt{17}}{\sin\phi}
 <\frac{32}{\sin\phi},
 \qquad
 p_0=\frac{\sqrt{t^2+\rho_0^2}}{\sin\phi}
 \le\frac{\sqrt{257}}{\sin\phi}
 <\frac{32}{\sin\phi}.
\]

\smallskip
\emph{Step 4: the shell family and the parameter family over the
normalized mass.} Let \(\mathcal I\supset[0,1]\) be an open interval. The
set
\[
 \mathcal Q=\{(t,z)\in\mathcal I\times\C^4:z\cdot z=t^2,\ z\ne0\}
\]
is a real-analytic manifold of dimension seven, because
\(d(z\cdot z)=2z\cdot dz\) maps onto \(\C\) at \(z\ne0\). For the same reason
the restriction of \(t\) to \(\mathcal Q\) is a submersion: for every
\(\delta t\), the equation \(2z\cdot\delta z=2t\,\delta t\) has a solution
\(\delta z\). Near each of its points, \(\mathcal Q\) is therefore the
product of an interval of values of \(t\) and a slice, and the slices are
the shells \(\mathcal Q_t\) without \(0\). In this sense, on the compact
normalized region \(\mathscr Z\), which lies in \(\mathcal Q\) because
\(|z|_{\mathrm E}\ge4\) there, \(t=\mu/\Energy=0\) is a regular parameter of
the shell family.

The parameter family \(\mathcal P\) consists of the
triples \((t,\mathcal L,\bp)\) with \(t\in\mathcal I\), \(\mathcal L\) proper
orthochronous and \(\bp\in\R^3\setminus\{0\}\), with
\(p=((t^2+|\bp|^2)^{1/2},\bp)\in H_t^+\); this parametrization is
real-analytic in \((t,\bp)\), also at \(t=0\), and
\(d\mathfrak m_t=d^3\bp/(2(2\pi)^3(t^2+|\bp|^2)^{1/2})\) depends continuously
on \(t\). The map \(\Xi(t,\mathcal L,\bp)=(t,\Xi_t(\mathcal L,p))\) sends
\(\mathcal P\) into \(\mathcal Q\). At the preimage
\((t_*,\mathcal L_*,\bp_*)\) of a point \((t_*,z_*)\in\mathscr Z\) found in
Step 1, the differential of \(\Xi\) is onto: its \(t\)-component is the
identity, and at fixed \(t\) it is the differential of \(\Xi_t\), onto by
Step 2. By the submersion theorem there are coordinates \((t,y)\),
\(y\in\R^6\), on \(\mathcal Q\) near \((t_*,z_*)\) and \((t,y,s)\),
\(s\in\R^3\), on \(\mathcal P\) near \((t_*,\mathcal L_*,\bp_*)\) in which
\(\Xi(t,y,s)=(t,y)\): complete \(t\) to coordinates \((t,y)\) on
\(\mathcal Q\), and \((t,y\circ\Xi)\) by three functions \(s\) to coordinates
on \(\mathcal P\).

\smallskip
\emph{Step 5: local pushforward densities, continuous in \(t\).} The
inequalities of Step 3 are strict, so the point
\((t_*,\mathcal L_*,\bp_*)\) has an open neighbourhood \(N\) whose closure is
compact, lies in the domain of the coordinates \((t,y,s)\) of Step 4, and carries \(p_0<32/{\sin\phi}\) and
\(\norm{\mathcal L^{-1}}_{\mathrm E}<32/{\sin\phi}\). Let \(\chi_N\) be a
continuous function with values in \([0,1]\) and compact support in \(N\),
positive at \((t_*,\mathcal L_*,\bp_*)\). With \(d\mathcal L\) for Haar
measure, the measure \(\chi_N\,d\mathcal L\,d\mathfrak m_t\) on the slice of
\(\mathcal P\) at \(t\) is \(m(t,y,s)\,dy\,ds\), with \(m\) continuous and
compactly supported. Fubini's theorem, which is the coarea formula in these
coordinates, shows that \((\Xi_t)_*(\chi_N\,d\mathcal L\,d\mathfrak m_t)\) is
\(\bigl(\int m(t,y,s)\,ds\bigr)\,dy\). Since \(d\vol_{\mathcal Q_t}\) has a
positive density with respect to \(dy\) that is continuous in \((t,y)\), the pushforward has a
density with respect to \(d\vol_{\mathcal Q_t}\) that is continuous in
\((t,y)\), by dominated convergence, and positive at \((t_*,z_*)\), because
\(m>0\) at \((t_*,\mathcal L_*,\bp_*)\). Hence there are an open
neighbourhood \(O_N\) of \((t_*,z_*)\) in \(\mathcal Q\) and \(\delta_N>0\)
such that this density is at least \(\delta_N\) on \(O_N\), for all values of
\(t\) that \(O_N\) contains, \(t=0\) included.

\smallskip
\emph{Step 6: a finite cover and one Lorentz density, which give
\eqref{eq:pushforward}.} The set \(\mathscr Z\) is compact, so finitely many of the neighbourhoods \(O_N\), say
\(O_{N_1},\dots,O_{N_n}\), cover \(\mathscr Z\). We put
\(\mathcal V=(O_{N_1}\cup\dots\cup O_{N_n})\cap{\{0\le t\le1\}}\). The
projections of the closures of \(N_1,\dots,N_n\) to the Lorentz group form
a compact set \(\mathscr L'\) inside the open set
\(\{\mathcal L:\norm{\mathcal L^{-1}}_{\mathrm E}<32/{\sin\phi}\}\). A smooth
function \(\wt\vartheta\ge0\) on the Lorentz group, positive on
\(\mathscr L'\), whose support \(\mathscr L_c\) is a compact subset of that
open set, gives the density
\(\vartheta_c=\wt\vartheta/{\int\wt\vartheta\,d\mathcal L}\), and
\(\sup_{\mathcal L\in\mathscr L_c}\norm{\mathcal L^{-1}}_{\mathrm E}
<32/{\sin\phi}\).

Since \(p_0<M_c\) on each \(N_i\) and
\(0\le\chi_{N_i}\le1\),
\[
 \vartheta_c(\mathcal L)\mathbf1_{\{p_0\le M_c\}}
 \ge\Bigl(\min_{\mathscr L'}\vartheta_c\Bigr)\chi_{N_i}(t,\mathcal L,\bp),
 \qquad i=1,\dots,n.
\]
Pushforward preserves the order of positive measures, so by Step 5, on
the slice \(\{z:(t,z)\in O_{N_i}\}\),
\[
 (\Xi_t)_*\left(d\vartheta_c\otimes
       \mathbf1_{\{p_0\le M_c\}}\,d\mathfrak m_t\right)
 \ge\Bigl(\min_{\mathscr L'}\vartheta_c\Bigr)\delta_{N_i}\,
 {d\vol_{\mathcal Q_t}}.
\]
These slices cover \(\mathcal V_t\), which gives \eqref{eq:pushforward} with
\(\delta_c=(\min_{\mathscr L'}\vartheta_c)\min_i\delta_{N_i}>0\). All choices
depend only on \(\phi=2\pi c\).

\smallskip
\emph{Step 7: holomorphic shell charts of a common radius.} Let
\((t_*,z_*)\in\mathscr Z\). Since \(\lvert\Rea\bz_*\rvert\ge4\), some spatial
coordinate satisfies \(|z_{*k}|\ge4/\sqrt3\), so
\(\partial_{z_k}(z\cdot z)=-2z_k\ne0\) near \(z_*\). The holomorphic implicit
function theorem, applied to \(z\cdot z=\tau\) with \(\tau=t^2\) as an
additional complex variable, writes \(\mathcal Q_t\) near \(z_*\) as a graph
\(z_k=g_t(w)\) over the other three coordinates \(w\in\C^3\), with \(g_t\)
holomorphic. For all \((t,z)\in\mathcal Q\) near \((t_*,z_*)\), the graph is
defined, with \(|\nabla g_t|\) bounded, on a neighbourhood of the closed
ball of one fixed radius about the point \(w\) formed by those three
coordinates of \(z\). Finitely many such neighbourhoods of points \((t_*,z_*)\) cover the compact
set \(\mathscr Z\), which gives one radius and one bound for
\(|\nabla g_t|\); a smaller radius \(r>0\) also keeps the image of the
chart about every \((t,z)\in\mathscr Z\) inside the slice \(\mathcal V_t\),
because the charts are uniformly Lipschitz and, \(\mathscr Z\) being compact
and \(\mathcal V\) open in the closed shell family, all points of the shell
family within some positive distance of \(\mathscr Z\) lie in \(\mathcal V\).

For
\((t,z)\in\mathscr Z\) the chart about \(z\) is the injective holomorphic map
\(\psi_{t,z}(w')=(w',g_t(w'))\), with the coordinates in the order of \(z\),
on the ball \(B_r(w)\), with \(\psi_{t,z}(w)=z\) and image
\(\mathcal B_z\subset\mathcal V_t\). The coordinate projection \(z\mapsto w\)
is a Euclidean contraction and inverts \(\psi_{t,z}\), so the volume
Jacobian \(J_{t,z}\) of \(\psi_{t,z}\) with respect to Lebesgue measure on
\(B_r(w)\subset\C^3=\R^6\) satisfies \(J_{t,z}\ge1\); the bound for
\(|\nabla g_t|\) gives \(J_{t,z}\le C_{\mathrm{sub}}\) for a constant
\(C_{\mathrm{sub}}\). Hence
\(\vol(B_r)\le\vol(\mathcal B_z)\le C_{\mathrm{sub}}\vol(B_r)\). For a
nonnegative plurisubharmonic \(F\) on \(\C^4\), the function
\(F\circ\psi_{t,z}\) is plurisubharmonic, hence subharmonic, near the
closed ball \(\overline{B_r(w)}\), and
\begin{align*}
 F(z)
 &\overset{\text{submean}}{\le}
   \frac1{\vol(B_r)}\int_{B_r(w)}F\circ\psi_{t,z}
 \overset{J_{t,z}\ge1}{\le}
   \frac1{\vol(B_r)}\int_{\mathcal B_z}F\,d\vol_{\mathcal Q_t}\\
 &\overset{J_{t,z}\le C_{\mathrm{sub}}}{\le}
   \frac{C_{\mathrm{sub}}}{\vol(\mathcal B_z)}
   \int_{\mathcal B_z}F\,{d\vol_{\mathcal Q_t}}.
\end{align*}

Finally, the bound on \(\mathscr L_c\) in Step 6 and
\[
 \ell_c=\frac{64}{\sin\phi}
 >\sqrt2\,\frac{32}{\sin\phi}
\]
give the stated wedge containment.
\end{proof}

\begin{lemma}[Null analytic discs]\label{lem:null-discs}
If \(z\in\mathcal Q_\mu\), \(\Energy\ge\mu\), and
\(|z|_{\mathrm E}\le2\Energy\), there is an analytic disc
\(\zeta\mapsto z(\zeta)\) in \(\mathcal Q_\mu\), \(|\zeta|\le1\), with
\(z(0)=z\), such that every boundary point \(z(\zeta)\), \(|\zeta|=1\),
satisfies
\[
 \lvert\Rea\bz(\zeta)\rvert\ge6\Energy,\qquad
 \lvert\Rea z_0(\zeta)\rvert\le2\Energy,\qquad
 (\Rea z(\zeta))^2\le-32\Energy^2,\qquad
 |z(\zeta)|_{\mathrm E}<14\Energy.
\]
In particular, \((\mu/\Energy,z(\zeta)/\Energy)\) belongs to the set
\(\mathscr Z\) of \eqref{eq:compact-target} for \(|\zeta|=1\).
\end{lemma}

\begin{proof}
There is a nonzero \(\mathbf n\in\C^3\) such that
\[
 \bz\cdot\mathbf n=0,\qquad \mathbf n\cdot\mathbf n=0.
\]
The first relation defines a complex subspace of dimension at least
two, and a quadratic form on a two-dimensional complex space has a
nonzero isotropic vector. Since
\[
 \mathbf n\cdot\mathbf n
 =\lvert\Rea\mathbf n\rvert^2-\lvert\Ima\mathbf n\rvert^2+2i\Rea\mathbf n\cdot\Ima\mathbf n
 =0,
\]
the real and imaginary parts of \(\mathbf n\) are orthogonal and of equal
positive length. After rescaling by a positive factor,
\[
 \lvert\Rea\mathbf n\rvert=\lvert\Ima\mathbf n\rvert=1,\qquad
 \Rea\mathbf n\cdot\Ima\mathbf n=0.
\]
Hence \(\lvert\Rea(e^{iu}\mathbf n)\rvert^2=\lvert\cos u\Rea\mathbf n-\sin u\Ima\mathbf n\rvert^2
=\cos^2u+\sin^2u=1\) for every real \(u\).

The disc
\[
 z(\zeta)=z+8\Energy\zeta(0,\mathbf n),\qquad |\zeta|\le1,
\]
lies on \(\mathcal Q_\mu\), because
\[
 z(\zeta)\cdot z(\zeta)
 =z\cdot z-16\Energy\zeta\,\bz\cdot\mathbf n
 -64\Energy^2\zeta^2\,\mathbf n\cdot\mathbf n
 =\mu^2.
\]
On its boundary \(\zeta=e^{iu}\), the time component is unchanged,
\(z_0(\zeta)=z_0\), and \(\lvert\Rea\bz\rvert\), \(\lvert\Rea z_0\rvert\) are at most
\(|z|_{\mathrm E}\le2\Energy\); hence
\[
 \lvert\Rea\bz(\zeta)\rvert\ge8\Energy\lvert\Rea(e^{iu}\mathbf n)\rvert-\lvert\Rea\bz\rvert
 \ge8\Energy-2\Energy=6\Energy,\qquad
 \lvert\Rea z_0(\zeta)\rvert\le2\Energy.
\]
These two inequalities imply the Lorentzian bound
\[
 (\Rea z(\zeta))^2=(\Rea z_0(\zeta))^2-\lvert\Rea\bz(\zeta)\rvert^2
 \le4\Energy^2-36\Energy^2=-32\Energy^2,
\]
while
\(|(0,\mathbf n)|_{\mathrm E}=\sqrt2\) and \(8\sqrt2<12\) give
\[
 |z(\zeta)|_{\mathrm E}
 \le|z|_{\mathrm E}+8\Energy|(0,\mathbf n)|_{\mathrm E}
 \le(2+8\sqrt2)\Energy<14\Energy.
\]
The normalized boundary points lie in the set \(\mathscr Z\) of
\eqref{eq:compact-target}: \(t=\mu/\Energy\in(0,1]\) because \(\Energy\ge\mu\), and
\(w=z(\zeta)/\Energy\) satisfies \(w\cdot w=t^2\),
\(|w|_{\mathrm E}<14\le16\) and \((\Rea w)^2\le-32\le-16\).
\end{proof}

\begin{proof}[Proof of Theorem~\ref{thm:evaluation}]
\emph{Step 1: the bound at the normalized points of \(\mathscr Z\), for
smooth sources.} Let \(\sx\) be smooth with compact support in \(U\), and
write \(\mu=t\Energy\), so that \(t\in(0,1]\). The substitution
\(p=\Energy p'\) maps \(H_t^+\) onto \(H_\mu^+\), and since
\(p_0=\Energy p_0'\) and \(d^3\bp=\Energy^3d^3\bp'\), the invariant measure
scales as
\[
 d\mathfrak m_\mu(\Energy p')=\Energy^2d\mathfrak m_t(p').
\]
For \((t,z)\in\mathscr Z\), let \(\mathcal B_z\subset\mathcal V_t\) be the
image of the holomorphic shell chart about \(z\) of
Lemma~\ref{lem:regular-patches}. The function
\(F=|\Psi_\sx(\Energy\,\cdot)|^2\) is plurisubharmonic on \(\C^4\), as the
squared modulus of an entire function. Then
\begin{align*}
 |\Psi_\sx(\Energy z)|^2
 &\overset{\text{Lem.~\ref{lem:regular-patches}}}{\le}
   \frac{C_{\mathrm{sub}}}{\vol(\mathcal B_z)}
   \int_{\mathcal B_z}|\Psi_\sx(\Energy z')|^2\,d\vol_{\mathcal Q_t}(z')\\
 &\overset{\eqref{eq:pushforward}}{\le}
   \frac{C_{\mathrm{sub}}}{\delta_c\vol(\mathcal B_z)}
   \int_{\mathscr L_c}\int_{p_0'\le M_c}
   |\Psi_\sx(\Energy\mathcal Lb_\phi(p'))|^2
   \,d\mathfrak m_t(p')\,d\vartheta_c(\mathcal L)\\
 &\overset{p=\Energy p'}{=}
   \frac{C_{\mathrm{sub}}\Energy^{-2}}{\delta_c\vol(\mathcal B_z)}
   \int_{\mathscr L_c}\int_{p_0\le\Energy M_c}
   |\Psi_\sx(\mathcal Lb_\phi(p))|^2
   \,d\mathfrak m_\mu(p)\,d\vartheta_c(\mathcal L)\\
 &\overset{p_0\le\Energy M_c}{\le}
   \frac{C_{\mathrm{sub}}\Energy^{-2}e^{2\ell_cM_cT\Energy\sin\phi}}
   {\delta_c\vol(\mathcal B_z)}
   \int_{\mathscr L_c}\int_{H_\mu^+}
   e^{-2\ell_cTp_0\sin\phi}
   |\Psi_\sx(\mathcal Lb_\phi(p))|^2
   \,d\mathfrak m_\mu(p)\,d\vartheta_c(\mathcal L)\\
 &\overset{\eqref{eq:averaged-wedges}}{\le}
   C_c\Energy^{-2}e^{2\ell_cM_cT\Energy\sin\phi}\norm{h_{c,U}\sx}^2.
\end{align*}
The second inequality applies \eqref{eq:pushforward} on
\(\mathcal B_z\subset\mathcal V_t\) to the function equal to \(F\) on
\(\mathcal B_z\) and to \(0\) elsewhere, which is at most \(F\) because \(F\ge0\):
\[
 \int_{\mathcal B_z}F\,d\vol_{\mathcal Q_t}
 \le\frac1{\delta_c}\int_{\mathscr L_c}\int_{H_t^+}
 \mathbf1_{\{p_0'\le M_c\}}F(\mathcal Lb_\phi(p'))
 \,d\mathfrak m_t(p')\,d\vartheta_c(\mathcal L),
\]
and \(F(\mathcal Lb_\phi(p'))=|\Psi_\sx(\Energy\mathcal Lb_\phi(p'))|^2\). The
equality substitutes \(p=\Energy p'\), with
\(\Energy\mathcal Lb_\phi(p')=\mathcal Lb_\phi(\Energy p')\) because
\(\mathcal L\) and \(b_\phi\) are linear. The third inequality inserts the
factor \(e^{2\ell_cT(\Energy M_c-p_0)\sin\phi}\ge1\) on
\(\{p_0\le\Energy M_c\}\) and then integrates over all of \(H_\mu^+\). The
last inequality is \eqref{eq:averaged-wedges} for
\(\mathscr L=\mathscr L_c\), \(\vartheta=\vartheta_c\) and
\(\ell_{\mathscr L}=\ell_c\), which Lemma~\ref{lem:regular-patches} permits,
with
\[
 C_c=\max\Bigl\{1,\frac{C_{\mathrm{sub}}}
 {\delta_c\inf_{(t,z)\in\mathscr Z}\vol(\mathcal B_z)}\Bigr\},
\]
which is finite because \(\vol(\mathcal B_z)\ge\vol(B_r)>0\) for every
\((t,z)\in\mathscr Z\) (Lemma~\ref{lem:regular-patches}).

The values of \(\ell_c\) and \(M_c\) in Lemma~\ref{lem:regular-patches} give
\[
 \ell_cM_c\sin\phi
 =\frac{64\cdot32}{\sin\phi}
 =\frac{2048}{\sin(2\pi c)}=\lambda_c,
\]
so that
\(|\Psi_\sx(\Energy z)|^2\le C_c\Energy^{-2}e^{2\lambda_cT\Energy}
\norm{h_{c,U}\sx}^2\) for every \((t,z)\in\mathscr Z\).

\smallskip
\emph{Step 2: every point of the statement, through analytic discs.} Let
\(z\in\mathcal Q_\mu\) with \(|z|_{\mathrm E}\le2\Energy\), and let
\(\zeta\mapsto z(\zeta)\) be the analytic disc of
Lemma~\ref{lem:null-discs}. The function \(\zeta\mapsto\Psi_\sx(z(\zeta))\)
is holomorphic near the closed unit disc, so Cauchy's mean formula gives
\[
 |\Psi_\sx(z)|^2
 =\Bigl|\frac1{2\pi}\int_0^{2\pi}\Psi_\sx(z(e^{iu}))\,du\Bigr|^2
 \le\max_{0\le u\le2\pi}|\Psi_\sx(z(e^{iu}))|^2.
\]
By Lemma~\ref{lem:null-discs}, \(z(e^{iu})=\Energy w_u\) with
\((\mu/\Energy,w_u)\in\mathscr Z\), so Step 1, at the same \(\Energy\) and
\(t=\mu/\Energy\), bounds every \(|\Psi_\sx(z(e^{iu}))|^2\), and hence \(|\Psi_\sx(z)|^2\), by
\(C_c\Energy^{-2}e^{2\lambda_cT\Energy}\norm{h_{c,U}\sx}^2\), which is
\eqref{eq:evaluation} for smooth sources, with the same constant.

\smallskip
\emph{Step 3: all sources, by density.} For each \(z\), the map
\(\sx\mapsto\Psi_\sx(z)\) is continuous on \(\Hh_{\mu,U}\) (beginning of
Section~\ref{sec:evaluation}), \(h_{c,U}\) is bounded, and smooth sources
are dense in \(\Hh_{\mu,U}\) by its definition. Both sides of
\eqref{eq:evaluation} are therefore continuous in \(\sx\), and density
extends the estimate from smooth sources to \(\Hh_{\mu,U}\).
The set \(U\) enters only through \(U\subset B_T(0)\) in
\eqref{eq:averaged-wedges}, and \(C_c\), \(\lambda_c\) come from
Lemma~\ref{lem:regular-patches} and \(\ell_c\) alone.
\end{proof}

\subsection{Weighted polynomial moments}

For \(1/2\le b<1\), we fix
\begin{equation}
 \Lambda_b=\frac{2^{15}}{\sin(\pi b)}
 =16\max\{1,\lambda_{(1-b)/2},\lambda_{b/2}\},
 \qquad B_b=5\Lambda_b.
 \label{eq:Lambda}
\end{equation}
The two expressions for \(\Lambda_b\) agree because, for the admissible
choice of \(\lambda_c\) in Theorem~\ref{thm:evaluation},
\(\lambda_{(1-b)/2}=\lambda_{b/2}=2048/{\sin(\pi b)}\).

Theorem~\ref{thm:evaluation} applies to sources in a ball
\(B_T(\bx_0)\) after the spatial translation by \(-\bx_0\), which
preserves the symplectic form, the covariance, and hence
\(h_{c,U}\).
For a multi-index \(\alpha\in\N^3\), define
\[
 \ell_\alpha^f(f,g)=\int\bx^\alpha f(\bx)\,d^3x,\qquad
 \ell_\alpha^g(f,g)=\int\bx^\alpha g(\bx)\,d^3x.
\]
We also write
\(\ell_{\bz}^f(f,g)=\wh f(\bz)\) and
\(\ell_{\bz}^g(f,g)=\wh g(\bz)\).

\begin{lemma}[Weighted moments]\label{lem:moments}
Let \(1/2\le b<1\), with \(\Lambda_b,B_b\) as in \eqref{eq:Lambda}.
There is a finite \(M_b\), depending only on \(b\), such that for all
\(\mu,T>0\), \(\bx_0\in\R^3\), every nonempty bounded open
\(U\subset B_T(\bx_0)\) and every \(\alpha\in\N^3\), the moments
\(\ell_\alpha^f,\ell_\alpha^g\), taken in coordinates centred at
\(\bx_0\), satisfy, with the uncertainty reserve \(\Res_U\) of
\eqref{eq:reserve} and in the sense of
Convention~\ref{conv:factorization},
\begin{align}
 \norm{\ell_\alpha^f\Res_U^{-b}}
 &\le M_b\frac{T}{1+\mu T}
      e^{\Lambda_b\mu T}(B_bT)^{|\alpha|},
 \label{eq:f-moments}\\
 \norm{\ell_\alpha^g\Res_U^{-b}}
 &\le M_b\frac{T^2}{(1+\mu T)^2}
      e^{\Lambda_b\mu T}(B_bT)^{|\alpha|}.
 \label{eq:g-moments}
\end{align}
\end{lemma}

\begin{proof}
\emph{Step 1: evaluation on the complex shell, bounded by the reserve.}
Throughout the proof we use coordinates centred at \(\bx_0\). The spectral measure of
\(G_U\) is split at
\(1\): \(P_1\) denotes the spectral projection of \(G_U\) for
\([0,1]\), and \(P_2=I-P_1\). By \eqref{eq:modular-source} and
\eqref{eq:reserve}, in the joint spectral calculus of the commuting
operators \(G_U,\overline G_U\),
\(h_{c,U}=G_U^{1/2-c}\overline G_U^{\,c}\) and
\(\Res_U^b=G_U^{b/2}\overline G_U^{\,b/2}\).

On the range of \(P_1\), where \(G_U\le I\) and
\(\overline G_U=2I-G_U\ge I\), we choose \(c=(1-b)/2\); on the range
of \(P_2\), where \(G_U\ge I\), we choose \(c=b/2\). Then
\begin{align*}
 h_{(1-b)/2,U}P_1
 &=G_U^{b/2}\,\overline G_U^{\,(1-b)/2}P_1
 =\overline G_U^{\,1/2-b}P_1\,\Res_U^b,\\
 h_{b/2,U}P_2
 &=G_U^{(1-b)/2}\,\overline G_U^{\,b/2}P_2
 =G_U^{1/2-b}P_2\,\Res_U^b,
\end{align*}
and the operators \(\overline G_U^{\,1/2-b}P_1\) and
\(G_U^{1/2-b}P_2\) are contractions because \(1/2-b\le0\). Both
values of \(c\) lie in \((0,1/2)\), as Theorem~\ref{thm:evaluation}
requires, because \(1/2\le b<1\).

Let \(\sx\in\Hh_{\mu,U}\), \(\Energy\ge\mu\), and
\(z\in\mathcal Q_\mu\) with \(|z|_{\mathrm E}\le2\Energy\). Since
\(\sx=P_1\sx+P_2\sx\) and \(\Psi_\sx(z)\) is linear in \(\sx\),
Theorem~\ref{thm:evaluation} at \(c=(1-b)/2\) for \(P_1\sx\) and at
\(c=b/2\) for \(P_2\sx\) gives
\begin{align*}
 |\Psi_\sx(z)|
 &\le|\Psi_{P_1\sx}(z)|+|\Psi_{P_2\sx}(z)|\\
 &\overset{\eqref{eq:evaluation}}{\le}
 \Energy^{-1}\Bigl(C_{(1-b)/2}^{1/2}e^{\lambda_{(1-b)/2}T\Energy}
   \norm{h_{(1-b)/2,U}P_1\sx}
   +C_{b/2}^{1/2}e^{\lambda_{b/2}T\Energy}
   \norm{h_{b/2,U}P_2\sx}\Bigr)\\
 &\le\bigl(C_{(1-b)/2}^{1/2}+C_{b/2}^{1/2}\bigr)\Energy^{-1}
   e^{\max\{\lambda_{(1-b)/2},\lambda_{b/2}\}T\Energy}
   \norm{\Res_U^b\sx},
\end{align*}
where the last step uses
\(\lVert h_{(1-b)/2,U}P_1\sx\rVert\le\lVert\Res_U^b\sx\rVert\) and
\(\lVert h_{b/2,U}P_2\sx\rVert\le\lVert\Res_U^b\sx\rVert\), which follow from the
two contractions. This bound reads, for every
\(\Energy\ge\mu\),
\begin{equation}
 \norm{\Psi_{(\cdot)}(z)\Res_U^{-b}}
 \le D_b^{\mathrm{ev}}\Energy^{-1}
       e^{\lambda_b^{\mathrm{ev}}T\Energy},
 \qquad
 z\in\mathcal Q_\mu,\quad |z|_{\mathrm E}\le2\Energy,
 \label{eq:symmetric-evaluation}
\end{equation}
where
\(\lambda_b^{\mathrm{ev}}
=\max\{\lambda_{(1-b)/2},\lambda_{b/2}\}\)
and \(D_b^{\mathrm{ev}}=C_{(1-b)/2}^{1/2}+C_{b/2}^{1/2}\) is finite.

\smallskip
\emph{Step 2: the bound on \(\wh f\).}
For \(\bz\in\C^3\), with Euclidean norm \(|\bz|\), we take
\begin{equation}
 \Energy=\mu+\frac1{\Lambda_bT}+|\bz|.
 \label{eq:energy-choice}
\end{equation}
Then \(\Energy\ge\mu\) and \(\mu+|\bz|<\Energy\). Let
\(\sx=(f,g)\in\Hh_{\mu,U}\). The two shell points
\(z_\pm=(\pm\sqrt{\mu^2+\bz\cdot\bz},\bz)\) lie on \(\mathcal Q_\mu\),
and \(\Psi_\sx(z_\pm)=\wh f(\bz)\pm i\sqrt{\mu^2+\bz\cdot\bz}\,\wh g(\bz)\).
Hence
\[
 \wh f(\bz)=\frac{\Psi_\sx(z_+)+\Psi_\sx(z_-)}2,\qquad
 \Psi_\sx(z_+)-\Psi_\sx(z_-)=2i\sqrt{\mu^2+\bz\cdot\bz}\,\wh g(\bz).
\]
Since \(|\bz\cdot\bz|\le|\bz|^2\),
\[
 |z_\pm|_{\mathrm E}^2
 =|\mu^2+\bz\cdot\bz|+|\bz|^2
 \le\mu^2+2|\bz|^2
 \le2(\mu+|\bz|)^2
 <2\Energy^2,
\]
so \(|z_\pm|_{\mathrm E}<\sqrt2\,\Energy\le2\Energy\), and
\eqref{eq:symmetric-evaluation} at the energy \(\Energy\) applies at
\(z_\pm\):
\[
 |\wh f(\bz)|
 \le\frac{|\Psi_\sx(z_+)|+|\Psi_\sx(z_-)|}2
 \overset{\eqref{eq:symmetric-evaluation}}{\le}
 D_b^{\mathrm{ev}}\Energy^{-1}e^{\lambda_b^{\mathrm{ev}}T\Energy}
 \norm{\Res_U^b\sx}.
\]

\smallskip
\emph{Step 3: the bound on \(\wh g\) across zero energy.}
The difference in Step 2 determines \(\wh g(\bz)\) only at nonzero
complex energy \(\sqrt{\mu^2+\bz\cdot\bz}\). To obtain a bound across
zero energy, \(z_1\) is varied on the circle \(z_1+8\Energy e^{iu}\),
\(u\in\R\), with \(\Energy\) fixed by \eqref{eq:energy-choice} at
\(\bz\). The point \(\bz_u=(z_1+8\Energy e^{iu},z_2,z_3)\) satisfies
\[
 \mu^2+\bz_u\cdot\bz_u
 =64\Energy^2e^{2iu}+16\Energy e^{iu}z_1+(\mu^2+\bz\cdot\bz),
\]
where \(|z_1|\le|\bz|<\Energy\) and
\(|\mu^2+\bz\cdot\bz|\le\mu^2+|\bz|^2<\Energy^2\). On this circle,
therefore,
\begin{align*}
 \left|\mu^2+(z_1+8\Energy e^{iu})^2+z_2^2+z_3^2\right|
 &\ge64\Energy^2-16\Energy|z_1|-|\mu^2+\bz\cdot\bz|\\
 &\ge64\Energy^2-16\Energy^2-\Energy^2=47\Energy^2.
\end{align*}

The two shell points
\(w_\pm=(\pm\sqrt{\mu^2+\bz_u\cdot\bz_u},\bz_u)\) over \(\bz_u\)
satisfy, by the bound on \(|z_\pm|_{\mathrm E}\) in Step 2 and \(|\bz_u|\le|\bz|+8\Energy\),
\begin{align*}
 |w_\pm|_{\mathrm E}^2
 &\le\mu^2+2|\bz_u|^2
 \le2(\mu+|\bz_u|)^2
 \le2(\mu+|\bz|+8\Energy)^2\\
 &<2(9\Energy)^2=162\Energy^2<(14\Energy)^2,
\end{align*}
so the two shell points have Euclidean norm less than
\(14\Energy=2\cdot7\Energy\). Since also \(7\Energy\ge\mu\),
\eqref{eq:symmetric-evaluation} applies at the energy \(7\Energy\).
Dividing the difference of Step 2 at \(\bz_u\) by
\(2i\sqrt{\mu^2+\bz_u\cdot\bz_u}\), whose modulus is at least
\(2\sqrt{47}\,\Energy\), gives
\begin{align*}
 |\wh g(\bz_u)|
 &=\frac{|\Psi_\sx(w_+)-\Psi_\sx(w_-)|}
        {2\,|\mu^2+\bz_u\cdot\bz_u|^{1/2}}
 \overset{\eqref{eq:symmetric-evaluation}}{\le}
 \frac{2D_b^{\mathrm{ev}}(7\Energy)^{-1}
       e^{7\lambda_b^{\mathrm{ev}}T\Energy}}
      {2\sqrt{47}\,\Energy}
 \norm{\Res_U^b\sx}\\
 &=\frac{D_b^{\mathrm{ev}}}{7\sqrt{47}}\,\Energy^{-2}
   e^{7\lambda_b^{\mathrm{ev}}T\Energy}\norm{\Res_U^b\sx}.
\end{align*}
The function \(\zeta\mapsto\wh g(\zeta,z_2,z_3)\) is entire, so the
Cauchy mean
\[
 \wh g(\bz)=\frac1{2\pi}\int_0^{2\pi}\wh g(\bz_u)\,du
\]
obeys the same bound.

By \eqref{eq:Lambda}, \(\Lambda_b\ge16\lambda_b^{\mathrm{ev}}\), so the
factor \(16\) covers both exponents:
\(\lambda_b^{\mathrm{ev}}T\Energy\le7\lambda_b^{\mathrm{ev}}T\Energy
\le\Lambda_bT\Energy\). Since also \(7\sqrt{47}>1\),
Steps 2 and 3 yield, with \(D_b'=D_b^{\mathrm{ev}}\),
\[
 \norm{\ell_{\bz}^f\Res_U^{-b}}
 \le D_b'\Energy^{-1}e^{\Lambda_bT\Energy},\qquad
 \norm{\ell_{\bz}^g\Res_U^{-b}}
 \le D_b'\Energy^{-2}e^{\Lambda_bT\Energy}.
\]

\smallskip
\emph{Step 4: the prefactors and the growth in \(\bz\).}
Formula \eqref{eq:energy-choice} gives
\(\Lambda_bT\Energy=\Lambda_b\mu T+1+\Lambda_bT|\bz|\ge1+\Lambda_b\mu T\).
With \(\Lambda_b\ge1\), it follows that
\[
 \Energy^{-1}\le\frac{\Lambda_bT}{1+\Lambda_b\mu T}
 \le\frac{\Lambda_bT}{1+\mu T},\qquad
 e^{\Lambda_bT\Energy}=e\,e^{\Lambda_b\mu T}e^{\Lambda_bT|\bz|}.
\]
Inserted in the two bounds of Step 3, these two relations produce the
dimensional prefactors and the growth \(e^{\Lambda_bT|\bz|}\) in \(\bz\):
for every \(\bz\in\C^3\),
\begin{align*}
 \norm{\ell_{\bz}^f\Res_U^{-b}}
 &\le eD_b'\Lambda_b\,\frac{T}{1+\mu T}\,
   e^{\Lambda_b\mu T}e^{\Lambda_bT|\bz|},\\
 \norm{\ell_{\bz}^g\Res_U^{-b}}
 &\le eD_b'\Lambda_b^2\,\frac{T^2}{(1+\mu T)^2}\,
   e^{\Lambda_b\mu T}e^{\Lambda_bT|\bz|}.
\end{align*}

\smallskip
\emph{Step 5: the moments by Cauchy estimates.}
Differentiation of \(\wh f(\bz)=\int e^{-i\bz\cdot\bx}f(\bx)\,d^3x\)
and of \(\wh g(\bz)\) at \(\bz=0\) gives, in the centred coordinates,
\[
 \partial^\alpha\wh f(0)=(-i)^{|\alpha|}\ell_\alpha^f(f,g),\qquad
 \partial^\alpha\wh g(0)=(-i)^{|\alpha|}\ell_\alpha^g(f,g).
\]
By
Convention~\ref{conv:factorization}, Step 4 bounds the entire
functions \(\wh f\) and \(\wh g\) on all of \(\C^3\) by multiples of
\(e^{\Lambda_bT|\bz|}\norm{\Res_U^b\sx}\).

For \(n=|\alpha|\ge1\), let \(r=n/(\sqrt3\Lambda_bT)\). On the closed
polydisc \(\{|z_1|,|z_2|,|z_3|\le r\}\) one has \(|\bz|\le\sqrt3\,r\),
hence \(e^{\Lambda_bT|\bz|}\le e^n\). The multivariable Cauchy
estimate on this polydisc gives
\begin{align*}
 |\ell_\alpha^f(f,g)|
 &\overset{\phantom{\text{Step 4}}}{=}|\partial^\alpha\wh f(0)|
 \le\frac{\alpha!}{r^n}\max_{|z_1|,|z_2|,|z_3|\le r}|\wh f(\bz)|\\
 &\overset{\text{Step 4}}{\le}
 \alpha!\left(\frac{\sqrt3\Lambda_bT}{n}\right)^n e^n
 \cdot eD_b'\Lambda_b\,\frac{T}{1+\mu T}\,e^{\Lambda_b\mu T}
 \norm{\Res_U^b\sx}.
\end{align*}
The derivative factor is bounded by
\[
 \alpha!\left(\frac{\sqrt3\Lambda_bT}{n}\right)^n e^n
 \le n^n\left(\frac{\sqrt3\Lambda_bT}{n}\right)^n e^n
 =(e\sqrt3\,\Lambda_bT)^n
 \le(5\Lambda_bT)^n=(B_bT)^n,
\]
using \(\alpha!\le n!\le n^n\) and \(e\sqrt3<5\).

The same estimates
for \(\wh g\), with the second bound of Step 4, give the factor
\(eD_b'\Lambda_b^2\,T^2/(1+\mu T)^2\) in place of
\(eD_b'\Lambda_b\,T/(1+\mu T)\). The remaining constants are absorbed
in \(M_b=eD_b'\Lambda_b^2\), which depends only on \(b\) and, since
\(\Lambda_b\ge1\), is at least \(eD_b'\Lambda_b\); with this \(M_b\),
the resulting bounds on \(\ell_\alpha^f\) and \(\ell_\alpha^g\) are \eqref{eq:f-moments} and
\eqref{eq:g-moments} for \(n\ge1\). The case
\(n=0\) follows directly from Step 4 at \(\bz=0\), with \(\Energy\) from
\eqref{eq:energy-choice}: there \(\ell_\alpha^f(f,g)=\wh f(0)\),
\(\ell_\alpha^g(f,g)=\wh g(0)\), and \(e^{\Lambda_bT|\bz|}=1\).
\end{proof}

\section{The forward estimate for separated balls}\label{sec:forward}

For unit-radius balls at separation \(s>2\) and mass \(\mu>0\), we set
\(\xi=\mu s\). Throughout this section, \(\Res_A\) and \(\Res_B\)
denote the uncertainty reserves \eqref{eq:reserve} of the two balls
\(B_1(0)\) and \(B_1(s\eone)\). The real cross-covariance operator in
the separate marginal source norms,
\[
 C_{\mu,s}:\Hh_{\mu,B_1(s\eone)}
 \longrightarrow\Hh_{\mu,B_1(0)},
\]
is defined on real sources by
\[
 \ip{\sx}{C_{\mu,s}\sy}_{\vv_{\mu,B_1(0)}}
 =\vv_\mu(\sx,\sy),
 \qquad
 \supp\sx\subset B_1(0),\quad \supp\sy\subset B_1(s\eone),
\]
where \(\vv_\mu(\sx,\sy)\) is the vacuum covariance form
\eqref{eq:covariance} of the two sources; by \eqref{eq:gram}, this
form is bounded by the product of the two marginal norms, and
\(C_{\mu,s}\) is extended complex-linearly.

Away from the diagonal, the field and momentum covariance kernels are
\begin{equation}
 \mathfrak X_\mu(\mathbf w)
 =\frac{\mu}{4\pi^2\varrho}K_1(\mu\varrho),\qquad
 \mathfrak P_\mu(\mathbf w)
 =-\frac{\mu^2}{4\pi^2\varrho^2}K_2(\mu\varrho),
 \qquad \varrho=\sqrt{\mathbf w\cdot\mathbf w},
 \label{eq:kernels}
\end{equation}
so that, for \(\sx=(f,g)\) and \(\sy=(f',g')\) with disjoint supports,
\[
 \vv_\mu(\sx,\sy)=\iint\bigl[f(\bx)\mathfrak X_\mu(\by-\bx)f'(\by)
 +g(\bx)\mathfrak P_\mu(\by-\bx)g'(\by)\bigr]\,d^3x\,d^3y,
\]
since the real covariance has zero mixed field--momentum kernel. The
square root \(\varrho\) has positive real part on the complex separated
neighbourhoods used below.

For real \(\mathbf w\ne0\), where \(\varrho=|\mathbf w|\), the
heat-kernel integral for \(h_\mu^{-1}/2\) gives the first formula in
\eqref{eq:kernels}. With
\(h_\mu^{-1}=\pi^{-1/2}\int_0^\infty t^{-1/2}e^{-th_\mu^2}\,dt\) and the
kernel \((4\pi t)^{-3/2}e^{-\varrho^2/(4t)-\mu^2t}\) of
\(e^{-th_\mu^2}\), the convolution kernel of \(h_\mu^{-1}/2\) is
\[
 \frac1{2\sqrt\pi}\int_0^\infty
 \frac{e^{-\varrho^2/(4t)-\mu^2t}}{t^{1/2}(4\pi t)^{3/2}}\,dt
 =\frac1{16\pi^2}\int_0^\infty t^{-2}e^{-\varrho^2/(4t)-\mu^2t}\,dt
 =\frac{\mu}{4\pi^2\varrho}K_1(\mu\varrho).
\]
The last equality is the representation
\(K_1(x)=(x/4)\int_0^\infty t^{-2}e^{-t-x^2/(4t)}\,dt\)
\cite[Eq.~10.32.10]{DLMF} at \(x=\mu\varrho\), after the substitution
\(t\mapsto t/\mu^2\).

Applying \(-\nabla^2+\mu^2\) to \(\mathfrak X_\mu\) gives the second formula. Off the
diagonal, the kernel of \(h_\mu/2=(-\nabla^2+\mu^2)\,h_\mu^{-1}/2\) is
\((-\nabla^2+\mu^2)\mathfrak X_\mu\). On radial functions
\(\nabla^2=\varrho^{-1}\partial_\varrho^2\,\varrho\), and
\(\varrho\,\mathfrak X_\mu=\mu K_1(x)/(4\pi^2)\) with \(x=\mu\varrho\).
Hence
\[
 \begin{split}
 (-\nabla^2+\mu^2)\mathfrak X_\mu
 &=-\frac{\partial_\varrho^2(\varrho\,\mathfrak X_\mu)}{\varrho}
 +\mu^2\mathfrak X_\mu
 =\frac{\mu^4}{4\pi^2x}\bigl(K_1(x)-K_1''(x)\bigr)\\
 &=-\frac{\mu^4}{4\pi^2x^2}K_2(x)
 =\mathfrak P_\mu(\mathbf w),
 \end{split}
\]
because \(K_1-K_1''=-x^{-1}K_2\) by the modified Bessel equation
\(K_1''=(1+x^{-2})K_1-x^{-1}K_1'\) and the recurrence
\(K_1'=-K_2+x^{-1}K_1\).

The integral representation \cite[Eq.~10.32.9]{DLMF}
\[
 K_j(x)=\int_0^\infty e^{-x\cosh u}\cosh(ju)\,du
\]
gives, for \(x>0\),
\[
 \begin{split}
 K_1(x)&\le e^{-x}
 \left(x^{-1}+\sqrt{\frac{\pi}{2x}}\right),\\
 K_2(x)&\le e^{-x}
 \left(2x^{-2}+4x^{-1}+\sqrt{\frac{\pi}{2x}}\right).
 \end{split}
\]
Indeed, the substitution \(v=\cosh u-1\), with
\(dv=\sqrt{v(v+2)}\,du\), \(\cosh u=1+v\) and \(\cosh2u=2v^2+4v+1\),
gives
\[
 K_1(x)=e^{-x}\int_0^\infty e^{-xv}\,\frac{1+v}{\sqrt{v(v+2)}}\,dv,
 \qquad
 K_2(x)=e^{-x}\int_0^\infty e^{-xv}\,\frac{2v^2+4v+1}{\sqrt{v(v+2)}}\,dv.
\]
The bound \(\sqrt{v(v+2)}\ge v\) on the terms of the numerators that
contain \(v\), and the bound \(\sqrt{v(v+2)}\ge\sqrt{2v}\) on their
constant terms, give
\[
 \frac{1+v}{\sqrt{v(v+2)}}\le1+(2v)^{-1/2},
 \qquad
 \frac{2v^2+4v+1}{\sqrt{v(v+2)}}\le2v+4+(2v)^{-1/2}.
\]
Integration against \(e^{-xv}\), with
\(\int_0^\infty e^{-xv}\,dv=x^{-1}\),
\(\int_0^\infty v\,e^{-xv}\,dv=x^{-2}\) and
\(\int_0^\infty(2v)^{-1/2}e^{-xv}\,dv=\sqrt{\pi/(2x)}\), yields the
bounds for \(K_1\) and \(K_2\) stated above.

Multiplying the bound for \(K_1\) by \(x\) and the bound for \(K_2\) by
\(x^2\) gives
\[
 \begin{split}
 xK_1(x)&\le\Bigl(1+\sqrt{\tfrac\pi2}\,x^{1/2}\Bigr)e^{-x}
 \le\Bigl(1+\sqrt{\tfrac\pi2}\Bigr)(1+x)^{1/2}e^{-x},\\
 x^2K_2(x)&\le\Bigl(2+4x+\sqrt{\tfrac\pi2}\,x^{3/2}\Bigr)e^{-x}
 \le\Bigl(6+\sqrt{\tfrac\pi2}\Bigr)(1+x)^{3/2}e^{-x},
 \end{split}
\]
because \(1,x^{1/2}\le(1+x)^{1/2}\) and \(1,x,x^{3/2}\le(1+x)^{3/2}\).
Since \(1+\sqrt{\pi/2}<3\) and \(6+\sqrt{\pi/2}<8\), in particular,
\begin{equation}
 xK_1(x)\le3(1+x)^{1/2}e^{-x},\qquad
 x^2K_2(x)\le8(1+x)^{3/2}e^{-x}.
 \label{eq:Bessel-scaled}
\end{equation}
For \(\Rea z>0\), the integral representation of \(K_j\) gives
\(|K_j(z)|\le K_j(\Rea z)\).

\begin{theorem}[Marginally weighted trace estimate]\label{thm:forward}
For each fixed \(1/2\le b<1\), there are finite \(S_b,F_b\),
independent of \(\mu\), such that for every \(\mu>0\) and \(s\ge S_b\)
the cross-covariance operator \(C_{\mu,s}\) factors through the reserves
\(\Res_A,\Res_B\) of the two balls as
\begin{equation}
 C_{\mu,s}=\Res_A^bX_{b,\mu,s}\Res_B^b,\qquad
 X_{b,\mu,s}\in\Sc_1,\qquad
 \norm{X_{b,\mu,s}}_1\le F_bs^{-2}.
 \label{eq:forward-trace}
\end{equation}
\end{theorem}

\begin{proof}
The coordinates \(\bx\) and \(\by\) are measured from the centres \(0\)
and \(s\eone\) of the two balls. With \(B_b\) as in \eqref{eq:Lambda},
we set
\[
 \mathfrak d_b=2B_b,\qquad
 d_{0,b}=2\sqrt3\,\mathfrak d_b=4\sqrt3B_b.
\]

\smallskip
\emph{Step 1: the square root on the polydisc.}
On the six-variable polydisc
\(|x_j|,|y_j|\le\mathfrak d_b\),
\[
 \mathbf w=s\eone+\by-\bx,\qquad
 |\mathbf w-s\eone|_{\mathrm E}\le d_{0,b},
\]
because each of the three components of \(\by-\bx\) has modulus at
most \(2\mathfrak d_b\). Let \(s\ge4d_{0,b}\), and put
\(q=d_{0,b}/s\le1/4\). Then
\[
 \begin{split}
 \mathbf w\cdot\mathbf w&=s^2(1+\epsilon),\qquad
 \epsilon=\frac{2(y_1-x_1)}s+\frac{(\by-\bx)\cdot(\by-\bx)}{s^2},\\
 |\epsilon|&\overset{|\by-\bx|_{\mathrm E}\le d_{0,b}}{\le}
 \frac{2d_{0,b}}s+\frac{d_{0,b}^2}{s^2}=2q+q^2\le\frac9{16}.
 \end{split}
\]
So \(1+\epsilon\) lies in the disc of radius \(9/16\) about \(1\), on
which the principal square root is holomorphic, and the square root
\(\varrho=s\sqrt{1+\epsilon}\) in \eqref{eq:kernels} is holomorphic on
the polydisc.

The real part of \(\varrho\) satisfies
\[
 \begin{split}
 \Rea\varrho
 &\overset{\phantom{q\le2/5}}{=}s\Rea\sqrt{1+\epsilon}
 \ge s\sqrt{1-|\epsilon|}
 \overset{|\epsilon|\le2q+q^2}{\ge}s\sqrt{1-2q-q^2}\\
 &\overset{q\le2/5}{\ge}s(1-2q)
 =s-2d_{0,b}
 \overset{q\le1/4}{\ge}\frac s2.
 \end{split}
\]
The first inequality holds because
\((\Rea\sqrt{1+\epsilon})^2=(|1+\epsilon|+1+\Rea\epsilon)/2\ge1-|\epsilon|\),
and the inequality marked \(q\le2/5\) is \((1-2q)^2\le1-2q-q^2\), that
is, \(5q^2\le2q\). Moreover,
\(|\varrho|=s\,|1+\epsilon|^{1/2}\le s\,(1+9/16)^{1/2}=5s/4\).

\smallskip
\emph{Step 2: the suprema of the two kernels.}
We write \(\sup\) for the supremum over the polydisc. At each of its
points, \(\mu\Rea\varrho>0\) by Step~1, so the bound
\(|K_j(z)|\le K_j(\Rea z)\) at \(z=\mu\varrho\) and
\eqref{eq:Bessel-scaled} at \(x=\mu\Rea\varrho\) give
\[
 \begin{split}
 |\mathfrak X_\mu(\mathbf w)|
 &\overset{\eqref{eq:kernels}}{=}
 \frac{\mu|K_1(\mu\varrho)|}{4\pi^2|\varrho|}
 \le\frac{\mu K_1(\mu\Rea\varrho)}{4\pi^2|\varrho|}
 \overset{\eqref{eq:Bessel-scaled}}{\le}
 \frac{3(1+\mu\Rea\varrho)^{1/2}e^{-\mu\Rea\varrho}}
 {4\pi^2|\varrho|\Rea\varrho},\\
 |\mathfrak P_\mu(\mathbf w)|
 &\overset{\eqref{eq:kernels}}{=}
 \frac{\mu^2|K_2(\mu\varrho)|}{4\pi^2|\varrho|^2}
 \le\frac{\mu^2K_2(\mu\Rea\varrho)}{4\pi^2|\varrho|^2}
 \overset{\eqref{eq:Bessel-scaled}}{\le}
 \frac{8(1+\mu\Rea\varrho)^{3/2}e^{-\mu\Rea\varrho}}
 {4\pi^2|\varrho|^2(\Rea\varrho)^2}.
 \end{split}
\]
By Step~1, \(|\varrho|\ge\Rea\varrho\ge s/2\),
\(1+\mu\Rea\varrho\le1+\mu|\varrho|\le1+5\xi/4\le5(1+\xi)/4\), and
\(e^{-\mu\Rea\varrho}\le e^{-\mu(s-2d_{0,b})}=e^{-\xi+2\mu d_{0,b}}\).
Thus
\[
 \begin{split}
 \sup|\mathfrak X_\mu|
 &\le Cs^{-2}(1+\xi)^{1/2}e^{-\xi+2\mu d_{0,b}},\\
 \sup|\mathfrak P_\mu|
 &\le Cs^{-4}(1+\xi)^{3/2}e^{-\xi+2\mu d_{0,b}}.
 \end{split}
\]

\smallskip
\emph{Step 3: the Taylor coefficients.}
By Step~1, \(\varrho\) is holomorphic on the polydisc with
\(\Rea\varrho\ge s/2\), and \(K_1\) and \(K_2\) are holomorphic on the
right half-plane. So \(\mathfrak X_\mu(s\eone+\by-\bx)\) and
\(\mathfrak P_\mu(s\eone+\by-\bx)\) are holomorphic functions of
\((\bx,\by)\) on the polydisc; let \(c^f_{\alpha\beta}\) and
\(c^g_{\alpha\beta}\), \(\alpha,\beta\in\N^3\), be their Taylor
coefficients at the origin. By the Cauchy estimates on the polydisc,
\[
 \begin{split}
 \mathfrak X_\mu(s\eone+\by-\bx)
 &=\sum_{\alpha,\beta\in\N^3}c^f_{\alpha\beta}\,\bx^\alpha\by^\beta,
 \qquad
 |c^f_{\alpha\beta}|
 \le\sup|\mathfrak X_\mu|\,\mathfrak d_b^{-|\alpha|-|\beta|},\\
 \mathfrak P_\mu(s\eone+\by-\bx)
 &=\sum_{\alpha,\beta\in\N^3}c^g_{\alpha\beta}\,\bx^\alpha\by^\beta,
 \qquad
 |c^g_{\alpha\beta}|
 \le\sup|\mathfrak P_\mu|\,\mathfrak d_b^{-|\alpha|-|\beta|}.
 \end{split}
\]
On the two unit balls, \(|x_j|,|y_j|<1<\mathfrak d_b\), so both series
converge absolutely and uniformly there.

\smallskip
\emph{Step 4: the rank-one moment series and
\eqref{eq:forward-refined}.}
Let \(\sx=(f,g)\) and \(\sy=(f',g')\) be smooth real sources supported
in \(B_1(0)\) and \(B_1(s\eone)\). In the coordinates of the two balls,
the formula after \eqref{eq:kernels} and term-by-term integration of
the series of Step~3 expand the covariance into products of the source
moments of Lemma~\ref{lem:moments}, one rank-one product for each
monomial kernel:
\[
 \begin{split}
 \vv_\mu(\sx,\sy)
 &=\iint\bigl[f(\bx)\,\mathfrak X_\mu(s\eone+\by-\bx)\,f'(s\eone+\by)\\
 &\qquad\quad+g(\bx)\,\mathfrak P_\mu(s\eone+\by-\bx)\,g'(s\eone+\by)\bigr]
 \,d^3x\,d^3y\\
 &=\sum_{\alpha,\beta\in\N^3}
 \bigl(c^f_{\alpha\beta}\,\ell^f_\alpha(\sx)\,\ell^f_\beta(\sy)
 +c^g_{\alpha\beta}\,\ell^g_\alpha(\sx)\,\ell^g_\beta(\sy)\bigr).
 \end{split}
\]

Lemma~\ref{lem:moments} at \(T=1\), for \(U=B_1(0)\) with \(\bx_0=0\)
and for \(U=B_1(s\eone)\) with \(\bx_0=s\eone\), turns
\eqref{eq:f-moments}--\eqref{eq:g-moments} into
\[
 \norm{\ell^f_\alpha\Res_A^{-b}}
 \le M_b\frac{e^{\Lambda_b\mu}}{1+\mu}B_b^{|\alpha|},
 \qquad
 \norm{\ell^g_\alpha\Res_A^{-b}}
 \le M_b\frac{e^{\Lambda_b\mu}}{(1+\mu)^2}B_b^{|\alpha|},
\]
and into the same bounds with \(\Res_B\) for the moments centred at
\(s\eone\). By Convention~\ref{conv:factorization} and the Riesz
representation, for \(\bullet\in\{f,g\}\) there are vectors
\(\mathsf u^\bullet_{A,\alpha}\in\Hh_{\mu,B_1(0)}\) and
\(\mathsf u^\bullet_{B,\beta}\in\Hh_{\mu,B_1(s\eone)}\) such that
\[
 \ell^\bullet_\alpha(\sx)
 =\ip{\mathsf u^\bullet_{A,\alpha}}{\Res_A^b\sx},
 \qquad
 \ell^\bullet_\beta(\sy)
 =\ip{\mathsf u^\bullet_{B,\beta}}{\Res_B^b\sy}
\]
for all \(\sx\in\Hh_{\mu,B_1(0)}\) and \(\sy\in\Hh_{\mu,B_1(s\eone)}\),
with \(\lVert\mathsf u^\bullet_{A,\alpha}\rVert
=\lVert\ell^\bullet_\alpha\Res_A^{-b}\rVert\) and
\(\lVert\mathsf u^\bullet_{B,\beta}\rVert
=\lVert\ell^\bullet_\beta\Res_B^{-b}\rVert\). We put
\[
 X_{b,\mu,s}
 =\sum_{\bullet\in\{f,g\}}\sum_{\alpha,\beta\in\N^3}
 c^\bullet_{\alpha\beta}\,
 \mathsf u^\bullet_{A,\alpha}\ip{\mathsf u^\bullet_{B,\beta}}{\cdot\,},
\]
a series of rank-one operators.

The term of \((\bullet,\alpha,\beta)\) has trace norm
\(|c^\bullet_{\alpha\beta}|\,\lVert\mathsf u^\bullet_{A,\alpha}\rVert\,
\lVert\mathsf u^\bullet_{B,\beta}\rVert\). By Step~3 and the moment bounds,
with \(B_b/\mathfrak d_b=1/2\) and
\(\sum_{\alpha,\beta\in\N^3}2^{-|\alpha|-|\beta|}
=\bigl(\sum_{k\ge0}2^{-k}\bigr)^6=64\), these trace norms are summable;
so the series converges absolutely in trace norm,
\(X_{b,\mu,s}\in\Sc_1\), and
\[
 \begin{split}
 \norm{X_{b,\mu,s}}_1
 &\le\sum_{\bullet,\alpha,\beta}|c^\bullet_{\alpha\beta}|\,
 \norm{\mathsf u^\bullet_{A,\alpha}}\,\norm{\mathsf u^\bullet_{B,\beta}}\\
 &\le M_b^2e^{2\Lambda_b\mu}
 \left[\frac{\sup|\mathfrak X_\mu|}{(1+\mu)^2}
 +\frac{\sup|\mathfrak P_\mu|}{(1+\mu)^4}\right]
 \sum_{\alpha,\beta\in\N^3}2^{-|\alpha|-|\beta|}\\
 &=64\,M_b^2e^{2\Lambda_b\mu}
 \left[\frac{\sup|\mathfrak X_\mu|}{(1+\mu)^2}
 +\frac{\sup|\mathfrak P_\mu|}{(1+\mu)^4}\right].
 \end{split}
\]
Inserting the bounds of Step~2 for the two suprema, with
\(2\Lambda_b\mu+2\mu d_{0,b}=(2\Lambda_b+8\sqrt3B_b)\mu\), gives
\begin{equation}
 \begin{split}
 \norm{X_{b,\mu,s}}_1
 \le{}&CM_b^2e^{\mathfrak k_b\mu-\xi}
 \left[
 \frac{s^{-2}(1+\xi)^{1/2}}{(1+\mu)^2}
 +\frac{s^{-4}(1+\xi)^{3/2}}{(1+\mu)^4}
 \right],\\
 &\mathfrak k_b=2\Lambda_b+8\sqrt3B_b.
 \end{split}
 \label{eq:forward-refined}
\end{equation}

\smallskip
\emph{Step 5: the threshold \(S_b\) and the constant \(F_b\).}
A sufficient geometric condition is
\begin{equation}
 S_b\ge\max\{4,4d_{0,b},2\mathfrak k_b\}.
 \label{eq:forward-geometry}
\end{equation}
Let \(s\ge S_b\). The entry \(4d_{0,b}\) is the condition of Step~1;
the entry \(4\) gives \(s>2\), the separation assumed throughout this
section, and \(s\ge1\); the entry \(2\mathfrak k_b\) gives
\(\mathfrak k_b\mu\le\mu s/2=\xi/2\). So the exponential in
\eqref{eq:forward-refined} is at most \(e^{-\xi/2}\).

Since \(1+\mu\ge1\) and \(s^{-4}\le s^{-2}\), the bracket in
\eqref{eq:forward-refined} is at most
\(s^{-2}\bigl[(1+\xi)^{1/2}+(1+\xi)^{3/2}\bigr]\le2s^{-2}(1+\xi)^{3/2}\).
Hence, with \(C\) the constant of \eqref{eq:forward-refined},
\[
 \norm{X_{b,\mu,s}}_1\le2CM_b^2s^{-2}(1+\xi)^{3/2}e^{-\xi/2}.
\]
The finite supremum \(3^{3/2}e^{-1}\) of \((1+\xi)^{3/2}e^{-\xi/2}\)
over \(\xi\ge0\), attained at \(\xi=2\), gives
\(F_b={2\cdot3^{3/2}e^{-1}CM_b^2}\) in \eqref{eq:forward-trace}. The
numbers \(S_b\) and \(F_b\) depend only on \(b\).

\smallskip
\emph{Step 6: the identification with \(C_{\mu,s}\).}
Let \(\sx,\sy\) be smooth real sources as in Step~4.
Their moments are real, so that
\[
 \ell^\bullet_\alpha(\sx)
 =\overline{\ip{\mathsf u^\bullet_{A,\alpha}}{\Res_A^b\sx}}
 =\ip{\Res_A^b\sx}{\mathsf u^\bullet_{A,\alpha}},
\]
and
\[
 \begin{split}
 \ip{\sx}{\Res_A^bX_{b,\mu,s}\Res_B^b\sy}
 &=\sum_{\bullet,\alpha,\beta}c^\bullet_{\alpha\beta}\,
 \ip{\Res_A^b\sx}{\mathsf u^\bullet_{A,\alpha}}\,
 \ip{\mathsf u^\bullet_{B,\beta}}{\Res_B^b\sy}\\
 &=\sum_{\bullet,\alpha,\beta}c^\bullet_{\alpha\beta}\,
 \ell^\bullet_\alpha(\sx)\,\ell^\bullet_\beta(\sy)
 \overset{\text{Step 4}}{=}\vv_\mu(\sx,\sy)
 =\ip{\sx}{C_{\mu,s}\sy}.
 \end{split}
\]

The bounded operators \(\Res_A^bX_{b,\mu,s}\Res_B^b\) and \(C_{\mu,s}\)
thus have the same sesquilinear form on pairs of smooth real sources. By
sesquilinearity the two forms agree on complex combinations of such
sources, which are dense in the complexified source spaces
(Section~\ref{subsec:sources}); so density identifies the full
cross-covariance operator, \(C_{\mu,s}=\Res_A^bX_{b,\mu,s}\Res_B^b\).
With the bound of Step~5, this factorization is
\eqref{eq:forward-trace}.
\end{proof}

\section{The reverse estimate in the joint covariance norm}
\label{sec:reverse}

For
\[
 U_s=B_1(0)\cup B_1(s\eone),
\]
the source Hilbert space \(\Hh_{\mu,U_s}\) now carries its joint vacuum
covariance norm. Its uncertainty reserve is denoted by
\(\JointRes_{\mu,s}\). The product covariance \(\vv_{\mathrm{prod},\mu,s}\)
is obtained from the joint covariance \(\vv_{\mu,U_s}\) by deleting the
two cross terms between the balls. The product-minus-joint
covariance operator is defined by
\begin{equation}
 \ip{\sx}{\JointD_{\mu,s}\sy}_{\vv_{\mu,U_s}}
 =
 \vv_{\mathrm{prod},\mu,s}(\sx,\sy)
 -\vv_{\mu,U_s}(\sx,\sy).
 \label{eq:D-joint}
\end{equation}
The form \(\vv_{\mathrm{prod},\mu,s}-\vv_{\mu,U_s}\) is bounded for all
sufficiently large \(s\) (Lemma~\ref{lem:polynomial}).

\subsection{Smooth kernels in the joint source norm}

The multiplier inequality
\(h_\mu\le\max\{1,\mu\}h_1\) gives
\begin{equation}
 \norm{h_1^{-1/2}f}_2
 \le\sqrt{\max\{1,\mu\}}\,
       \norm{h_\mu^{-1/2}f}_2.
 \label{eq:f-comparison}
\end{equation}
For \(g\) supported in \(U_s\), a set of volume at most \(8\pi/3\), twice
the volume \(4\pi/3\) of a unit ball, the unitary Fourier transform and
the Cauchy--Schwarz inequality give, for every \(\bp\),
\begin{align*}
 |\mathcal Fg(\bp)|^2
 &\le\frac1{(2\pi)^3}\Bigl(\int_{U_s}|g|\,d^3x\Bigr)^2
 \le\frac{\vol(U_s)}{(2\pi)^3}\norm{g}_2^2\\
 &\le\frac{8\pi/3}{8\pi^3}\norm{g}_2^2
 =\frac1{3\pi^2}\norm{g}_2^2,
\end{align*}
and integration over the unit ball \(\{|\bp|<1\}\), of volume
\(4\pi/3\), gives
\[
 \int_{|\bp|<1}|\mathcal Fg(\bp)|^2\,d^3p
 \le\frac{4\pi}3\cdot\frac1{3\pi^2}\norm{g}_2^2
 =\frac4{9\pi}\norm{g}_2^2
 <\frac12\norm{g}_2^2.
\]
Consequently, by Plancherel, the part of \(\norm{g}_2^2\) from the
frequencies \(|\bp|\ge1\) exceeds \(\tfrac12\norm{g}_2^2\), so that
\begin{equation}
 \norm{g}_2^2
 \le2\int|\bp|\,|\mathcal Fg(\bp)|^2\,d^3p
 \le2\norm{h_\mu^{1/2}g}_2^2.
 \label{eq:g-comparison}
\end{equation}

We fix a smooth cutoff \(\chi_0\) on \(\R^3\) with
\(0\le\chi_0\le1\), equal to \(1\) on \(B_1(0)\) and with compact
support in \(B_2(0)\), and put
\(\chi(\bx)=\chi_0(\bx)+\chi_0(\bx-s\eone)\). For \(s\ge5\), the
supports of the two terms are disjoint, so \(\chi=1\) on \(U_s\), and
the derivative bounds of \(\chi\) are those of \(\chi_0\), independent
of \(s\).

Let \(K\) be a scalar kernel on
\((B_2(0)\cup B_2(s\eone))^2\), with derivatives through order two in
each variable bounded by \(M\):
\[
 |\partial_\bx^\alpha\partial_\by^\beta K(\bx,\by)|\le M,
 \qquad |\alpha|,|\beta|\le2.
\]
In the joint source norm \(\norm{\sx}=\vv_{\mu,U_s}(\sx,\sx)^{1/2}\),
the field-source form and the momentum-source form of \(K\) satisfy,
for smooth real sources \(\sx=(f,g)\) and \(\sy=(f',g')\) supported in
\(U_s\),
\[
 \begin{split}
 \Bigl|\iint f(\bx)K(\bx,\by)f'(\by)\,d^3x\,d^3y\Bigr|
 &\le C\max\{1,\mu\}M\norm{\sx}\norm{\sy},\\
 \Bigl|\iint g(\bx)K(\bx,\by)g'(\by)\,d^3x\,d^3y\Bigr|
 &\le CM\norm{\sx}\norm{\sy}
 \end{split}
\]
(the \emph{smooth-kernel estimate}).

Let \(K_\chi(\bx,\by)=\chi(\bx)K(\bx,\by)\chi(\by)\), extended by zero
to \(\R^6\); it equals \(K\) on \(U_s\times U_s\), and the integral
operator with kernel \(K_\chi\) is denoted by the same symbol. By the
Leibniz rule, \(|\partial_\bx^\alpha\partial_\by^\beta K_\chi|\le CM\)
for \(|\alpha|,|\beta|\le2\), and \(K_\chi\) is supported in a set of
volume at most \((2\vol B_2(0))^2\).

For the field form, multiplication by the cutoffs and insertion of
\(h_1^{1/2}\) on both sides give the Hilbert--Schmidt operator
\(h_1^{1/2}K_\chi h_1^{1/2}\), controlled by the kernel's \(H^2\) norm
in each variable. With \(\mathcal FK_\chi\) the unitary Fourier
transform of \(K_\chi\) on \(\R^6\),
\begin{align*}
 \norm{h_1^{1/2}K_\chi h_1^{1/2}}_2^2
 &\overset{\text{Plancherel}}{=}
   \iint(1+|\bp|^2)^{1/2}(1+|\mathbf q|^2)^{1/2}
   |\mathcal FK_\chi(\bp,\mathbf q)|^2\,d^3p\,d^3q\\
 &\overset{\phantom{\text{Plancherel}}}{\le}\iint(1+|\bp|^2)^{2}(1+|\mathbf q|^2)^{2}
   |\mathcal FK_\chi(\bp,\mathbf q)|^2\,d^3p\,d^3q\\
 &\overset{\text{Plancherel}}{\le}C\sum_{|\alpha|,|\beta|\le2}
   \norm{\partial_\bx^\alpha\partial_\by^\beta K_\chi}_{L^2(\R^6)}^2
 \overset{\text{Leibniz}}{\le}CM^2 .
\end{align*}
Since \(f=\chi f\) and \(f'=\chi f'\), and \(h_1^{\pm1/2}\) are
self-adjoint, \eqref{eq:f-comparison} completes the estimate:
\begin{align*}
 \Bigl|\iint fKf'\Bigr|
 &\overset{\phantom{\eqref{eq:f-comparison}}}{=}\bigl|\ip{h_1^{-1/2}f}{(h_1^{1/2}K_\chi h_1^{1/2})\,h_1^{-1/2}f'}\bigr|
 \le\norm{h_1^{1/2}K_\chi h_1^{1/2}}_2\,
   \norm{h_1^{-1/2}f}_2\,\norm{h_1^{-1/2}f'}_2\\
 &\overset{\eqref{eq:f-comparison}}{\le}\max\{1,\mu\}\,
   \norm{h_1^{1/2}K_\chi h_1^{1/2}}_2\,
   \norm{h_\mu^{-1/2}f}_2\,\norm{h_\mu^{-1/2}f'}_2
 \overset{\eqref{eq:covariance}}{\le}C\max\{1,\mu\}M\norm{\sx}\norm{\sy},
\end{align*}
where the last step uses \(\lVert h_\mu^{-1/2}f\rVert_2^2\le2\lVert\sx\rVert^2\)
and the bound \(CM\) of the Hilbert--Schmidt norm.

For the momentum form, the Hilbert--Schmidt kernel norm and
\eqref{eq:g-comparison} suffice: the Hilbert--Schmidt norm of
\(K_\chi\) is \(\norm{K_\chi}_{L^2(\R^6)}\le CM\), and
\begin{align*}
 \Bigl|\iint gKg'\Bigr|
 &\overset{\phantom{\eqref{eq:g-comparison}}}{=}|\ip{g}{K_\chi g'}|
 \le\norm{K_\chi}_{L^2(\R^6)}\norm{g}_2\norm{g'}_2\\
 &\overset{\eqref{eq:g-comparison}}{\le}
   2\norm{K_\chi}_{L^2(\R^6)}\norm{h_\mu^{1/2}g}_2\norm{h_\mu^{1/2}g'}_2
 \overset{\eqref{eq:covariance}}{\le}CM\norm{\sx}\norm{\sy}.
\end{align*}

All cutoff supports and derivative norms have fixed size, so the
constants \(C\) are independent of \(s\ge5\) and of \(\mu\). Density
extends the estimates to \(\mathcal X_{\mu,U_s}\).

\subsection{Polynomial approximation on the disconnected support}

The polynomial selector
\begin{equation}
 \mathsf p_N(z)
 =(1-z)^N\sum_{k=0}^{N-1}\binom{N+k-1}{k}z^k
 \label{eq:selectors}
\end{equation}
has degree \(2N-1\); \(\mathsf p_N-1\) has a zero of order \(N\) at
zero, and \(\mathsf p_N\) has a zero of order \(N\) at one. For the zero of \(\mathsf p_N-1\)
at zero, the expansion
\((1-z)^{-N}=\sum_{k\ge0}\binom{N+k-1}{k}z^k\), \(|z|<1\), gives
\[
 \mathsf p_N(z)-1
 =-(1-z)^N\sum_{k\ge N}\binom{N+k-1}{k}z^k .
\]

For \(N\ge1\) and \(0<r<1/2\), the identity for \(\mathsf p_N-1\)
above, the definition \eqref{eq:selectors}, the binomial bound \(\binom{N+k-1}{k}\le2^{N+k-1}\) and
\(\sum_{k=0}^{N-1}\binom{N+k-1}{k}=\binom{2N-1}{N-1}\le2^{2N-1}\) give
the endpoint bounds
\[
 \begin{aligned}
 |\mathsf p_N(z)-1|
 &\le(1+r)^N\sum_{k\ge N}2^{N+k-1}r^k
 =\frac{(4r(1+r))^N}{2(1-2r)},&\qquad&|z|\le r,\\
 |\mathsf p_N(z)|
 &\le r^N(1+r)^{N-1}\binom{2N-1}{N-1}
 \le(4r(1+r))^N,&&|z-1|\le r.
 \end{aligned}
\]

The coefficient sum of a polynomial, the sum of the absolute values of
its coefficients, is submultiplicative. Expansion of \((1-z)^N\) gives
the coefficient sum \(2^N\), and the binomial bound gives
\(\binom{2N-1}{N-1}\le2^{2N-1}\) for the sum in \eqref{eq:selectors};
so the coefficient sum of \(\mathsf p_N\) is at most \(2^{3N}\). The
fixed affine shifts used below are \(z\mapsto z+1/2\) and
\(z\mapsto1/2-z\). Each turns \(z^k\) into a polynomial of coefficient
sum \((3/2)^k\), so after either shift the coefficient sum of
\(\mathsf p_N\) is at most \((3/2)^{2N-1}\,2^{3N}\le C^N\).

\begin{lemma}[Polynomial increments]\label{lem:polynomial}
There are universal finite constants \(C_1,C_2,c_{\mathrm{pol}},K_{\mathrm{pol}}\)
and a sufficiently large universal \(S_{\mathrm{pol}}\) with the
following properties. For every \(\mu>0\) and \(s\ge S_{\mathrm{pol}}\),
with \(\xi=\mu s\), there are real self-adjoint finite-rank forms
\(\mathsf F_n\) with polynomial kernels, \(n\ge1\), with
\(\mathsf F_0=0\), such that \(\mathsf F_n\to\JointD_{\mu,s}\) in
operator norm, where \(\JointD_{\mu,s}\) is the product-minus-joint
covariance operator \eqref{eq:D-joint}. Their increments
\(\mathsf G_n=\mathsf F_{n+1}-\mathsf F_n\), \(n\ge0\), satisfy
\begin{equation}
 \norm{\mathsf G_n}_1
 \le C_1(n+1)^3s^{-2}(1+\xi)^2e^{-\xi/2}(c_{\mathrm{pol}}/s)^n.
 \label{eq:increment-unweighted}
\end{equation}
For \(\Lambda_0=\Lambda_{1/2}\), the constant \eqref{eq:Lambda} at
\(b=1/2\), and \(T=s/2+1\), there are finite-rank \(B_n\) such that,
with the joint reserve \(\JointRes_{\mu,s}\),
\begin{equation}
 \mathsf G_n
 =\JointRes_{\mu,s}^{1/2}B_n\JointRes_{\mu,s}^{1/2},
 \qquad
 \norm{B_n}_1
 \le C_2e^{-\xi/2}e^{2\Lambda_0\mu T}K_{\mathrm{pol}}^n.
 \label{eq:increment-weighted}
\end{equation}
\end{lemma}

\begin{proof}
\emph{Step 1: the cross kernels on a polydisc.} In dimensionless
coordinates \(u=\bx/s\), \(v=\by/s\), the support centres are \(0\)
and \(\eone\). We use a polydisc of radius \(r_0=1/100\) about the
cross centre \((0,\eone)\), and the slightly larger polydisc of radius
\(2r_0\). On the latter, \(v-u=\eone+\boldsymbol\delta\) with
\(\boldsymbol\delta\in\C^3\) and
\(|\boldsymbol\delta|_{\mathrm E}\le4\sqrt3\,r_0<1/10\), so
\[
 (v-u)\cdot(v-u)=1+\epsilon,\qquad
 |\epsilon|=|2\delta_1+\boldsymbol\delta\cdot\boldsymbol\delta|
 \le2|\boldsymbol\delta|_{\mathrm E}+|\boldsymbol\delta|_{\mathrm E}^2
 <\tfrac14 .
\]
Hence the square root \(\varrho=s\sqrt{1+\epsilon}\) in
\eqref{eq:kernels} is holomorphic there, and, with
\(\Rea\sqrt{1+\epsilon}\ge\sqrt{1-|\epsilon|}\) as in the proof of
Theorem~\ref{thm:forward},
\[
 \tfrac s2\le\Rea\varrho\le|\varrho|\le2s .
\]

The Bessel bounds \eqref{eq:Bessel-scaled} give the polydisc suprema
\[
 M_f=Cs^{-2}(1+\xi)^{1/2}e^{-\xi/2},\qquad
 M_g=Cs^{-4}(1+\xi)^{3/2}e^{-\xi/2}
\]
of the two cross kernels \(\mathfrak X_\mu(s(v-u))\) and
\(\mathfrak P_\mu(s(v-u))\): on the polydisc of radius \(2r_0\), the
pointwise bounds of Step~2 of the proof of Theorem~\ref{thm:forward}
hold at \(\mathbf w=s(v-u)\), and their right-hand sides are at most
\(M_f\) and \(M_g\) because \(\tfrac s2\le\Rea\varrho\le|\varrho|\le2s\),
so that \(1+\mu\Rea\varrho\le2(1+\xi)\) and
\(e^{-\mu\Rea\varrho}\le e^{-\xi/2}\).

\smallskip
\emph{Step 2: the polynomial forms.} On the polydisc of radius
\(2r_0\), the two cross kernels are holomorphic, and real at real
points; so their Taylor expansions at the cross centre,
\[
 \mathfrak X_\mu(s(v-u))
 =\sum_{\alpha,\beta\in\N^3}a^f_{\alpha\beta}\,u^\alpha(v-\eone)^\beta,
 \qquad
 \mathfrak P_\mu(s(v-u))
 =\sum_{\alpha,\beta\in\N^3}a^g_{\alpha\beta}\,u^\alpha(v-\eone)^\beta,
\]
have real coefficients, and Cauchy's estimate on the polydisc of
radius \(r_0\) gives
\(|a^\bullet_{\alpha\beta}|\le M_\bullet r_0^{-|\alpha|-|\beta|}\).

Let \(\mathcal T_{n,f}(u,v-\eone)\) and \(\mathcal T_{n,g}(u,v-\eone)\)
denote the Taylor polynomials of total degree \(n-1\) of these two
cross kernels at the cross centre \((0,\eone)\), in the coordinates
\((u,v-\eone)\); thus \(\mathcal T_{0,\bullet}=0\). For
\(\bullet\in\{f,g\}\), the
\emph{oriented polynomial}, whose selector factors equal \(1\) to high
order on the cross component \(u\approx0\), \(v\approx\eone\) and
vanish to high order on the other three components of the support, is
\[
 \mathsf p_{2n}(u_1)\,
 \mathsf p_{2n}(1-v_1)\,
 \mathcal T_{n,\bullet}(u,v-\eone).
\]

Adding the transpose of the oriented polynomial and taking the negative
sign from \eqref{eq:D-joint} gives the polynomial kernel
\[
 f_{n,\bullet}(\bx,\by)
 =-\mathsf p_{2n}(u_1)\,\mathsf p_{2n}(1-v_1)\,
   \mathcal T_{n,\bullet}(u,v-\eone)
  -\mathsf p_{2n}(v_1)\,\mathsf p_{2n}(1-u_1)\,
   \mathcal T_{n,\bullet}(v,u-\eone),
\]
which is real and symmetric in \((\bx,\by)\), and defines
\(\mathsf F_n\) by
\[
 \ip{\sx}{\mathsf F_n\sy}_{\vv_{\mu,U_s}}
 =\iint\bigl[f(\bx)f_{n,f}(\bx,\by)f'(\by)
 +g(\bx)f_{n,g}(\bx,\by)g'(\by)\bigr]\,d^3x\,d^3y
\]
for smooth real sources supported in \(U_s\); in particular
\(\mathsf F_0=0\).

Since \(\mathsf p_{2n}\) has degree \(4n-1\), the degree of
\(f_{n,\bullet}\) is at most \(2(4n-1)+(n-1)=9n-3\le9n\). Expanded in
the monomials \(\bx^\alpha\by^\beta\), the form of \(\mathsf F_n\) is a
finite real combination of the products
\(\ell^f_\alpha(\sx)\ell^f_\beta(\sy)\) and
\(\ell^g_\alpha(\sx)\ell^g_\beta(\sy)\), \(|\alpha|,|\beta|\le9n\), of
the moments of Lemma~\ref{lem:moments}. That lemma, applied at
\(b=1/2\) in a ball about \(0\) containing \(U_s\), bounds each of these
moments at \(\sx\) by a finite multiple of
\(\lVert\JointRes_{\mu,s}^{1/2}\sx\rVert\), and
\(\lVert\JointRes_{\mu,s}^{1/2}\sx\rVert\le\lVert\sx\rVert\) since
\(\JointRes_{\mu,s}\le I\). So the moments are bounded on
\(\mathcal X_{\mu,U_s}\), and \(\mathsf F_n\) is a bounded real
self-adjoint operator of finite rank. Its
form depends on \(\sx\) only through the moments \(\ell^f_\alpha(\sx)\)
and \(\ell^g_\alpha(\sx)\), two for each of the \(\binom{9n+3}{3}\)
multi-indices \(\alpha\in\N^3\) with \(|\alpha|\le9n\), so that
\[
 \rank\mathsf F_n\le2\binom{9n+3}{3}\le C(n+1)^3 .
\]

\smallskip
\emph{Step 3: the differentiated error on the four components.} Let
\(\mathfrak K_\bullet\) be the target product-minus-joint kernel, equal
to the negative covariance kernel on the two cross components and zero
on the two same-region components. The definition \eqref{eq:D-joint} reads
\[
 \ip{\sx}{\JointD_{\mu,s}\sy}_{\vv_{\mu,U_s}}
 =\iint\bigl[f(\bx)\mathfrak K_f(\bx,\by)f'(\by)
 +g(\bx)\mathfrak K_g(\bx,\by)g'(\by)\bigr]\,d^3x\,d^3y
\]
for smooth real sources supported in \(U_s\). Let
\(\mathcal N_s=(B_2(0)\cup B_2(s\eone))^2\).
Then
\begin{equation}
 \max_{|\alpha|,|\beta|\le2}
 \sup_{\mathcal N_s}
 \left|
 \partial_\bx^\alpha\partial_\by^\beta
 (\mathfrak K_\bullet-f_{n,\bullet})
 \right|
 \le CM_\bullet(c_{\mathrm{pol}}/s)^n.
 \label{eq:differentiated-error}
\end{equation}

The selector and Taylor estimates for \eqref{eq:differentiated-error}
hold on slightly larger complex physical neighbourhoods of all four
components. Let
\(\mathcal N_s^{\C}\subset\C^6\) be the union of the four open
polydiscs of radius \(3\), in the six coordinates of \((\bx,\by)\),
about the points of \(\{0,s\eone\}^2\). It contains the closed unit
polydisc about every point of \(\mathcal N_s\), and in the coordinates
\((u,v)\) its four parts are the polydiscs of radius \(3/s\) about the
points of \(\{0,\eone\}^2\).

The universal \(S_{\mathrm{pol}}\) is taken with
\(3/S_{\mathrm{pol}}\le r_0/2\) and larger than the constant \(c_{\mathrm{pol}}\)
fixed at the end of this step. For \(s\ge S_{\mathrm{pol}}\), then,
\(s>6\), so the four parts are disjoint, and the cross part about
\((0,\eone)\) lies in the polydisc of radius \(r_0/2\) about the cross
centre. The exchange \((u,v)\mapsto(v,u)\) carries this cross part onto
the other one; \(\mathfrak K_\bullet\) is symmetric under the exchange,
because the kernels in \eqref{eq:kernels} are even in \(\mathbf w\),
and so is \(f_{n,\bullet}\). Continued as \(-\mathfrak X_\mu(s(v-u))\)
or \(-\mathfrak P_\mu(s(v-u))\) on the cross parts and as zero on the
others, \(\mathfrak K_\bullet\) is holomorphic on \(\mathcal N_s^{\C}\)
by Step~2.

At points where the six coordinates of \((u,v-\eone)\) have moduli at
most \(\rho\), the bound on \(a^\bullet_{\alpha\beta}\) from Step~2 and
the number \(\binom{k+5}{5}\le(k+1)^5\) of monomials of degree \(k\)
give
\[
 \sum_{|\alpha|+|\beta|\ge n}
 |a^\bullet_{\alpha\beta}u^\alpha(v-\eone)^\beta|
 \le M_\bullet\sum_{k\ge n}(k+1)^5\Bigl(\frac\rho{r_0}\Bigr)^k,
 \qquad
 |\mathcal T_{n,\bullet}(u,v-\eone)|
 \le M_\bullet\sum_{k<n}(k+1)^5\Bigl(\frac\rho{r_0}\Bigr)^k .
\]
On the cross part about \((0,\eone)\), \(\rho=3/s\le r_0/2\); with
\((k+1)^5\le(n+1)^5(k-n+1)^5\) for \(k\ge n\), the Taylor remainder is
at most \(C(n+1)^5M_\bullet(C/s)^n\) there, and
\(|\mathcal T_{n,\bullet}(u,v-\eone)|\le CM_\bullet\). On every part,
the arguments \((u,v-\eone)\) and \((v,u-\eone)\) of
\(\mathcal T_{n,\bullet}\) have coordinates of moduli at most
\(1+3/s\le2\), and there
\(|\mathcal T_{n,\bullet}|\le n^6(2/r_0)^nM_\bullet\le(n+1)^6C^nM_\bullet\).

Let \(n\ge1\). By the endpoint bounds after \eqref{eq:selectors}, with
\(N=2n\) and \(r=3/s\le1/200\): at arguments within \(3/s\) of zero,
\(|\mathsf p_{2n}-1|\le(C/s)^{2n}\le1\), so \(|\mathsf p_{2n}|\le2\);
at arguments within \(3/s\) of one, \(|\mathsf p_{2n}|\le(C/s)^{2n}\).

On the cross part about \((0,\eone)\), where
\(\mathfrak K_\bullet=-\sum_{\alpha,\beta}a^\bullet_{\alpha\beta}
u^\alpha(v-\eone)^\beta\), the arguments \(u_1\) and \(1-v_1\) of the
first two selectors of \(f_{n,\bullet}\) lie within \(3/s\) of zero,
and the arguments \(v_1\) and \(1-u_1\) of the other two within \(3/s\)
of one. There
\begin{align*}
 \mathfrak K_\bullet-f_{n,\bullet}
 &=\bigl(\mathfrak K_\bullet+\mathcal T_{n,\bullet}(u,v-\eone)\bigr)
 +\bigl(\mathsf p_{2n}(u_1)\mathsf p_{2n}(1-v_1)-1\bigr)
  \mathcal T_{n,\bullet}(u,v-\eone)\\
 &\quad+\mathsf p_{2n}(v_1)\mathsf p_{2n}(1-u_1)\,
  \mathcal T_{n,\bullet}(v,u-\eone),
\end{align*}
where the first term is minus the Taylor remainder, and
\[
 |\mathsf p_{2n}(u_1)\mathsf p_{2n}(1-v_1)-1|
 \le|\mathsf p_{2n}(u_1)-1|\,|\mathsf p_{2n}(1-v_1)|
 +|\mathsf p_{2n}(1-v_1)-1|\le3(C/s)^{2n}.
\]
Hence, on this part and,
by the symmetry, on the other cross part,
\[
 |\mathfrak K_\bullet-f_{n,\bullet}|
 \le C(n+1)^5M_\bullet\Bigl(\frac Cs\Bigr)^n
 +3\Bigl(\frac Cs\Bigr)^{2n}CM_\bullet
 +\Bigl(\frac Cs\Bigr)^{4n}(n+1)^6C^nM_\bullet .
\]

On the two same-region parts, \(\mathfrak K_\bullet=0\), and in each of
the two terms of \(f_{n,\bullet}\) one selector has its argument
within \(3/s\) of one and is at most \((C/s)^{2n}\), while the other is
at most \(2\); so
\(|\mathfrak K_\bullet-f_{n,\bullet}|\le4(C/s)^{2n}(n+1)^6C^nM_\bullet\)
there. Since \((C/s)^{jn}C^n\le(C^{j+1}/s)^n\) for \(s\ge1\) and
\(j\ge1\), the bounds on the four parts combine to
\[
 \sup_{\mathcal N_s^{\C}}|\mathfrak K_\bullet-f_{n,\bullet}|
 \le C(n+1)^6M_\bullet(c/s)^n
\]
with a universal \(c\). For \(n=0\), where \(f_{0,\bullet}=0\), the
same bound holds, since \(|\mathfrak K_\bullet|\le M_\bullet\) on
\(\mathcal N_s^{\C}\) by Step~1.

Cauchy estimates then bound the derivatives in
\eqref{eq:differentiated-error}: for \((\bx,\by)\in\mathcal N_s\) and
\(|\alpha|,|\beta|\le2\), the closed unit polydisc about
\((\bx,\by)\) lies in \(\mathcal N_s^{\C}\), so
\[
 |\partial_\bx^\alpha\partial_\by^\beta
 (\mathfrak K_\bullet-f_{n,\bullet})(\bx,\by)|
 \le\alpha!\,\beta!\,
 \sup_{\mathcal N_s^{\C}}|\mathfrak K_\bullet-f_{n,\bullet}|
 \le4C(n+1)^6M_\bullet(c/s)^n .
\]
Polynomial factors in \(n\), here the factor \((n+1)^6\) from the
coefficient counts, are absorbed in \(c_{\mathrm{pol}}^n\): since
\(n+1\le2^n\), \((n+1)^6(c/s)^n\le(2^6c/s)^n\), and
\eqref{eq:differentiated-error} holds with \(c_{\mathrm{pol}}=2^6c\).
The constants in the Taylor bounds above depend only on the polydisc
radius \(r_0\), and those in the selector bounds are absolute; so \(c\)
and \(c_{\mathrm{pol}}\) depend only on \(r_0\), and in particular not on
\(\mu\) or \(s\).

\smallskip
\emph{Step 4: the midpoint coefficients.} If the polynomial kernel
\(f_{n,\bullet}\) is expressed in midpoint coordinates,
\[
 f_{n,\bullet}(\bx,\by)
 =\sum_{\alpha,\beta}c_{\alpha\beta}^{(n,\bullet)}
 \left(\frac{\bx-s\eone/2}{s}\right)^\alpha
 \left(\frac{\by-s\eone/2}{s}\right)^\beta,
\]
then, for a universal finite \(K_{\mathrm{mid}}\),
\begin{equation}
 \sum_{\alpha,\beta}
 |c_{\alpha\beta}^{(n,\bullet)}|
 \le CM_\bullet K_{\mathrm{mid}}^n.
 \label{eq:coefficient-bound}
\end{equation}

Write \(u'=u-\eone/2\) and
\(v'=v-\eone/2\) for the midpoint coordinates; then \(u=u'+\eone/2\)
and \(v-\eone=v'-\eone/2\), so only the first coordinates are shifted,
by \(\pm1/2\), and the coefficient sums below are taken in
\((u',v')\). By \(|a^\bullet_{\alpha\beta}|\le M_\bullet
r_0^{-|\alpha|-|\beta|}\) and the count of monomials, the coefficient
sum of \(\mathcal T_{n,\bullet}\) in \((u,v-\eone)\) is at most
\(M_\bullet\sum_{k<n}(k+1)^5r_0^{-k}\le(n+1)^6r_0^{-n}M_\bullet\); the
shift multiplies it by at most \((3/2)^{n-1}\), since a monomial of
degree \(k\le n-1\) acquires at most \(k\) factors \(u_1'+1/2\) or
\(v_1'-1/2\). The two selectors become \(\mathsf p_{2n}(u_1'+1/2)\) and
\(\mathsf p_{2n}(1/2-v_1')\), with coefficient sums at most \(C^n\) by
the bound after \eqref{eq:selectors}. The transposed term exchanges
\(u'\) and \(v'\) and has the same bound. The product of these bounds
for the Taylor polynomial, the shift and the two selectors gives
\(K_{\mathrm{mid}}^n\):
\[
 \sum_{\alpha,\beta}|c^{(n,\bullet)}_{\alpha\beta}|
 \le2\,C^n\,C^n\Bigl(\frac32\Bigr)^{n-1}(n+1)^6r_0^{-n}M_\bullet
 \le2(n+1)^6\Bigl(\frac C{r_0}\Bigr)^nM_\bullet
 \le2M_\bullet K_{\mathrm{mid}}^n
\]
with \(K_{\mathrm{mid}}=2^6C/r_0\), by \((n+1)^6\le2^{6n}\).

\smallskip
\emph{Step 5: the operator-norm error.} By Steps~2 and~3, for
\(n\ge0\) and smooth real sources supported in \(U_s\), the form of
\(\JointD_{\mu,s}-\mathsf F_n\) is the sum of the field-source form of
the kernel \(\mathfrak K_f-f_{n,f}\) and the momentum-source form of
\(\mathfrak K_g-f_{n,g}\); at \(n=0\), where \(\mathsf F_0=0\), it is
the form \eqref{eq:D-joint} itself. By \eqref{eq:differentiated-error},
the smooth-kernel estimate above applies to these kernels with
\(M=CM_\bullet(c_{\mathrm{pol}}/s)^n\); together with
\(\max\{1,\mu\}\le1+\mu\le1+\xi\) for \(s\ge1\), it now gives
\begin{align*}
 |\ip{\sx}{(\JointD_{\mu,s}-\mathsf F_n)\sy}|
 &\le C\bigl(\max\{1,\mu\}M_f+M_g\bigr)
   (c_{\mathrm{pol}}/s)^n\norm{\sx}\norm{\sy}\\
 &\le C\bigl(s^{-2}(1+\xi)^{3/2}+s^{-4}(1+\xi)^{3/2}\bigr)e^{-\xi/2}
   (c_{\mathrm{pol}}/s)^n\norm{\sx}\norm{\sy}
\end{align*}
on smooth real sources, and by density on \(\mathcal X_{\mu,U_s}\).
The forms are real and symmetric, so the operators they define on
\(\Hh_{\mu,U_s}\) are the complexifications of self-adjoint operators
on \(\mathcal X_{\mu,U_s}\), with the same norms. With
\(s^{-4}\le s^{-2}\) and \((1+\xi)^{3/2}\le(1+\xi)^2\), the bound of the
forms gives
\begin{equation}
 \norm{\JointD_{\mu,s}-\mathsf F_n}
 \le Cs^{-2}(1+\xi)^2e^{-\xi/2}(c_{\mathrm{pol}}/s)^n.
 \label{eq:polynomial-error}
\end{equation}
At \(n=0\), \eqref{eq:polynomial-error} shows that the form in
\eqref{eq:D-joint} is bounded. It also gives
\(\mathsf F_n\to\JointD_{\mu,s}\) in operator norm, because
\(S_{\mathrm{pol}}>c_{\mathrm{pol}}\) (Step~3),
so that \((c_{\mathrm{pol}}/s)^n\to0\).

\smallskip
\emph{Step 6: the unweighted increments.} The kernel of
\(\mathsf G_n=\mathsf F_{n+1}-\mathsf F_n\) is
\(f_{n+1,\bullet}-f_{n,\bullet}\), of degree at most \(9(n+1)\); as in
Step~2, the form of \(\mathsf G_n\) depends on \(\sx\) only through the
moments \(\ell^f_\alpha(\sx)\) and \(\ell^g_\alpha(\sx)\) with
\(|\alpha|\le9(n+1)\), so
\[
 \rank\mathsf G_n\le2\binom{9n+12}{3}\le C(n+1)^3 .
\]
The bounds \eqref{eq:polynomial-error} at \(n\) and \(n+1\),
multiplied by the rank of
\(\mathsf G_n=(\JointD_{\mu,s}-\mathsf F_n)-(\JointD_{\mu,s}-\mathsf F_{n+1})\),
give \eqref{eq:increment-unweighted} for every \(n\ge0\): the trace norm
of a finite-rank
operator is at most its rank times its operator norm, and
\(c_{\mathrm{pol}}/s\le1\), so
\[
 \norm{\mathsf G_n}_1
 \le\rank(\mathsf G_n)\bigl(\norm{\JointD_{\mu,s}-\mathsf F_n}
 +\norm{\JointD_{\mu,s}-\mathsf F_{n+1}}\bigr)
 \le C(n+1)^3\cdot2Cs^{-2}(1+\xi)^2e^{-\xi/2}(c_{\mathrm{pol}}/s)^n .
\]

\smallskip
\emph{Step 7: the weighted increments.} In the midpoint expansion of
Step~4, the form of \(\mathsf G_n\) is a combination of products of the
moments \(s^{-|\alpha|}\ell^f_\alpha\) and \(s^{-|\alpha|}\ell^g_\alpha\),
centred at \(\bx_0=s\eone/2\) and taken over the whole union. The union
lies in the ball \(B_T(\bx_0)\) of radius \(T=s/2+1\), and
\(s/2\le T\le s\) for \(s\ge2\). Lemma~\ref{lem:moments}, at \(b=1/2\),
applies to each of these moments, for \(U=U_s\subset B_T(\bx_0)\) and
its reserve \(\JointRes_{\mu,s}\).
Since \(T/s\le1\), every normalized degree costs at most
\(B_{1/2}=5\Lambda_{1/2}\):
\[
 \norm{s^{-|\alpha|}\ell^f_\alpha\JointRes_{\mu,s}^{-1/2}}
 \le M_{1/2}\frac{T}{1+\mu T}e^{\Lambda_0\mu T}B_{1/2}^{|\alpha|},
 \qquad
 \norm{s^{-|\alpha|}\ell^g_\alpha\JointRes_{\mu,s}^{-1/2}}
 \le M_{1/2}\frac{T^2}{(1+\mu T)^2}e^{\Lambda_0\mu T}B_{1/2}^{|\alpha|}.
\]

The Riesz
representation gives vectors \(\psi^\bullet_\alpha\), real because
\(\JointRes_{\mu,s}\) and the moments are real, with
\[
 s^{-|\alpha|}\ell^\bullet_\alpha(\sx)
 =\ip{\psi^\bullet_\alpha}{\JointRes_{\mu,s}^{1/2}\sx}
\]
and norms bounded by the right-hand sides above.

The kernel \(f_{n+1,\bullet}-f_{n,\bullet}\) of \(\mathsf G_n\) has the
midpoint coefficients
\(c^{(n+1,\bullet)}_{\alpha\beta}-c^{(n,\bullet)}_{\alpha\beta}\), which
vanish unless \(|\alpha|+|\beta|\le9(n+1)\) and, by
\eqref{eq:coefficient-bound}, have the coefficient sum at most
\(CM_\bullet(1+K_{\mathrm{mid}})K_{\mathrm{mid}}^n\). Put
\[
 B_n=\sum_{\bullet\in\{f,g\}}\sum_{\alpha,\beta}
 \bigl(c^{(n+1,\bullet)}_{\alpha\beta}-c^{(n,\bullet)}_{\alpha\beta}\bigr)
 \,|\psi^\bullet_\alpha\rangle\langle\psi^\bullet_\beta| .
\]
For smooth real sources,
\(\ip{\sx}{\JointRes_{\mu,s}^{1/2}B_n\JointRes_{\mu,s}^{1/2}\sy}\) is the
form of \(\mathsf G_n\) written in the midpoint moments, so
\(\mathsf G_n=\JointRes_{\mu,s}^{1/2}B_n\JointRes_{\mu,s}^{1/2}\) as
bounded operators. The two weighted moment norms of each rank-one
monomial give its trace norm
\(\lVert\psi^\bullet_\alpha\rVert\,\lVert\psi^\bullet_\beta\rVert\), with
\(B_{1/2}^{|\alpha|+|\beta|}\le B_{1/2}^{9(n+1)}\); multiplied by the
coefficient sum, these trace norms give
\[
 \norm{B_n}_1
 \le C(1+K_{\mathrm{mid}})K_{\mathrm{mid}}^nB_{1/2}^{9(n+1)}
 M_{1/2}^2e^{2\Lambda_0\mu T}
 \biggl[M_f\frac{T^2}{(1+\mu T)^2}+M_g\frac{T^4}{(1+\mu T)^4}\biggr].
\]

The dimensional prefactors satisfy, for \(s/2\le T\le s\),
\[
 \begin{split}
 s^{-2}\frac{T^2}{(1+\mu T)^2}(1+\xi)^{1/2}&\le C,\\
 s^{-4}\frac{T^4}{(1+\mu T)^4}(1+\xi)^{3/2}&\le C,
 \end{split}
\]
because \(T\le s\) and \(1+\mu T\ge1+\xi/2\ge(1+\xi)/2\). With
\(M_f\) and \(M_g\) from Step~1, the bracket is therefore at most
\(Ce^{-\xi/2}\), and
\[
 \norm{B_n}_1
 \le C(1+K_{\mathrm{mid}})B_{1/2}^9M_{1/2}^2\,
 e^{-\xi/2}e^{2\Lambda_0\mu T}\bigl(K_{\mathrm{mid}}B_{1/2}^9\bigr)^n .
\]
The bound on \(\norm{B_n}_1\) is thus
\eqref{eq:increment-weighted} with
\(K_{\mathrm{pol}}=K_{\mathrm{mid}}B_{1/2}^9\) and
\(C_2=C(1+K_{\mathrm{mid}})B_{1/2}^9M_{1/2}^2\). The operator
\(B_n\) is real and self-adjoint: the
midpoint coefficients are real and symmetric in \((\alpha,\beta)\),
because \(f_{n,\bullet}\) is real and symmetric in \((\bx,\by)\), and
the vectors \(\psi^\bullet_\alpha\) are real.
\end{proof}

\subsection{Interpolation at a small joint reserve power}

The constants in \eqref{eq:Lambda} fix the parameters
\begin{equation}
 \Lambda_0=2^{15},\qquad
 \theta=\frac1{4\Lambda_0}=2^{-17},\qquad
 a=\frac{\theta}{2}=2^{-18},\qquad
 \gamma=2(1-\theta)=2(1-2^{-17}).
 \label{eq:reverse-exponents}
\end{equation}
In particular, \(0<\theta\le1/64\), \(0<a\le1/8\), and
\(\gamma>0\). The bound \(a\le1/8\) places \(a\) in the range
\(0<a\le1/2\) of Lemmas~\ref{lem:thermal-comparison}
and~\ref{lem:compression}, and the order \(r_-=1/(1-a)\) of
Section~\ref{sec:completion} then satisfies \(1<r_-<2\), as required in
Theorem~\ref{thm:probability}.

The factor \(16\) in the moment enlargement
\eqref{eq:Lambda} makes \(\Lambda_0=16\lambda_{1/4}=2^{15}\) a power of two,
and with it \(\theta\) and \(a\). The interpolation weight
\(\theta=1/(4\Lambda_0)\) is the value at which the growth factor
\(e^{2\Lambda_0\mu T}\) of \eqref{eq:increment-weighted}, raised to
the power \(\theta\), becomes \(e^{\mu T/2}\), which the decay
\(e^{-\xi/2}\) of \eqref{eq:increment-unweighted} dominates for
\(T\le3s/4\); the reserve power \(a=\theta/2\) arises because the
interpolating family in the proof of Theorem~\ref{thm:reverse} is
\(\JointRes_{\mu,s}^{(1-z)/2}B_n\JointRes_{\mu,s}^{(1-z)/2}\).

\begin{theorem}[Jointly weighted reverse trace estimate]\label{thm:reverse}
There are finite \(S_{\mathrm{rev}},\mathcal J_0\), independent of \(\mu\),
such that for every \(\mu>0\) and \(s\ge S_{\mathrm{rev}}\) the
product-minus-joint covariance operator \(\JointD_{\mu,s}\) of
\eqref{eq:D-joint} factors, with \(a,\gamma\) as in
\eqref{eq:reverse-exponents}, as
\begin{equation}
 \JointD_{\mu,s}
 =\JointRes_{\mu,s}^{a}Y_{\mu,s}\JointRes_{\mu,s}^{a},
 \qquad
 Y_{\mu,s}=Y_{\mu,s}^*\in\Sc_1,\qquad
 \norm{Y_{\mu,s}}_1\le \mathcal J_0s^{-\gamma}.
 \label{eq:reverse-trace}
\end{equation}
Here \(\JointRes_{\mu,s}\) is the uncertainty reserve of the joint
vacuum covariance norm on the union \(U_s\) of the two unit balls.
\end{theorem}

\begin{proof}
\emph{Step 1: interpolation for each increment.} For each increment
\(\mathsf G_n\) of Lemma~\ref{lem:polynomial}, let
\(Z_n=\JointRes_{\mu,s}^{(1-\theta)/2}B_n\JointRes_{\mu,s}^{(1-\theta)/2}\),
with \(B_n\) from \eqref{eq:increment-weighted}. Trace-norm interpolation
gives
\begin{equation}
 \mathsf G_n
 =\JointRes_{\mu,s}^{a}Z_n\JointRes_{\mu,s}^{a},
 \qquad
 \norm{Z_n}_1
 \le\norm{\mathsf G_n}_1^{1-\theta}\norm{B_n}_1^\theta.
 \label{eq:interpolation}
\end{equation}
Indeed, the analytic family
\[
 z\longmapsto
 \JointRes_{\mu,s}^{(1-z)/2}
 B_n
 \JointRes_{\mu,s}^{(1-z)/2},
 \qquad 0\le\Rea z\le1,
\]
has boundary trace norms \(\norm{\mathsf G_n}_1\) and
\(\norm{B_n}_1\). Since
\(0<\JointRes_{\mu,s}\le I\), the powers \(\JointRes_{\mu,s}^w\) with
\(\Rea w\ge0\), given by the spectral calculus on \((0,1]\), are
contractions, and unitary for \(\Rea w=0\). For \(\Rea z=0\),
\(\JointRes_{\mu,s}^{(1-z)/2}
=\JointRes_{\mu,s}^{1/2}\JointRes_{\mu,s}^{-z/2}\), and for
\(\Rea z=1\), \(\JointRes_{\mu,s}^{(1-z)/2}\) is unitary; with
\eqref{eq:increment-weighted},
\[
 \norm{\JointRes_{\mu,s}^{(1-z)/2}B_n\JointRes_{\mu,s}^{(1-z)/2}}_1
 =\begin{cases}
 \norm{\JointRes_{\mu,s}^{-z/2}\mathsf G_n\JointRes_{\mu,s}^{-z/2}}_1
 =\norm{\mathsf G_n}_1, & \Rea z=0,\\
 \norm{B_n}_1, & \Rea z=1.
 \end{cases}
\]
At \(z=\theta\), the value of the family is \(Z_n\), and
multiplication by \(\JointRes_{\mu,s}^{\theta/2}=\JointRes_{\mu,s}^a\)
on both sides recovers
\(\JointRes_{\mu,s}^{1/2}B_n\JointRes_{\mu,s}^{1/2}=\mathsf G_n\).

Spectral cutoffs and their strong limits prove the continuity and
holomorphy of the family, also when zero lies in the spectrum of
\(\JointRes_{\mu,s}\). Applied to any fixed vector,
\(w\mapsto\JointRes_{\mu,s}^w\) is continuous on \(\Rea w\ge0\) and
holomorphic on \(\Rea w>0\). Since
\(B_n\) is a finite combination of the rank-one operators
\(|\psi^\bullet_\alpha\rangle\langle\psi^\bullet_\beta|\) of the proof
of Lemma~\ref{lem:polynomial}, the family is the same combination of
\(|\JointRes_{\mu,s}^{(1-z)/2}\psi^\bullet_\alpha\rangle
\langle\JointRes_{\mu,s}^{(1-\bar z)/2}\psi^\bullet_\beta|\). It is
therefore continuous in trace norm on the closed strip and holomorphic
inside, with trace norm at most \(\norm{B_n}_1\).

\smallskip
\emph{Step 2: the exponent.} Inserting
\eqref{eq:increment-unweighted} and \eqref{eq:increment-weighted} into
the bound of \eqref{eq:interpolation} gives
\begin{align*}
 \norm{Z_n}_1
 &\le\Bigl(C_1(n+1)^3s^{-2}(1+\xi)^2e^{-\xi/2}
   (c_{\mathrm{pol}}/s)^n\Bigr)^{1-\theta}
   \Bigl(C_2e^{-\xi/2}e^{2\Lambda_0\mu T}K_{\mathrm{pol}}^n\Bigr)^{\theta}\\
 &=C_1^{1-\theta}C_2^\theta\,s^{-2(1-\theta)}(1+\xi)^{2(1-\theta)}
   e^{-\xi/2+2\Lambda_0\theta\mu T}(n+1)^{3(1-\theta)}
   \left[(c_{\mathrm{pol}}/s)^{1-\theta}K_{\mathrm{pol}}^\theta\right]^n,
\end{align*}
since \(-(1-\theta)\xi/2-\theta\xi/2=-\xi/2\). Here
\(s^{-2(1-\theta)}=s^{-\gamma}\) and \(2\Lambda_0\theta=1/2\) by
\eqref{eq:reverse-exponents}. For \(s\ge4\), the enclosing radius
satisfies \(T\le3s/4\), so \(\mu T/2\le3\xi/8\) and
\[
 -\xi/2+2\Lambda_0\theta\mu T
 =-\xi/2+\mu T/2\le-\xi/8.
\]
With this bound on the exponent,
the estimate of \(\norm{Z_n}_1\) above becomes
\begin{equation}
 \begin{split}
 \norm{Z_n}_1\le{}&
 C_1^{1-\theta}C_2^\theta
 s^{-\gamma}(1+\xi)^{2(1-\theta)}e^{-\xi/8}
 (n+1)^{3(1-\theta)}\\
 &\quad{}\times
 \left[(c_{\mathrm{pol}}/s)^{1-\theta}K_{\mathrm{pol}}^\theta\right]^n.
 \end{split}
 \label{eq:Z-bound}
\end{equation}

\smallskip
\emph{Step 3: summation.} The threshold \(S_{\mathrm{rev}}\) is chosen
so that every \(s\ge S_{\mathrm{rev}}\) satisfies \(s\ge S_{\mathrm{pol}}\),
\(s\ge5\), and
\[
 (c_{\mathrm{pol}}/s)^{1-\theta}K_{\mathrm{pol}}^\theta\le\tfrac12.
\]
With \((n+1)^{3(1-\theta)}\le(n+1)^3\) and
\((1+\xi)^{2(1-\theta)}\le(1+\xi)^2\), \eqref{eq:Z-bound} then gives
\[
 \sum_{n\ge0}\norm{Z_n}_1
 \le C_1^{1-\theta}C_2^\theta s^{-\gamma}
 \Bigl(\sup_{\xi\ge0}(1+\xi)^2e^{-\xi/8}\Bigr)
 \sum_{n\ge0}(n+1)^3\,2^{-n}.
\]

The bounds
\[
 \sum_{n\ge0}(n+1)^3\,2^{-n}=52,\qquad
 \sup_{\xi\ge0}(1+\xi)^2e^{-\xi/8}\le256
\]
follow from \(\sum_{n\ge0}(n+1)^3x^n=(1+4x+x^2)/(1-x)^4\) at
\(x=1/2\), and from the maximum of \((1+\xi)^2e^{-\xi/8}\) at
\(\xi=15\), where it equals \(256e^{-15/8}\). They permit
\begin{equation}
 \mathcal J_0=52\cdot256\,C_1^{1-\theta}C_2^\theta.
 \label{eq:J0}
\end{equation}
Thus \(Y_{\mu,s}=\sum_{n\ge0}Z_n\) converges in trace norm and
satisfies the trace bound
\(\norm{Y_{\mu,s}}_1\le\mathcal J_0s^{-\gamma}\).

\smallskip
\emph{Step 4: the factorization.} Multiplication by the bounded powers
\(\JointRes_{\mu,s}^a\) on both sides is continuous in trace norm, and
it identifies \(\JointRes_{\mu,s}^aY_{\mu,s}\JointRes_{\mu,s}^a\) with
\(\sum_n\mathsf G_n=\JointD_{\mu,s}\):
\[
 \JointRes_{\mu,s}^{a}Y_{\mu,s}\JointRes_{\mu,s}^{a}
 =\sum_{n\ge0}\JointRes_{\mu,s}^{a}Z_n\JointRes_{\mu,s}^{a}
 \overset{\eqref{eq:interpolation}}{=}\sum_{n\ge0}\mathsf G_n
 =\lim_{N\to\infty}\mathsf F_{N+1}
 \overset{\text{Lem.~\ref{lem:polynomial}}}{=}\JointD_{\mu,s},
\]
since \(\sum_{n=0}^N\mathsf G_n=\mathsf F_{N+1}-\mathsf F_0
=\mathsf F_{N+1}\) and \(\mathsf F_n\to\JointD_{\mu,s}\) in operator
norm. Each \(Z_n\) is self-adjoint, because \(B_n\) is, and
self-adjointness passes to the trace-norm limit.
\end{proof}

\section{Completion of the regional Bell bound}\label{sec:completion}

\subsection{Comparisons on every finite source system}

For the parameters \eqref{eq:reverse-exponents}, we fix
\begin{equation}
 r_-=\frac1{1-a}=\frac{2^{18}}{2^{18}-1},\qquad
 r_+=\left\lceil\frac4a\right\rceil=2^{20},\qquad
 b=1-\frac1{r_+}=1-2^{-20}.
 \label{eq:orders}
\end{equation}
Here \(r_+\) is an integer, as required in
\eqref{eq:high-compression}. Then \(1<r_-<2\), \(r_+\ge32\),
\(ar_+=4\), and, by \eqref{eq:probability-exponents},
\[
 \eta=1-\frac{1-b}{ab^2}
 =1-\frac1{4b^2}>0.
\]

The reverse exponent \(a\) constrains the forward order through the
condition \(\eta>0\): since \(1-b=1/r_+\), positivity of \(\eta\)
requires \(r_+>1/(ab^2)\), so the small value of \(a\) forces a large
order; the choice \(r_+=4/a\) gives the exponent \(b=1-a/4\) close to
\(1\) and, through \eqref{eq:Cprob}, the size of \(C_*\).
Theorem~\ref{thm:forward} is used at this fixed \(b\).

For arbitrary finite canonical source systems \(F_A,F_B\) in the two
unit balls, the symplectic form \eqref{eq:symplectic} vanishes between
a source in \(F_A\) and a source in \(F_B\), because the balls are
disjoint. On
\(F=F_A\oplus F_B\), the form is therefore the direct sum of its
nondegenerate restrictions to \(F_A\) and \(F_B\). Thus \(F\) is a finite
canonical source system, and, by the Weyl relations, the Weyl operators
of \(F_A\) commute with those of \(F_B\).

Let \(\rho_F\) be the joint Gaussian vacuum state of \(F\), and let
\[
 \sigma_F=(\rho_F)_{F_A}\otimes(\rho_F)_{F_B}
\]
be the product of its marginals. Both states are centred Gaussian
states on \(F\). They are faithful by \eqref{eq:strict-uncertainty},
applied to \(F\) for \(\rho_F\) and to \(F_A\) and \(F_B\) for the two
factors of \(\sigma_F\).

The thermal comparison constants \eqref{eq:scalar-reserve}, the
constant \(F_b\) of Theorem~\ref{thm:forward} and the constant
\(\mathcal J_0\) of Theorem~\ref{thm:reverse} give
\begin{equation}
 A_+=c_b^{-1}F_b=(r_+-1)F_b,\qquad
 A_-=c_a^{-1}\mathcal J_0=2^{2a-1}\mathcal J_0.
 \label{eq:Aconstants}
\end{equation}
The comparison envelopes are
\begin{equation}
 \varepsilon_+(s)=A_+s^{-2},\qquad
 \varepsilon_-(s)=A_-s^{-\gamma}.
 \label{eq:envelopes}
\end{equation}

For the forward comparison, let
\(\iota_A\colon F_A\to\mathcal X_{\mu,B_1(0)}\) and
\(\iota_B\colon F_B\to\mathcal X_{\mu,B_1(s\eone)}\) be the inclusions;
they are isometries when \(F_A\) and \(F_B\) carry the marginal
covariance inner products \(\vv_{\mu,B_1(0)}\) and
\(\vv_{\mu,B_1(s\eone)}\). On \(F\) we use the orthogonal sum of these
two inner products, which is the product covariance
\(\vv_{\mathrm{prod},\mu,s}\), and coordinates orthonormal for it. In
these product-normalized coordinates, \(V_{\sigma_F}=I\), because
\(\sigma_F\) has the product covariance, and the definition of the
cross-covariance operator \(C_{\mu,s}\) in Section~\ref{sec:forward}
gives
\[
 V_{\rho_F}-V_{\sigma_F}
 =\begin{pmatrix}
  0&\iota_A^*C_{\mu,s}\iota_B\\
  (\iota_A^*C_{\mu,s}\iota_B)^*&0
 \end{pmatrix}.
\]
Let \(\mathsf Z_{+,F}\)
denote the operator \(\mathsf Z\) of Theorem~\ref{thm:gaussian} for
\((\rho,\sigma,r)=(\rho_F,\sigma_F,r_+)\) in these coordinates.
By Step~1 of the proof of Theorem~\ref{thm:gaussian}, the norms of
\(\mathsf Z\) may be computed in these coordinates, and in the joint
coordinates below.

Here \(\tau=(r_+-1)/r_+=b\). Since \(\sigma_F\) is a product state, its
thermal power \(\sigma_F^b\) is the product of the thermal powers of its
factors, so the thermal increment of \(V_{\sigma_F}=I\) is block
diagonal, with blocks \(H_{A,F}=\Therm_b(I_{F_A})\) and
\(H_{B,F}=\Therm_b(I_{F_B})\). Hence
\[
 \mathsf Z_{+,F}
 =\begin{pmatrix}0&\mathsf Z_{AB}\\ \mathsf Z_{AB}^*&0\end{pmatrix},
 \qquad
 \mathsf Z_{AB}=H_{A,F}^{-1/2}\iota_A^*C_{\mu,s}\iota_BH_{B,F}^{-1/2}.
\]

The high-order compression \eqref{eq:high-compression}, applied in each
ball with the integer \(r_+\), gives
\(H_{A,F}\ge c_b\iota_A^*\Res_A^{2b}\iota_A\) and
\(H_{B,F}\ge c_b\iota_B^*\Res_B^{2b}\iota_B\), the hypotheses of
\eqref{eq:trace-compression}. For \(s\ge S_b\), the trace estimate
\eqref{eq:forward-trace} at this \(b\) then gives
\begin{align*}
 \norm{\mathsf Z_{AB}}_1
 &\overset{\eqref{eq:forward-trace}}{=}
 \norm{H_{A,F}^{-1/2}\iota_A^*\Res_A^bX_{b,\mu,s}\Res_B^b
 \iota_BH_{B,F}^{-1/2}}_1\\
 &\overset{\eqref{eq:trace-compression}}{\le}
 c_b^{-1}\norm{X_{b,\mu,s}}_1
 \overset{\eqref{eq:forward-trace}}{\le}
 c_b^{-1}F_bs^{-2}
 \overset{\eqref{eq:Aconstants},\,\eqref{eq:envelopes}}{=}
 \varepsilon_+(s).
\end{align*}
The eigenvalues of \(\mathsf Z_{+,F}\) are plus and minus the singular
values of \(\mathsf Z_{AB}\), together with zeros. Hence
\(\norm{\mathsf Z_{+,F}}=\norm{\mathsf Z_{AB}}\le\norm{\mathsf Z_{AB}}_1\)
and \(\norm{\mathsf Z_{+,F}}_1=2\norm{\mathsf Z_{AB}}_1\), which gives
\begin{equation}
 \norm{\mathsf Z_{+,F}}\le\varepsilon_+(s),\qquad
 \norm{\mathsf Z_{+,F}}_1\le2\varepsilon_+(s).
 \label{eq:finite-forward}
\end{equation}

The number \(k_r\) of Theorem~\ref{thm:gaussian} is here
\(k_{r_+}=\max\{2,1/b\}=2\), since \(b\ge1/2\). If
\(\varepsilon_+(s)\le1/8\), then \eqref{eq:finite-forward} gives the
hypothesis \(\norm{\mathsf Z_{+,F}}\le1/8=(4k_{r_+})^{-1}\) of
Theorem~\ref{thm:gaussian}, and \eqref{eq:gaussian-bound} gives
\(\log\Qtil_{r_+}(\rho_F\|\sigma_F)\le2r_+k_{r_+}\norm{\mathsf Z_{+,F}}_1
\le2r_+\cdot2\cdot2\varepsilon_+(s)\), that is,
\begin{equation}
 \log\Qtil_{r_+}(\rho_F\|\sigma_F)
 \le8r_+\varepsilon_+(s).
 \label{eq:Qforward}
\end{equation}

For the reverse comparison, the combined finite source space \(F\) is
embedded isometrically in the joint covariance norm by the inclusion
\(\iota_F\colon F\to\mathcal X_{\mu,U_s}\), which preserves \(\varsigma\).
The change to coordinates orthonormal for this norm, the joint
coordinates, acts on the same finite system and may mix the
coordinates of \(F_A\) and \(F_B\). In the joint coordinates,
\begin{equation}
 \begin{gathered}
 V_{\rho_F}=I_F,\qquad
 V_{\sigma_F}-V_{\rho_F}
 =\iota_F^*\JointD_{\mu,s}\iota_F,\\
 \Therm_a(I_F)\ge c_a\Res_F^{2a}
 \ge c_a\iota_F^*\JointRes_{\mu,s}^{2a}\iota_F.
 \end{gathered}
 \label{eq:joint-compression}
\end{equation}
In the first line, \(V_{\rho_F}=I_F\) because \(\iota_F\) is isometric
for the joint covariance, and the difference is given by
\eqref{eq:D-joint} because \(\sigma_F\) has the product covariance
\(\vv_{\mathrm{prod},\mu,s}\). In the last line,
\(\Res_F\) is the reserve of
Lemma~\ref{lem:compression} for \(\iota=\iota_F\); the line uses
\eqref{eq:scalar-reserve} on the finite space and
\eqref{eq:low-compression}, with \(a\le1/2\). Theorem~\ref{thm:gaussian} is now applied
with the roles of the two states exchanged:
\(\mathsf Z_{-,F}\) denotes the operator \(\mathsf Z\) of
Theorem~\ref{thm:gaussian} for
\((\rho,\sigma,r)=(\sigma_F,\rho_F,r_-)\) in the joint coordinates
\eqref{eq:joint-compression}.

Here \(\tau=(r_--1)/r_-=a\). For \(s\ge S_{\mathrm{rev}}\),
Theorem~\ref{thm:reverse} gives the factorization
\eqref{eq:reverse-trace}, and with the first line of
\eqref{eq:joint-compression},
\begin{align*}
 \mathsf Z_{-,F}
 &=\Therm_a(I_F)^{-1/2}\iota_F^*\JointD_{\mu,s}\iota_F\Therm_a(I_F)^{-1/2}
 \overset{\eqref{eq:reverse-trace}}{=}\mathsf N^*Y_{\mu,s}\mathsf N,\\
 \mathsf N&=\JointRes_{\mu,s}^a\iota_F\Therm_a(I_F)^{-1/2}.
\end{align*}
The last line of \eqref{eq:joint-compression} is the hypothesis of
the one-space normalization bound stated before
\eqref{eq:trace-compression}, with \(H_F=\Therm_a(I_F)\) and
\(\iota=\iota_F\). That bound gives \(\norm{\mathsf N}^2\le c_a^{-1}\),
and therefore
\[
 \norm{\mathsf Z_{-,F}}_1
 \le\norm{\mathsf N}^2\norm{Y_{\mu,s}}_1
 \overset{\eqref{eq:reverse-trace}}{\le}c_a^{-1}\mathcal J_0s^{-\gamma}
 \overset{\eqref{eq:Aconstants},\,\eqref{eq:envelopes}}{=}\varepsilon_-(s),
\]
that is,
\begin{equation}
 \norm{\mathsf Z_{-,F}}_1\le\varepsilon_-(s).
 \label{eq:finite-reverse}
\end{equation}

Here \(k_{r_-}=\max\{2,1/a\}=1/a\), since \(a\le1/2\). If
\(\varepsilon_-(s)\le a/4\), then \eqref{eq:finite-reverse} gives
\(\norm{\mathsf Z_{-,F}}\le\norm{\mathsf Z_{-,F}}_1\le a/4=(4k_{r_-})^{-1}\),
the hypothesis of Theorem~\ref{thm:gaussian}, and
\eqref{eq:gaussian-bound} gives
\(\log\Qtil_{r_-}(\sigma_F\|\rho_F)\le2r_-k_{r_-}\norm{\mathsf Z_{-,F}}_1
\le2r_-a^{-1}\varepsilon_-(s)\), that is,
\begin{equation}
 \log\Qtil_{r_-}(\sigma_F\|\rho_F)
 \le\frac{2r_-}{a}\varepsilon_-(s).
 \label{eq:Qreverse}
\end{equation}
All constants in \eqref{eq:finite-forward}--\eqref{eq:Qreverse}
are independent of the finite source systems and their symplectic
spectra.

\subsection{Regional events and bounded observables}

Let \(E\) and \(P\) be effects of the von Neumann algebras generated by
the Weyl operators of \(F_A\) and of \(F_B\), and write
\(e_A=\omega(E)\), \(e_B=\omega(P)\) and \(w=\omega(EP)\). The two
algebras commute, so the product
\(EP\) is an effect. Its expectations in the two states are
\(\rho_F(EP)=w\), since \(\rho_F\) is the vacuum state on \(F\), and
\(\sigma_F(EP)=e_Ae_B\), since \(\sigma_F\) is the product of the
marginals of \(\rho_F\).

When \eqref{eq:Qforward} and \eqref{eq:Qreverse} hold, the two Rényi
quantities there are finite, and the effect bound
\eqref{eq:effect-renyi} of Lemma~\ref{lem:renyi-consequences} with
\(M=EP\), applied to \((\rho,\sigma,r)=(\rho_F,\sigma_F,r_+)\) and to
\((\sigma_F,\rho_F,r_-)\), gives
\[
 w^{r_+}\le\Qtil_{r_+}(\rho_F\|\sigma_F)\,(e_Ae_B)^{r_+-1},
 \qquad
 (e_Ae_B)^{r_-}\le\Qtil_{r_-}(\sigma_F\|\rho_F)\,w^{r_--1}.
\]
Bounding the two Rényi quantities by \eqref{eq:Qforward} and
\eqref{eq:Qreverse}, we obtain
\begin{align}
 w^{r_+}
 &\le e^{8r_+\varepsilon_+(s)}(e_Ae_B)^{r_+-1},
 \label{eq:forward-events}\\
 (e_Ae_B)^{r_-}
 &\le e^{2r_-\varepsilon_-(s)/a}w^{r_--1}.
 \label{eq:reverse-events}
\end{align}

Since \eqref{eq:forward-events} and \eqref{eq:reverse-events} contain no
denominators, they have the form \(\Phi(e_A,e_B,w)\le0\) with
\(\Phi\) continuous on \([0,1]^3\), zero probabilities included, because
\(r_+-1>0\) and \(r_--1>0\); their constants depend only on \(s\).
Lemma~\ref{lem:exhaustion} therefore passes these continuous
inequalities to all effects in the two regional algebras. The suprema
in \eqref{eq:K} range over projections with
\(e_Ae_B>0\). For such projections \(E\) and \(P\),
\eqref{eq:forward-events} bounds the quotient in \(K_{r_+,+}\) by
\(e^{8r_+\varepsilon_+(s)}\), and
\eqref{eq:reverse-events} forces \(w>0\) and bounds the quotient in
\(K_{r_-,-}\) by \(e^{2r_-\varepsilon_-(s)/a}\). Hence the constants \(K_{r_+,+}(\mu,s)\) and
\(K_{r_-,-}(\mu,s)\) of \eqref{eq:K} for the two regional algebras at mass
\(\mu\) and separation \(s\) satisfy
\begin{equation}
 \log K_{r_+,+}(\mu,s)\le8r_+\varepsilon_+(s),\qquad
 \log K_{r_-,-}(\mu,s)\le\frac{2r_-}{a}\varepsilon_-(s).
 \label{eq:full-K}
\end{equation}

The trace-distance conclusion \eqref{eq:trace-distance}, which holds
for \(r_+\ge2\), and \eqref{eq:Qforward} give
\[
 \norm{\rho_F-\sigma_F}_1
 \overset{\eqref{eq:trace-distance}}{\le}
 \sqrt{\Qtil_{r_+}(\rho_F\|\sigma_F)^{1/(r_+-1)}-1}
 \overset{\eqref{eq:Qforward}}{\le}
 \sqrt{\exp\Bigl(\frac{8r_+\varepsilon_+(s)}{r_+-1}\Bigr)-1},
\]
and, since \(r_+/(r_+-1)=1/b\),
\begin{equation}
 \norm{\rho_F-\sigma_F}_1
 \le\sqrt{\exp(8\varepsilon_+(s)/b)-1}.
 \label{eq:finite-distance}
\end{equation}
For self-adjoint contractions \(A\) and \(B\) of the von Neumann
algebras generated by the Weyl operators of \(F_A\) and of \(F_B\), one
has \(\omega(AB)=\rho_F(AB)\), \(\omega(A)\omega(B)=\sigma_F(AB)\) and
\(\norm{AB}\le1\), so that
\[
 |\omega(AB)-\omega(A)\omega(B)|
 =\lvert\Tr((\rho_F-\sigma_F)AB)\rvert
 \le\norm{\rho_F-\sigma_F}_1
 \overset{\eqref{eq:finite-distance}}{\le}
 \sqrt{\exp(8\varepsilon_+(s)/b)-1}.
\]
The right-hand side depends only on \(s\), and
Lemma~\ref{lem:exhaustion} extends the bound to the full regional
algebras.

If \(\varepsilon_+(s)\le b/8\), then \(x=8\varepsilon_+(s)/b\)
lies in \([0,1]\), and the inequality \(e^x-1\le2x\) for \(0\le x\le1\)
gives
\[
 \sqrt{\exp(8\varepsilon_+(s)/b)-1}
 \le\sqrt{16\varepsilon_+(s)/b}
 \overset{\eqref{eq:envelopes}}{=}\frac{4\sqrt{A_+/b}}{s}.
\]
Hence every pair of self-adjoint contractions \(A\) and \(B\) of the
two regional algebras satisfies
\begin{equation}
 |\omega(AB)-\omega(A)\omega(B)|
 \le\frac{Q_0}{s},\qquad
 Q_0=4\sqrt{A_+/b}.
 \label{eq:mass-uniform-cluster}
\end{equation}

\subsection{The sufficient separation ratio}

The geometric thresholds can be chosen to satisfy
\begin{equation}
 \begin{gathered}
 S_{\mathrm{rev}}\ge
 \max\left\{
 S_{\mathrm{pol}},\,5,\,
 c_{\mathrm{pol}}\,2^{1/(1-\theta)}
 K_{\mathrm{pol}}^{\theta/(1-\theta)}
 \right\},\\
 S_{\mathrm{geo}}\ge\max\{5,S_b,S_{\mathrm{rev}}\}.
 \end{gathered}
 \label{eq:geometric-thresholds}
\end{equation}
The forward geometric condition \eqref{eq:forward-geometry}, through
\(S_{\mathrm{geo}}\ge S_b\), requires
\(S_{\mathrm{geo}}\ge2\mathfrak k_b=(4+80\sqrt3)\Lambda_b\), because
\(\mathfrak k_b=2\Lambda_b+8\sqrt3B_b\) and \(B_b=5\Lambda_b\). Since
\(\Lambda_b=2^{15}/{\sin(\pi b)}\ge2^{15}\) by \eqref{eq:Lambda}, the
bound \((4+80\sqrt3)\Lambda_b\) is at least
\((4+80\sqrt3)2^{15}>4.6\times10^6\) for every \(b\in[1/2,1)\). For the
present choice \(b=1-2^{-20}\), one has
\(\sin(\pi b)=\sin(\pi/2^{20})\le\pi/2^{20}\), which gives the numerical
lower requirement
\[
 S_{\mathrm{geo}}
 \ge(4+80\sqrt3)\Lambda_b
 =\frac{(4+80\sqrt3)2^{15}}{\sin(\pi/2^{20})}
 \ge\frac{(4+80\sqrt3)2^{35}}{\pi}
 >1.5\times10^{12}.
\]

For the orders \eqref{eq:orders}, the probability exponents
\eqref{eq:nus} become
\begin{equation}
 \nu_+=\frac{1+b}{4b^2},\qquad
 \nu_-=2^{18}-1,\qquad
 \eta=1-\frac1{4b^2}.
 \label{eq:completion-exponents}
\end{equation}
We denote the probability contribution to the separation ratio by
\begin{equation}
 C_{\mathrm{prob}}
 =Q_0\exp\!\left(
 \frac{2\log2/(ab)+\nu_++\nu_-}{\eta}\right).
 \label{eq:Cprob}
\end{equation}

Table~\ref{tab:constants} gives the dependency chain. Constants defined
by finite covers, cutoff norms, and approximation bounds are known
only to be finite.

{\small
\begin{longtable}{@{}p{0.17\textwidth}p{0.60\textwidth}p{0.17\textwidth}@{}}
\caption*{\normalsize\scshape\tablename~\thetable.\upshape\enspace
Constants entering the sufficient ratio \(C_*\).}
\label{tab:constants}\\
\toprule
Constant & Choice or definition & Source\\
\midrule
\endfirsthead
\toprule
Constant & Choice or definition & Source\\
\midrule
\endhead
\bottomrule
\endfoot
\(C_c,\lambda_c\)
 & \(C_c\) from the finite coarea and submean cover;
   \(\lambda_c=2048/{\sin(2\pi c)}\).
 & Theorem~\ref{thm:evaluation}\\
\(\Lambda_b,B_b\)
 & \(\Lambda_b=2^{15}/{\sin(\pi b)}\), \(B_b=5\Lambda_b\).
 & \eqref{eq:Lambda}\\
\(M_b\)
 & Weighted moment coefficient, depending only on \(b\).
 & Lemma~\ref{lem:moments}\\
\(\mathfrak d_b,d_{0,b},\mathfrak k_b\)
 & \(2B_b,\ 4\sqrt3B_b,\ 2\Lambda_b+8\sqrt3B_b\).
 & Section~\ref{sec:forward}\\
\(S_b,F_b\)
 & \(S_b\ge\max\{4,4d_{0,b},2\mathfrak k_b\}\);
   \(F_b\) from \eqref{eq:forward-refined}.
 & Theorem~\ref{thm:forward}\\
\(r_0,S_{\mathrm{pol}}\)
 & \(r_0=1/100\); \(S_{\mathrm{pol}}\) covers the fixed analytic
   neighbourhoods and polynomial error estimates.
 & Lemma~\ref{lem:polynomial}\\
\(C_1,C_2,c_{\mathrm{pol}},K_{\mathrm{pol}}\)
 & Constants of \eqref{eq:increment-unweighted} and
   \eqref{eq:increment-weighted}; \(c_{\mathrm{pol}}\) depends only on
   \(r_0\), and the others include the fixed cutoff norms, coefficient
   bound \(K_{\mathrm{mid}}\), rank bound, and \(M_{1/2}\).
 & Lemma~\ref{lem:polynomial}\\
\(\Lambda_0,\theta,a,\gamma\)
 & \(2^{15},\ 2^{-17},\ 2^{-18},\ 2(1-2^{-17})\).
 & \eqref{eq:reverse-exponents}\\
\(S_{\mathrm{rev}},\mathcal J_0\)
 & \(S_{\mathrm{rev}}\) as in \eqref{eq:geometric-thresholds};
   \(\mathcal J_0=52\cdot256\,C_1^{1-\theta}C_2^\theta\).
 & \eqref{eq:geometric-thresholds}, \eqref{eq:J0}\\
\(r_-,r_+,b\)
 & \(2^{18}/(2^{18}-1),\ 2^{20},\ 1-2^{-20}\).
 & \eqref{eq:orders}\\
\(c_a,c_b\)
 & \(2^{1-2a},\ (1-b)/b=1/(2^{20}-1)\).
 & \eqref{eq:scalar-reserve}\\
\(\eta,\nu_+,\nu_-\)
 & \(1-1/(4b^2),\ (1+b)/(4b^2),\ 2^{18}-1\).
 & \eqref{eq:completion-exponents}\\
\(A_+,A_-\)
 & \((r_+-1)F_b,\ 2^{2a-1}\mathcal J_0\).
 & \eqref{eq:Aconstants}\\
\(Q_0\)
 & \(4\sqrt{A_+/b}\).
 & \eqref{eq:mass-uniform-cluster}\\
\(S_{\mathrm{geo}}\)
 & Any constant at least \(\max\{5,S_b,S_{\mathrm{rev}}\}\).
 & \eqref{eq:geometric-thresholds}\\
\(C_{\mathrm{prob}}\)
 & \(Q_0\exp[(2\log2/(ab)+\nu_++\nu_-)/\eta]\).
 & \eqref{eq:Cprob}\\
\(C_*\)
 & Any finite number satisfying the four requirements below.
 & \eqref{eq:Cstar}\\
\end{longtable}
}

We choose a finite \(C_*\) satisfying
\begin{equation}
 C_*\ge
 \max\left\{
 S_{\mathrm{geo}},\
 \sqrt{8r_+A_+},\
 \left(\frac{4A_-}{a}\right)^{1/\gamma},\
 C_{\mathrm{prob}}
 \right\}.
 \label{eq:Cstar}
\end{equation}

\subsection{Proof of the main theorem and the fixed-regions corollary}
\label{subsec:main-proof}

\begin{mainrestatement}[Separation threshold]
There is a finite constant \(C_*>2\), independent of \(m\) and \(R\),
such that
\[
 \forall m,R>0,\qquad
 L\ge C_*R\quad\Longrightarrow\quad \beta_m(R,L)=1.
\]
\end{mainrestatement}

\begin{proof}[Proof of Theorem~\ref{thm:main}]
We first use unit-radius balls and arbitrary mass \(\mu>0\), and we fix
\(s\ge C_*\).

\smallskip
\emph{Step 1: the range of \(s\).} By the first entry of
\eqref{eq:Cstar} and by \eqref{eq:geometric-thresholds},
\[
 s\ge C_*\ge S_{\mathrm{geo}}\ge\max\{5,S_b,S_{\mathrm{rev}}\}.
\]
Hence \(s\ge S_b\) and \(s\ge S_{\mathrm{rev}}\), the ranges of
Theorems~\ref{thm:forward} and~\ref{thm:reverse}, so
\eqref{eq:finite-forward} and \eqref{eq:finite-reverse} hold for all
finite canonical source systems in the two balls. The same chain gives
\(C_*\ge5>2\).

\smallskip
\emph{Step 2: the envelopes.} The second and third entries of
\eqref{eq:Cstar} give \(s^2\ge8r_+A_+\) and \(s^\gamma\ge4A_-/a\). By
\eqref{eq:envelopes},
\[
 \varepsilon_+(s)=\frac{A_+}{s^2}\le\frac1{8r_+}\le\frac b8<\frac18,
 \qquad
 \varepsilon_-(s)=\frac{A_-}{s^\gamma}\le\frac a4,
\]
where \(1/r_+\le b\) because \(r_+b=r_+-1\ge1\). The bounds
\(\varepsilon_+(s)\le1/8\) and \(\varepsilon_-(s)\le a/4\) are the
conditions of \eqref{eq:Qforward} and \eqref{eq:Qreverse}, which
therefore hold, and with them
\eqref{eq:forward-events}--\eqref{eq:finite-distance}. The bound
\(\varepsilon_+(s)\le b/8\) is the condition of
\eqref{eq:mass-uniform-cluster}, which therefore holds for all
self-adjoint contractions \(A\in\A_\mu(D_1)\) and
\(B\in\A_\mu(D_1+(0,s\eone))\), the double-cone algebras of the two
balls.

\smallskip
\emph{Step 3: the event constants.} By \eqref{eq:full-K} and Step~2,
\[
 \log K_{r_+,+}\le8r_+\varepsilon_+(s)\le8r_+\cdot\frac1{8r_+}=1,
 \qquad
 \log K_{r_-,-}\le\frac{2r_-}{a}\varepsilon_-(s)
 \le\frac{2r_-}{a}\cdot\frac a4=\frac{r_-}2<1,
\]
the last inequality because \(r_-<2\). In particular, \(K_{r_+,+}\) and
\(K_{r_-,-}\) are finite.

\smallskip
\emph{Step 4: the cutoff \eqref{eq:probability-cutoff}.} By Step~2, the
correlation bound \eqref{eq:additive-cluster} holds with
\(q_{\mathrm{cl}}=Q_0/s\). The exponents \(\eta,\nu_+,\nu_-\) of
\eqref{eq:completion-exponents} are positive, and the last entry of
\eqref{eq:Cstar} is the number \(C_{\mathrm{prob}}\) of
\eqref{eq:Cprob}, so \(s\ge C_{\mathrm{prob}}\). Hence
\[
 \biggl(\frac{Q_0}{s}\biggr)^\eta
 K_{r_+,+}^{\nu_+}K_{r_-,-}^{\nu_-}
 \overset{\text{Step 3}}{\le}
 \biggl(\frac{Q_0}{s}\biggr)^\eta e^{\nu_++\nu_-}
 \overset{s\ge C_{\mathrm{prob}}}{\le}
 \biggl(\frac{Q_0}{C_{\mathrm{prob}}}\biggr)^\eta e^{\nu_++\nu_-}
 \overset{\eqref{eq:Cprob}}{=}
 2^{-2/(ab)},
\]
which is \eqref{eq:probability-cutoff} with \(q_{\mathrm{cl}}=Q_0/s\).

\smallskip
\emph{Step 5: the unit-radius conclusion.} By \eqref{eq:time-slice},
the algebras \(\A_\mu(D_1)\) and \(\A_\mu(D_1+(0,s\eone))\) are
generated by the Weyl operators of sources in the disjoint balls
\(B_1(0)\) and \(B_1(s\eone)\). These Weyl operators commute, because
the symplectic form vanishes between the two balls, so the two
algebras commute. The orders \eqref{eq:orders} satisfy \(r_+>2\) and
\(1<r_-<2\), and they give \(b=1-1/r_+\) and \(a=1-1/r_-\), as in
\eqref{eq:probability-exponents}, with \(\eta>0\). With Steps~3
and~4, every hypothesis of Theorem~\ref{thm:probability} holds, and
that theorem gives
\[
 \forall\mu>0,\qquad
 s\ge C_*\quad\Longrightarrow\quad\beta_\mu(1,s)=1.
\]

\smallskip
\emph{Step 6: scaling.} Let \(m,R>0\) and \(L\ge C_*R\). Then \(mR>0\)
and \(L/R\ge C_*\), so Lemma~\ref{lem:scaling} and Step~5 give
\[
 \beta_m(R,L)=\beta_{mR}(1,L/R)=1.
\]
This proves \eqref{eq:main}. Every entry of \eqref{eq:Cstar} is
independent of \(\mu\), so the finite constant \(C_*\) does not depend
on \(m\) or \(R\).
\end{proof}

\begin{proof}[Proof of Corollary~\ref{cor:fixed}]
Let \(m>0\), let \(O_1,O_2\) be bounded regions, and let \(e\) be the
spacelike vector in the given direction with \(e\cdot e=-1\), so that
the translates of \(O_2\) in that direction are \(O_2+Le\), \(L>0\).

\smallskip
\emph{Step 1: a spatial direction.} Since \(e\cdot e=-1\), a boost
followed by a rotation gives a proper orthochronous Lorentz
transformation \(\Lambda\) with \(\Lambda e=(0,\eone)\). By Poincar\'e
covariance of the free field, a unitary operator that fixes the vacuum
vector carries \(\A_m(O)\) onto \(\A_m(\Lambda O)\) for every region
\(O\), so, as in Step~4 of the proof of Lemma~\ref{lem:scaling},
\[
 \beta\bigl(\omega;\A_m(O_1),\A_m(O_2+Le)\bigr)
 =\beta\bigl(\omega;\A_m(\Lambda O_1),\A_m(\Lambda O_2+(0,L\eone))\bigr).
\]

\smallskip
\emph{Step 2: double cones centred at time zero.} The bounded set
\(\Lambda O_1\cup\Lambda O_2\) lies in a double cone \(D_R\) centred at
the origin, so that
\[
 \Lambda O_1\subset D_R,\qquad
 \Lambda O_2+(0,L\eone)\subset D_R+(0,L\eone).
\]
The radius \(R\) does not depend on \(L\), and the centres \(0\) and
\((0,L\eone)\) lie on the time-zero slice at distance \(L\). For
\(L>2R\), the bases \(B_R(0)\) and \(B_R(L\eone)\) are disjoint, and
the algebras \(\A_m(D_R)\) and \(\A_m(D_R+(0,L\eone))\) commute, as in
Step~5 of the proof of Theorem~\ref{thm:main}.

\smallskip
\emph{Step 3: isotony.} If one region is contained in another, its
local algebra is contained in that of the other.
The supremum \eqref{eq:beta} can only increase when both algebras are
enlarged, so the inclusions of Step~2 give
\[
 \beta\bigl(\omega;\A_m(\Lambda O_1),\A_m(\Lambda O_2+(0,L\eone))\bigr)
 \le\beta\bigl(\omega;\A_m(D_R),\A_m(D_R+(0,L\eone))\bigr)
 =\beta_m(R,L).
\]

\smallskip
\emph{Step 4: conclusion.} Let \(L\ge C_*R\), which exceeds \(2R\). By
Steps~1 to~3, one unitary carries \(\A_m(O_1)\) and \(\A_m(O_2+Le)\) into
the commuting algebras \(\A_m(D_R)\) and \(\A_m(D_R+(0,L\eone))\), so
\(\A_m(O_1)\) and \(\A_m(O_2+Le)\) commute, and their Bell correlation
is at most \(\beta_m(R,L)\), which equals \(1\) by
Theorem~\ref{thm:main}. The identities \(A_0=A_1=B_0=B_1=I\) give the
value \(1\) in \eqref{eq:beta}. Hence the Bell correlation of \(\A_m(O_1)\)
and \(\A_m(O_2+Le)\) equals \(1\) for all \(L\ge C_*R\).
\end{proof}

\section{Open problems}\label{sec:discussion}

\begin{enumerate}[label=\textup{(\arabic*)},leftmargin=*]
\item \textbf{Wedges at a positive gap and growing regions.}
For complementary wedges, the Summers--Werner results give maximal
Bell correlation \cite{SW87c,SW88}. Wedges separated by a positive gap
are unbounded localization regions, and their Bell correlation at large
gap is open. Determining this Bell correlation requires control of
source families extending arbitrarily far along the wedge. The Bell
correlation at large
separation is also open for bounded regions whose diameters grow in
proportion to their separation, at ratios that Theorem~\ref{thm:main}
does not cover.

\item \textbf{Other Bell inequalities, measurement settings, and local filtering.}
The probability criterion of Theorem~\ref{thm:probability} uses the
four correlators of CHSH, or equivalently two binary settings at each
site. For larger numbers of settings or outcomes, the problem is to
find analogous event criteria.
The relation between the eventual value \(\beta=1\) of the CHSH
correlation \eqref{eq:beta}, other Bell inequalities, and filtered output
states remains to be determined.

\item \textbf{Other dimensions and fields.}
The null-disc construction of Lemma~\ref{lem:null-discs} needs at least
three complex spatial coordinates, so one or two spatial dimensions
require a replacement for that construction. Every dimension other than three
requires the appropriate source-space and covariance estimates. Extensions to
other free fields and to interacting models pose corresponding
questions about local modular structure and covariance comparison.
The independence of \(C_*\) from the mass also motivates a continuity
analysis of regional Bell correlations at \(m=0\), including the
source-space identifications at that endpoint.

\item \textbf{Numerical enclosure of the sufficient ratio.}
The finite Lorentz cover determines \(C_c\) existentially; enclosing
the Jacobians and submean constants of the cover is the first
quantitative task.
These constants then enter \(M_b,F_b,\mathcal J_0\), and ultimately \(C_*\).
Even before those enclosures, the probability contribution
\eqref{eq:Cprob} obeys, along \(r_+=4/a\) (so that \(b=1-a/4\),
\(2\log2/(ab)=2\log2/a+O(1)\), and, by \eqref{eq:probability-exponents}
and \eqref{eq:nus}, \(\nu_+=(1+b)/(4b^2)=O(1)\), \(\nu_-=1/a-1\) and
\(\eta=1-1/(4b^2)=3/4+O(a)\)),
\begin{align*}
 \log_{10}(C_{\mathrm{prob}}/Q_0)
 &\overset{\eqref{eq:Cprob}}{=}
 \frac{2\log2/(ab)+\nu_++\nu_-}{\eta\log10}
 =\frac{(2\log2+1)/a+O(1)}{(3/4+O(a))\log10}\\
 &\overset{\phantom{\eqref{eq:Cprob}}}{=}
 \frac{4(2\log2+1)}{3\log10}\,\frac1a+O(1)
 =
 \frac{1.38180596\ldots}{a}+O(1).
\end{align*}
For the parameters in Table~\ref{tab:constants},
\[
 C_{\mathrm{prob}}
 =Q_0\cdot10^{362\,232.284166\ldots},
 \qquad C_*\ge C_{\mathrm{prob}}.
\]
A numerical enclosure of the full \(C_*\) is an open
problem, and so is the true threshold, the infimum of the \(L_0\) such
that \(\beta_m(R,L)=1\) for all \(L\ge L_0\).

\item \textbf{The regime \(\boldsymbol{L<C_*R}\).}
The dependence of \(\beta_m(R,L)\) on \(mR\) and \(L/R\), including
the region of strict violation and the approach to the classical
value, remains to be determined for \(2R<L<C_*R\).
Tangent-region maximality provides the short-distance reference
point \cite{SW88}. Results on selected smeared modes, such as the
two-mode separability scale in \cite{Zych10}, give mode-dependent
scales. Relating these scales to the
supremum over the full local algebras is a further quantitative
problem.
\end{enumerate}

\appendix
\section{Bargmann-kernel derivation of the Gaussian sandwich formulas}
\label{app:gaussian}

The formulas \eqref{eq:sandwich} follow from the Bargmann kernel
\cite{Bargmann61}; the resulting covariance identities agree with the
Gaussian product formalism by which Seshadreesan, Lami and Wilde
compute \(\Qtil_r(\rho\|\sigma)\) \cite[Theorem~22]{SLW18}. The sandwich
\(\Theta=\sigma^{-\tau/2}\rho\sigma^{-\tau/2}\) of \eqref{eq:sandwich},
which contains inverse powers of \(\sigma\), is defined as a positive
quadratic form, and its coherent-state diagonal supplies the trace-class
criterion.

For the Bargmann kernel and the matrices built from it, quadratures are ordered as
\((q_1,\ldots,q_n,\allowbreak p_1,\ldots,p_n)\). This ordering differs
from the ordering \((q_1,p_1,\ldots,q_n,p_n)\), in which the commutation
matrix is \(J_n\), by an orthogonal permutation matrix \(\Pi\). The change
of coordinates conjugates every covariance matrix by \(\Pi\) (Step~1 of
the proof of Theorem~\ref{thm:gaussian}), and conjugation by \(\Pi\)
commutes with sums, products, inverses and positive square roots and
preserves determinants; so \eqref{eq:Vtau-derived} and
\eqref{eq:trace-derived}, whose matrices are formed from covariance
matrices in this way, hold in either ordering.

The annihilation and number
operators are \(a_j=(q_j+ip_j)/\sqrt2\) and \(N_j=a_j^*a_j\). The Fock
vacuum \(|0\rangle\) is annihilated by every \(a_j\) and has covariance
\(I/2\), and \(|k\rangle\), \(k\in\N^n\), is the number basis,
\(N_j|k\rangle=k_j|k\rangle\).

For a centred Gaussian state
\(\rho\), define
\[
 \mathsf C_\rho=V_\rho+\tfrac12I,\qquad
 D_{\mathrm B}=\frac1{\sqrt2}
 \begin{pmatrix}I&I\\ iI&-iI\end{pmatrix},\qquad
 S_{\mathrm B}=\begin{pmatrix}0&I\\I&0\end{pmatrix}.
\]
For \(z\in\C^n\), the unnormalized exponential vectors and the
normalized coherent states are
\[
 |e(z)\rangle
 =\exp\!\left(\sum_jz_ja_j^*\right)|0\rangle
 =\sum_{k\in\N^n}\frac{z^k}{\sqrt{k!}}\,|k\rangle,\qquad
 |z\rangle=e^{-|z|^2/2}|e(z)\rangle,
\]
with \(z^k=\prod_jz_j^{k_j}\) and \(k!=\prod_jk_j!\). The kernel of
\(\rho\) between exponential vectors is
\begin{equation}
 \begin{split}
 \langle e(z)|\rho|e(w)\rangle
 &=
 \frac1{\sqrt{\det\mathsf C_\rho}}
 \exp\!\left[
 \frac12
 \begin{pmatrix}\bar z\\w\end{pmatrix}^{\!T}
 M_\rho
 \begin{pmatrix}\bar z\\w\end{pmatrix}
 \right],\\
 M_\rho&=S_{\mathrm B}
       -D_{\mathrm B}^T\mathsf C_\rho^{-1}D_{\mathrm B}.
 \end{split}
 \label{eq:Bargmann}
\end{equation}

\smallskip
\emph{Step 1: the kernel of a centred Gaussian state.}
To obtain this formula, the overlap with a normalized coherent state
is computed as the Fourier transform of the product of the Gaussian
characteristic function and the vacuum characteristic function. With
the characteristic function \(\chi_A(u)=\Tr(Ae^{iu^T(q,p)})\),
\(u\in\R^{2n}\), of a trace-class operator \(A\), where \((q,p)\) is the
vector of quadratures in this order, any two trace-class operators
satisfy the trace formula
\[
 \Tr(AB)=\frac1{(2\pi)^n}\int_{\R^{2n}}\chi_A(u)\,\chi_B(-u)\,d^{2n}u.
\]
The state \(\rho\) has \(\chi_\rho(u)=e^{-\frac12u^TV_\rho u}\), and the
coherent projection has
\(\chi_{|z\rangle\langle z|}(u)=e^{-\frac14u^Tu}e^{iu^Tr(z)}\), the
characteristic function of \(|0\rangle\) times a phase, where
\(r(z)=\sqrt2(\Rea z,\Ima z)\) is the vector of quadrature means of
\(|z\rangle\). With
\(e^{-\frac12u^TV_\rho u}e^{-\frac14u^Tu}=e^{-\frac12u^T\mathsf C_\rho u}\),
the trace formula and the Fourier Gaussian integral give
\begin{align*}
 \langle z|\rho|z\rangle
 &=\Tr\bigl(\rho\,|z\rangle\langle z|\bigr)
 =\frac1{(2\pi)^n}\int_{\R^{2n}}
   e^{-\frac12u^TV_\rho u}\,e^{-\frac14u^Tu}\,e^{-iu^Tr(z)}\,d^{2n}u\\
 &=\frac1{\sqrt{\det\mathsf C_\rho}}
 \exp[-\tfrac12 r(z)^T\mathsf C_\rho^{-1}r(z)].
\end{align*}

Since \(r(z)=D_{\mathrm B}(\bar z,z)\) and
\(D_{\mathrm B}^TD_{\mathrm B}=S_{\mathrm B}\), the exponent of
\eqref{eq:Bargmann} at \(w=z\) is
\[
 \frac12
 \begin{pmatrix}\bar z\\z\end{pmatrix}^{\!T}
 M_\rho
 \begin{pmatrix}\bar z\\z\end{pmatrix}
 =\frac12r(z)^Tr(z)-\frac12r(z)^T\mathsf C_\rho^{-1}r(z)
 =|z|^2-\frac12r(z)^T\mathsf C_\rho^{-1}r(z).
\]
Multiplication of the overlap by \(e^{|z|^2}\), which turns
\(\langle z|\rho|z\rangle\) into \(\langle e(z)|\rho|e(z)\rangle\),
therefore gives \eqref{eq:Bargmann} on the diagonal. Both sides of
\eqref{eq:Bargmann} are entire functions of \((\bar z,w)\in\C^{2n}\):
the left side is the series
\(\sum_{k,l}\bar z^kw^l\langle k|\rho|l\rangle/\sqrt{k!\,l!}\), which
converges everywhere because \(|\langle k|\rho|l\rangle|\le1\), and the
right side is the exponential of a quadratic polynomial.
Polarization of the entire sesquianalytic kernel gives
\eqref{eq:Bargmann} off the diagonal.

\smallskip
\emph{Step 2: the kernel and the coherent diagonal of the sandwich.}
In Williamson coordinates, the faithful reference state is
\[
 \sigma=Z_\sigma\prod_j\zeta_j^{N_j},\qquad
 \zeta_j=\frac{\upsilon_j-1/2}{\upsilon_j+1/2}\in(0,1),\qquad
 Z_\sigma=\prod_j(1-\zeta_j).
\]
For \(0<\tau<1\), let
\[
 \mathbb T_\sigma=\diag(\zeta_j^{-\tau/2})\oplus\diag(\zeta_j^{-\tau/2}).
\]
Since \(\sigma|k\rangle=Z_\sigma\zeta^k|k\rangle\), with
\(\zeta^k=\prod_j\zeta_j^{k_j}\) as for \(z^k\), the series of
\(|e(z)\rangle\) gives
\(\sigma^{-\tau/2}|e(z)\rangle=Z_\sigma^{-\tau/2}|e(z')\rangle\) with
\(z'_j=\zeta_j^{-\tau/2}z_j\); in particular, \(|e(z)\rangle\) lies in
the domain of \(\sigma^{-\tau/2}\). With \(w'_j=\zeta_j^{-\tau/2}w_j\)
likewise, \((\bar z',w')=\mathbb T_\sigma(\bar z,w)\), and the quadratic
form \(\Theta=\sigma^{-\tau/2}\rho\sigma^{-\tau/2}\) has the kernel
\[
 \langle e(z)|\Theta|e(w)\rangle
 =\langle\sigma^{-\tau/2}e(z)|\rho|\sigma^{-\tau/2}e(w)\rangle
 =Z_\sigma^{-\tau}\langle e(z')|\rho|e(w')\rangle,
\]
which is \eqref{eq:Bargmann} with \(z,w\) replaced by \(z',w'\) and its
prefactor multiplied by \(Z_\sigma^{-\tau}\).

On the diagonal \(w=z\), the commutation
\(D_{\mathrm B}\mathbb T_\sigma=\mathbb T_\sigma D_{\mathrm B}\) gives
\(D_{\mathrm B}(\bar z',z')=\mathbb T_\sigma r(z)\), while
\((\bar z',z')^TS_{\mathrm B}(\bar z',z')=2\sum_j\zeta_j^{-\tau}|z_j|^2
=r(z)^T\mathbb T_\sigma^2r(z)\) and \(2|z|^2=r(z)^Tr(z)\). The coherent
diagonal of \(\Theta\) is therefore
\begin{align*}
 \langle z|\Theta|z\rangle
 &\overset{\phantom{\eqref{eq:Bargmann}}}{=}e^{-|z|^2}\langle e(z)|\Theta|e(z)\rangle\\
 &\overset{\eqref{eq:Bargmann}}{=}
  \frac{Z_\sigma^{-\tau}}{\sqrt{\det\mathsf C_\rho}}
  \exp\Bigl[-\tfrac12r(z)^Tr(z)+\tfrac12r(z)^T\mathbb T_\sigma^2r(z)
  -\tfrac12r(z)^T\mathbb T_\sigma\mathsf C_\rho^{-1}\mathbb T_\sigma r(z)\Bigr]\\
 &\overset{\phantom{\eqref{eq:Bargmann}}}{=}\frac{Z_\sigma^{-\tau}}{\sqrt{\det\mathsf C_\rho}}
  \exp\bigl[-\tfrac12r(z)^T\mathsf C_\tau^{-1}r(z)\bigr],
\end{align*}
where, in real quadrature coordinates, \(\mathsf C_\tau^{-1}\) denotes
the real symmetric matrix
\begin{equation}
 \mathsf C_\tau^{-1}
 =I-\mathbb T_\sigma^2
  +\mathbb T_\sigma\mathsf C_\rho^{-1}\mathbb T_\sigma.
 \label{eq:Ctau-inverse}
\end{equation}

\smallskip
\emph{Step 3: positivity of \(\mathsf C_\tau^{-1}\).} The matrix
\(\mathbb T_\sigma\) is diagonal with entries \(\zeta_j^{-\tau/2}>1\),
so \(I-\mathbb T_\sigma^{-2}\) is diagonal with entries
\(1-\zeta_j^\tau\in(0,1)\), and \eqref{eq:Ctau-inverse} reads
\[
 \mathsf C_\tau^{-1}
 =\mathbb T_\sigma\bigl(\mathsf C_\rho^{-1}-(I-\mathbb T_\sigma^{-2})\bigr)
  \mathbb T_\sigma .
\]
Congruence by \(\mathbb T_\sigma^{-1}\), which gives
\(\mathsf C_\tau^{-1}>0\) exactly when
\(\mathsf C_\rho^{-1}>I-\mathbb T_\sigma^{-2}\), and inversion of
positive matrices, which reverses their order, show that
\begin{equation}
 \mathsf C_\tau^{-1}>0
 \quad\Longleftrightarrow\quad
 \mathsf C_\rho<(I-\mathbb T_\sigma^{-2})^{-1}
 \quad\Longleftrightarrow\quad
 V_\rho<\mathsf T_\tau,\qquad
 \mathsf T_\tau=V_\sigma^{[\tau]},
 \label{eq:gaussian-domain}
\end{equation}
where
\((I-\mathbb T_\sigma^{-2})^{-1}=\mathsf T_\tau+I/2\). This identity
holds because
\(\sigma^\tau/{\Tr\sigma^\tau}=\prod_j(1-\zeta_j^\tau)\zeta_j^{\tau N_j}\)
is \(\sigma\) with every \(\zeta_j\) replaced by \(\zeta_j^\tau\), while
the definition of \(\zeta_j\) is equivalent to
\(\upsilon_j+\tfrac12=(1-\zeta_j)^{-1}\). Hence
\(\mathsf T_\tau=V_\sigma^{[\tau]}\) is diagonal in Williamson coordinates,
with entries \(\upsilon_j^{[\tau]}\) satisfying
\(\upsilon_j^{[\tau]}+\tfrac12=(1-\zeta_j^\tau)^{-1}\), the diagonal
entries of \((I-\mathbb T_\sigma^{-2})^{-1}\).

\smallskip
\emph{Step 4: the trace and the operator \(\Theta\).} On finite
linear combinations \(\psi\) of number states,
\(\langle\psi|\Theta|\psi\rangle
=\langle\sigma^{-\tau/2}\psi|\rho|\sigma^{-\tau/2}\psi\rangle\ge0\), so
the sandwich defines a positive matrix form, with entries
\[
 \langle k|\Theta|l\rangle
 =Z_\sigma^{-\tau}(\zeta^k\zeta^l)^{-\tau/2}\langle k|\rho|l\rangle,
 \qquad k,l\in\N^n .
\]
By the series of \(|e(z)\rangle\), the coherent diagonal of \(\Theta\)
is the absolutely convergent series
\(\langle z|\Theta|z\rangle
=e^{-|z|^2}\sum_{k,l}\bar z^kz^l\langle k|\Theta|l\rangle/\sqrt{k!\,l!}\).
Angular integration over the phases of \(z_1,\ldots,z_n\) removes the
terms with \(k\ne l\), and
\(\pi^{-n}\int_{\C^n}e^{-|z|^2}|z^k|^2\,d^{2n}z=k!\), where
\(d^{2n}z=\prod_jd\Rea z_j\,d\Ima z_j\). Monotone convergence of the
nonnegative diagonal coefficients then identifies the integral with the
sum of the number-state diagonal entries:
\[
 \sum_{k\in\N^n}\langle k|\Theta|k\rangle
 =\frac1{\pi^n}\int_{\C^n}\langle z|\Theta|z\rangle\,d^{2n}z
 =\frac{Z_\sigma^{-\tau}}{(2\pi)^n\sqrt{\det\mathsf C_\rho}}
  \int_{\R^{2n}}e^{-\frac12r^T\mathsf C_\tau^{-1}r}\,d^{2n}r ,
\]
the last step by Step~2 and \(d^{2n}z=2^{-n}d^{2n}r\) for \(r=r(z)\).

Under \eqref{eq:gaussian-domain}, the matrix \eqref{eq:Ctau-inverse} is
positive definite, with inverse \(\mathsf C_\tau\), and the last
integral equals \((2\pi)^n(\det\mathsf C_\tau)^{1/2}\), so the diagonal
sum is finite. A positive matrix form with finite diagonal sum extends
uniquely to a positive trace-class operator. Hence \(\Theta\) is a
positive trace-class operator whose trace is the diagonal sum, and
\begin{equation}
 \Tr\Theta
 =Z_\sigma^{-\tau}
  \sqrt{\frac{\det\mathsf C_\tau}{\det\mathsf C_\rho}}
 <\infty.
 \label{eq:traceTheta}
\end{equation}

Since \(\sigma^{\tau/2}|k\rangle=Z_\sigma^{\tau/2}(\zeta^k)^{\tau/2}|k\rangle\),
the entries of \(\Theta\) give
\[
 \langle k|\sigma^{\tau/2}\Theta\sigma^{\tau/2}|l\rangle
 =Z_\sigma^{\tau}(\zeta^k\zeta^l)^{\tau/2}\langle k|\Theta|l\rangle
 =\langle k|\rho|l\rangle ,
\]
so \(\rho=\sigma^{\tau/2}\Theta\sigma^{\tau/2}\), both sides being
bounded operators with the same matrix in the number basis.

\smallskip
\emph{Step 5: the covariance of the normalized sandwich.} The
matrices \(\mathbb T_\sigma\), \(\mathsf T_\tau\) and \(\mathsf L\) are
diagonal, and \(\mathbb T_\sigma^2>I\). By Step~3,
\(\mathsf T_\tau+\tfrac12I=(I-\mathbb T_\sigma^{-2})^{-1}
=\mathbb T_\sigma^2(\mathbb T_\sigma^2-I)^{-1}
=I+(\mathbb T_\sigma^2-I)^{-1}\), and
\(\mathsf L^2=(\mathsf T_\tau-\tfrac12I)(\mathsf T_\tau+\tfrac12I)\);
hence
\[
 \mathsf T_\tau-\tfrac12I=(\mathbb T_\sigma^2-I)^{-1},\qquad
 \mathsf L=(\mathsf T_\tau^2-I/4)^{1/2}
 =\mathbb T_\sigma(\mathbb T_\sigma^2-I)^{-1}.
\]

These identities and the Woodbury matrix identity, in the form of
\cite[Lemma~5]{SLW18}, applied to \eqref{eq:Ctau-inverse} with the
invertible summands \(I-\mathbb T_\sigma^2\) and
\(\mathbb T_\sigma\mathsf C_\rho^{-1}\mathbb T_\sigma\), give
\begin{align*}
 \mathsf C_\tau
 &=(I-\mathbb T_\sigma^2)^{-1}
  -(I-\mathbb T_\sigma^2)^{-1}
   \bigl((I-\mathbb T_\sigma^2)^{-1}
   +\mathbb T_\sigma^{-1}\mathsf C_\rho\mathbb T_\sigma^{-1}\bigr)^{-1}
   (I-\mathbb T_\sigma^2)^{-1}\\
 &=-(\mathsf T_\tau-\tfrac12I)
  -\mathsf L(V_\rho-\mathsf T_\tau)^{-1}\mathsf L
 =-(\mathsf T_\tau-\tfrac12I)
  +\mathsf L(\mathsf T_\tau-V_\rho)^{-1}\mathsf L .
\end{align*}
The second line uses
\((I-\mathbb T_\sigma^2)^{-1}=-(\mathsf T_\tau-\tfrac12I)\),
\((I-\mathbb T_\sigma^2)^{-1}\mathbb T_\sigma=-\mathsf L\) and
\[
 (I-\mathbb T_\sigma^2)^{-1}
 +\mathbb T_\sigma^{-1}\mathsf C_\rho\mathbb T_\sigma^{-1}
 =\mathbb T_\sigma^{-1}
  \bigl(\mathsf C_\rho-\mathsf T_\tau-\tfrac12I\bigr)\mathbb T_\sigma^{-1}
 =\mathbb T_\sigma^{-1}(V_\rho-\mathsf T_\tau)\mathbb T_\sigma^{-1},
\]
which is invertible under \eqref{eq:gaussian-domain}.

By Step~2 and \eqref{eq:traceTheta}, the density operator
\(\Theta/{\Tr\Theta}\) has the coherent diagonal
\[
 \frac{\langle z|\Theta|z\rangle}{\Tr\Theta}
 =\frac1{\sqrt{\det\mathsf C_\tau}}
  \exp\bigl[-\tfrac12r(z)^T\mathsf C_\tau^{-1}r(z)\bigr]
 =\frac1{(2\pi)^n}\int_{\R^{2n}}
  e^{-\frac12u^T\mathsf C_\tau u}\,e^{-iu^Tr(z)}\,d^{2n}u .
\]
For every density operator, as in Step~1, the coherent diagonal is
\((2\pi)^{-n}\) times the Fourier transform at \(r(z)\) of the product
of its characteristic function and the vacuum characteristic function
\(e^{-\frac14u^Tu}\), which has no zeros. Uniqueness of the Fourier
transform therefore gives \(\Theta/{\Tr\Theta}\) the characteristic
function \(e^{-\frac12u^T(\mathsf C_\tau-\frac12I)u}\), so
\(\Theta/{\Tr\Theta}\) is the centred Gaussian state with covariance
\(V_\tau=\mathsf C_\tau-\tfrac12I\). With the formula for
\(\mathsf C_\tau\) above, the covariance of the normalized sandwich is
\begin{equation}
 V_\tau
 =-\mathsf T_\tau+\mathsf L(\mathsf T_\tau-V_\rho)^{-1}\mathsf L.
 \label{eq:Vtau-derived}
\end{equation}

\smallskip
\emph{Step 6: the trace formula.} The form of
\eqref{eq:Ctau-inverse} in Step~3 and
\(I-\mathbb T_\sigma^{-2}=(\mathsf T_\tau+\tfrac12I)^{-1}\) give
\[
 \mathsf C_\tau^{-1}
 =\mathbb T_\sigma\bigl(\mathsf C_\rho^{-1}
  -(\mathsf T_\tau+\tfrac12I)^{-1}\bigr)\mathbb T_\sigma
 =\mathbb T_\sigma\mathsf C_\rho^{-1}(\mathsf T_\tau-V_\rho)
  (\mathsf T_\tau+\tfrac12I)^{-1}\mathbb T_\sigma ,
\]
since \(\mathsf T_\tau+\tfrac12I-\mathsf C_\rho=\mathsf T_\tau-V_\rho\).
Taking determinants in the product form of \eqref{eq:Ctau-inverse}, and using
\(\mathsf T_\tau+\tfrac12I=\mathbb T_\sigma^2(\mathbb T_\sigma^2-I)^{-1}\)
from Step~5, gives
\[
 \frac{\det\mathsf C_\tau}{\det\mathsf C_\rho}
 =\frac{\det\bigl((\mathsf T_\tau+\tfrac12I)\mathbb T_\sigma^{-2}\bigr)}
  {\det(\mathsf T_\tau-V_\rho)}
 =
 \frac1{
 \det(\mathbb T_\sigma^2-I)\det(\mathsf T_\tau-V_\rho)}.
\]
All factors of \(\Tr\Theta\) except \(\det(\mathsf T_\tau-V_\rho)\) depend only on
\(\sigma,\tau\):
\[
 \Tr\Theta
 \overset{\eqref{eq:traceTheta}}{=}
 \frac{Z_\sigma^{-\tau}}{\sqrt{\det(\mathbb T_\sigma^2-I)}}\,
 \frac1{\sqrt{\det(\mathsf T_\tau-V_\rho)}} .
\]

At \(\rho=\sigma\), \(\Theta=\sigma^{1-\tau}\), and
\eqref{eq:gaussian-domain} holds because
\(\mathsf T_\tau-V_\sigma=\Therm_\tau(V_\sigma)>0\). Evaluation of the
formula for \(\Tr\Theta\) at \(\rho=\sigma\) therefore gives
\(Z_\sigma^{-\tau}\det(\mathbb T_\sigma^2-I)^{-1/2}
=\Tr\sigma^{1-\tau}(\det\mathsf H)^{1/2}\), with
\[
 \mathsf H=\mathsf T_\tau-V_\sigma,\qquad
 \mathsf Z=
 \mathsf H^{-1/2}(V_\rho-V_\sigma)\mathsf H^{-1/2}.
\]
The factorization
\(\mathsf T_\tau-V_\rho=\mathsf H-(V_\rho-V_\sigma)
=\mathsf H^{1/2}(I-\mathsf Z)\mathsf H^{1/2}\) gives
\(\det(\mathsf T_\tau-V_\rho)=\det\mathsf H\,\det(I-\mathsf Z)\). With
this value of \(Z_\sigma^{-\tau}\det(\mathbb T_\sigma^2-I)^{-1/2}\), the
formula for \(\Tr\Theta\) becomes the trace formula
\begin{equation}
 \Tr\Theta
 =\Tr\sigma^{1-\tau}
  \sqrt{\frac{\det\mathsf H}{\det(\mathsf T_\tau-V_\rho)}}
 =\frac{\Tr\sigma^{1-\tau}}{\sqrt{\det(I-\mathsf Z)}}.
 \label{eq:trace-derived}
\end{equation}
Equations \eqref{eq:Vtau-derived} and \eqref{eq:trace-derived} are
the two identities in \eqref{eq:sandwich}.

\section*{Data availability}

This article has no associated data.

\bibliographystyle{amsplain}
\bibliography{vacuum-bell}

\end{document}